\documentclass{amsart}

\usepackage{lineno}
%\linenumbers
\usepackage{BOONDOX-cal}
\usepackage{fontawesome}
\usepackage{wrapfig}
%\usepackage{boondox-cal} % \mathcal \mathbcal
%\usepackage{unilogo}
%\pagestyle{empty}
\usepackage{epigraph}
\usepackage{amssymb}
%\usepackage{mathbbol}
\usepackage{environ}
\usepackage{amsmath}
\usepackage[mathscr]{eucal}
\usepackage{mathrsfs}
\usepackage{xspace}
\usepackage{adjustbox}
\usepackage{tabu}
%\usepackage{thmtools}
\usepackage[utf8]{inputenc}
\usepackage[english]{babel}
\usepackage[all,2cell]{xy}
\usepackage{xcolor}
\usepackage{graphicx}
\usepackage{wasysym} % music
\usepackage{hyperref}
%%%%%%% ETIQUETES
\usepackage[notcite, notref]{showkeys}
\usepackage{pdflscape}
\usepackage{caption}
\usepackage{multicol}
\setlength{\multicolsep}{6pt plus 2pt minus 1.5pt}
\usepackage{longtable}

\newcommand\hmmax{0}  % default 3, for "heavy" symbols
\newcommand\bmmax{1}  % default 4, for bold symbols

\usepackage{bm}

\makeatletter
\newcommand{\leqnomode}{\tagsleft@true}
\newcommand{\reqnomode}{\tagsleft@false}

\newsavebox{\measure@tikzpicture}
\NewEnviron{scaletikzpicturetowidth}[1]{%
	\def\tikz@width{#1}%
	\def\tikzscale{1}\begin{lrbox}{\measure@tikzpicture}%
		\BODY
	\end{lrbox}%
	\pgfmathparse{#1/\wd\measure@tikzpicture}%
	\edef\tikzscale{\pgfmathresult}%
	\BODY
}

\makeatother

\usepackage{array}
\newcolumntype{L}[1]{>{\raggedright\let\newline\\\arraybackslash\hspace{0pt}}m{#1}}
\newcolumntype{C}[1]{>{\centering\let\newline\\\arraybackslash\hspace{0pt}}m{#1}}
\newcolumntype{R}[1]{>{\raggedleft\let\newline\\\arraybackslash\hspace{0pt}}m{#1}}


\definecolor{sg}{HTML}{df78ef}
\definecolor{sg1}{HTML}{ab47bc}
\definecolor{sg2}{HTML}{790e8b}


\usepackage{tikz}
\usetikzlibrary{mindmap,backgrounds, calc}
\usetikzlibrary{matrix,chains,scopes,positioning,arrows,fit}
\usetikzlibrary{positioning, shapes, patterns}
\usetikzlibrary{automata}
\usetikzlibrary{shapes.geometric, arrows,arrows.meta}
\usetikzlibrary{positioning,decorations.markings}
\usepackage{hyperref}
\usepackage[acronym]{glossaries}

\usepackage{imakeidx}
\makeindex


\usepackage{stmaryrd} % symbols computer sciences
\UseAllTwocells

%%%%%%%%%%%%%%%%%%%%%%%%%%%%%
\makeatletter
\newcommand*\bigcdot{\mathpalette\bigcdot@{.5}}
\newcommand*\bigcdot@[2]{\mathbin{\vcenter{\hbox{\scalebox{#2}{$\m@th#1\bullet$}}}}}
\makeatother

% Macros
\newcommand{\sg}[1]{\ensuremath{\protect\underline{#1}}}   % signatura (S,Sigma)
\newcommand{\fm}[1]{\ensuremath{\mathcal{#1}}}             % familia A
\newcommand{\al}[1]{\ensuremath{\protect\underline{#1}}}   % algebra (A,F)
\newcommand{\eq}[1]{\ensuremath{\protect\underline{#1}}}   % ecuacion (S,Sigma,E)

\newcommand{\cs}[1]{\ensuremath{\mathcal{#1}}}             % sistema de clausura
\newcommand{\co}[1]{\ensuremath{\mathrm{#1}}}              % operador clausura
\newcommand{\ct}[1]{\ensuremath{\mathbf{#1}}}              % categoria
\newcommand{\cd}[1]{\ensuremath{\mathfrak{#1}}}            % cardinal
\newcommand{\at}[1]%
{\ensuremath{\protect\underline{\mathbf{#1}}}} % algebraic theory
\newcommand{\mn}[1]{\ensuremath{\mathbf{#1}}}              % monada
\newcommand{\ad}[1]{\ensuremath{\mathbf{#1}}}              % adjuncion
\newcommand{\ec}[1]{\ensuremath{\mathcal{#1}}}             % familia de ecuaciones

\newcommand{\cl}[1]{\ensuremath{\overline{#1}}}            % Clausura
\newcommand{\op}[1]{\ensuremath{\operatorname{#1}}}        % Operador
\newcommand{\lop}[1]{\ensuremath{\operatorname*{#1}}}      % Operador grande
\newcommand{\tp}[1]{\langle#1\rangle}                      % tupla <a,b,c> <a|b>
\newcommand{\gu}[1]{\ensuremath{\boldsymbol{\mathcal{#1}}}}% Grothendieck Universe
\newcommand{\farg}{\cdot}                                  % farg. el argumento de una funcion. alt:{\_}

\newcommand{\mawh}[1]{\widehat{#1}}                        % \section{Math accents}
\newcommand{\mawt}[1]{\widetilde{#1}}

\newcommand{\mab}[1]{\Bar{#1}}
\newcommand{\mabb}[1]{\Bar{\Bar{#1}}}
\newcommand{\mah}[1]{\Hat{#1}}
\newcommand{\mahh}[1]{\Hat{\Hat{#1}}}
\newcommand{\mat}[1]{\Tilde{#1}}
\newcommand{\matt}[1]{\Tilde{\Tilde{#1}}}
\newcommand{\maol}[1]{\overline{#1}}

\newcommand{\dmd}{\diamond}                                % diamante
\newcommand{\trl}{\triangleleft}                           % triangleleft
\newcommand{\trr}{\triangleright}                          % triangleright
\newcommand{\prt}{\partial}                                % partial d

% $S$-palabra $S$\nbd conjunto  era:\newcommand{\s}[1]{$S$\nbd}
\newcommand{\h}[1][]                                       % \h-conjunto $\Omega^{\h}$
{\ifthenelse{\boolean{mmode}}%
	{$\mathrm{h}$}%
	{h\nobreakdash#1\hspace{0pt}}}

%\section{Arrows}
\newcommand{\rar}{\rightarrow}
\newcommand{\Rar}{\Rightarrow}

\newcommand{\true}{\top}                 % true
\newcommand{\false}{\bot}                % false
% neg negacion
\newcommand{\con}{\wedge}                % con (and definido)
\newcommand{\dis}{\vee}                  % dis (or definido)
\newcommand{\lcon}{\bigwedge}            % con grande
\newcommand{\ldis}{\bigvee}              % dis grande
\newcommand{\imp}{\rightarrow}           % imp -->
\newcommand{\coimp}{\leftrightarrow}     % coimp <->
\newcommand{\nec}{\mathbin{\square}}               % necesidad
\newcommand{\pos}{\lozenge}              % posibilidad

\newcommand{\ded}{\vdash}                % \ded A |- a
\newcommand{\nded}{\nvdash}              % \ded A |/- a
\newcommand{\rded}{\dashv}               % \rded a -| A

\newcommand{\nmodels}{\nvDash}           % \nmodels |/=
\newcommand{\cons}{\Vdash}               % \cons consecuencia semantica

\newcommand{\mrule}[2]{
	\begin{array}{@{}c@{}}
		#1 \\
		\hline
		#2
\end{array} }
% \brule para 2
\newcommand{\brule}[2]{
	\begin{array}{@{}c@{}}
		#1 \\
		\hline \hline
		#2
\end{array} }

\newcommand{\serial}{\scriptstyle{+}} % diagramas serialmente conmutativos

\DeclareMathOperator{\N}{\op{N}}   % Numeros naturales
\DeclareMathOperator{\U}{\boldsymbol{\mathcal{U}}}   % Universo de Grothendieck
\DeclareMathOperator{\Ob}{Ob}      % Objetos de una categoria Ob(C)
\DeclareMathOperator{\Mor}{Mor}    % Morfismos de una categoria Mor(C)
\DeclareMathOperator{\dom}{dom}    % funcion dominio
\DeclareMathOperator{\cod}{cod}    % funcion codominio
\newcommand{\id}{\mathrm{id}}        % id(entidad)
\DeclareMathOperator{\Hom}{Hom}    % Homomorfismos Hom_{C}(A,B)
\DeclareMathOperator{\End}{End}    % Endomorfismos End_{C}(A)
\DeclareMathOperator{\Aut}{Aut}    % Automorfismos Aut_{C}(A)
\DeclareMathOperator{\Exp}{Exp}    % Exponenciaci\'{o}n
% \subsection{Composicion}
\newcommand{\comp}{\circ}          % composicion
\newcommand{\dcomp}{\diamond}      % diamond composition
\newcommand{\hcomp}{\ast}          % composici\'{o}n horizontal de trans. nat.
\newcommand{\vcomp}{\circ}         % composici\'{o}n vertical de trans. nat.
\newcommand{\adcomp}%
{\overset{\operatorname{ad}}{\comp}} % Ad-composicion cuadrados adjuntos
\newcommand{\funcomp}%
{\overset{\operatorname{fn}}{\comp}}% Fun-composici\'{o}n cuadrados adjuntos
\newcommand{\tcomp}{\mathbin{\Tilde{\circ}}}  % Tilde composition
\newcommand{\thcomp}{\mathbin{\Tilde{\ast}}} % Tilde horizontal composition

\DeclareMathOperator{\Id}{Id}      % Id Functor identidad
\DeclareMathOperator{\arityr}{Arr}    % Arr categoria de flechas
\DeclareMathOperator{\Fun}{Fun}    % Fun(A,B) categoria functores y t.n. de A en B
\DeclareMathOperator{\Nat}{Nat}    % Nat(F,G) t.n. de F en G
\DeclareMathOperator{\In}{In}      % In Functor inclusion
% Pr Functor proyeccion
\DeclareMathOperator{\opp}{op}     % C^{op} categoria dual
\newcommand{\dl}[1]{#1^{\partial}} % si f:A->B en C, \dl{f}:B->A en C^{op}

\newcommand{\ccat}{\mathbin{\downarrow}}

\newcommand{\sccat}
{\mathbin{\kern-1pt\raisebox{6pt}{.}\kern-5pt
		\downarrow\kern-5pt\raisebox{6pt}{.}\kern-1pt}}

\newcommand{\parrow}[1]%           % Projective arrow
{\underset{{\displaystyle \raisebox{5pt}%
			{$\longleftarrow$}}}{\op{#1}}{\,}}
\newcommand{\iarrow}[1]%           % Inductive arrow
{\underset{{\displaystyle \raisebox{5pt}%
			{$\longrightarrow$}}}{\op{#1}}{\,}}

\newcommand{\Cns}{\parrow{Cns}}    % Conos (con flecha)
\newcommand{\CoCns}{\iarrow{Cns}}  % Coconos (con flecha)

\newcommand{\Lim}{\varprojlim}     % Limites con flecha
\newcommand{\CoLim}{\varinjlim}    % Limites con flecha

\DeclareMathOperator{\Equ}{Eq}     % Igualador
\DeclareMathOperator{\equ}{eq}     % igualador flecha
\DeclareMathOperator{\Coeq}{Coeq}  % Coigualador
\DeclareMathOperator{\coeq}{coeq}  % coigualador flecha
\DeclareMathOperator{\Pb}{Pb}      % pullbach producto fibrado (fibered product)
\DeclareMathOperator{\Po}{Po}      % pushout suma amalgamada (amalgamated sum)

\newcommand{\ladj}{\!\dashv\!}         % ladj A -I B
\newcommand{\radj}{\!\vdash\!}         % radj A l- B
\newcommand{\uadj}{\top}           % uadj T
\newcommand{\dadj}{\bot}           % dadj _I_
\DeclareMathOperator{\lad}{lad}    % f^{\lad}
\DeclareMathOperator{\rad}{rad}    % f^{\rad}
\DeclareMathOperator{\adj}{adj}    % adj

\newcommand{\Grtctv}[1]{\textstyle{\int^{\scriptscriptstyle#1}}}
\newcommand{\Grtcv}[1]{\textstyle{\int_{\scriptscriptstyle#1}}}
\newcommand{\Grtc}{\textstyle{\int^{\ast}\!}}
\newcommand{\Grt} {\textstyle{\int_{\ast}  }}
\DeclareMathOperator{\Shp}{Shp}    % Functor shape
\DeclareMathOperator{\wreath}{wr}  % wreath product (producto orlado)
\DeclareMathOperator{\St}{St}      % funcion St(ructuras) de los constructos
\DeclareMathOperator{\Ad}{Ad}      % funcion Ad(misibles) de los constructos
\DeclareMathOperator{\Con}{Con}    % Con(K) Categor\'{i}as concretas sobre K
\DeclareMathOperator{\Struct}{Struct}% Struct(K) Categor\'{i}as de K-objetos con estructura

\DeclareMathOperator{\Pow}{Pow}    %

\DeclareMathOperator{\Set}{Set}    % Functor Set
\DeclareMathOperator{\Diag}{Diag}  % Functor Diag
\DeclareMathOperator{\Kl}{Kl}      % Kleisli
\DeclareMathOperator{\EM}{EM}      % Eilenberg-Moore

\DeclareMathOperator{\Sig}{Sig}    % Functor Sig
\DeclareMathOperator{\BTh}{BTh}    % Funtor Thr teorias algebraicas de Benabou
\DeclareMathOperator{\Thr}{Thr}    % Funtor Thr teorias algebraicas
\DeclareMathOperator{\thp}{th}     % Presentaciones de teorias

\DeclareMathOperator{\cng}{cn}     % Categor\'{i}a conjugada
\DeclareMathOperator{\trn}{tr}     % Categor\'{i}a transpuesta
\DeclareMathOperator{\sym}{sm}     % Categor\'{i}a simetrica
\DeclareMathOperator{\tns}{tn}     % transformaciones
\DeclareMathOperator{\alg}{alg}    % algebraico
\DeclareMathOperator{\LSq}{LSq}    % Lax squares. lax cuadrados
\DeclareMathOperator{\PSq}{PSq}    % pseudo squares. pseuado cuadrados.
\newcommand{\inv}{^{\text{-}1}}    % inverso
\DeclareMathOperator{\Sub}{Sub}    % Subobjeto (algebra, conjunto) Sub (A)
\DeclareMathOperator{\ch}{ch}      % caracter de una funcion
\DeclareMathOperator{\Rel}{Rel}    % Relaciones Rel(A,B)
\DeclareMathOperator{\ND}{ND}      % Hom_{\ND} aplicaciones no deterministas
\DeclareMathOperator{\Fnc}{Fnc}    % Funciones Fnc(A,B)
\DeclareMathOperator{\Eqv}{Eqv}    % Eqv(ivalencias)  Equivalencia
\DeclareMathOperator{\Eg}{Eg}      % Eqv(ivalencias)  Equivalencia generada
\newcommand{\grph}[1]{\Gamma_{#1}} % Grafo de una aplicacion
\DeclareMathOperator{\Img}{Im}     % Imagen de una aplicaci\'{o}n (Im definido)
\DeclareMathOperator{\inc}{in}     % inclusion (in definido)
\DeclareMathOperator{\supp}{supp}  % soporte
\newcommand{\eset}{\emptyset}      % emptyset conjunto vac\'{i}o
\newcommand{\lmapsto}{\longmapsto} % lmapsto
\newcommand{\rest}%
{\mathnormal{\restriction}}        % restriction

\newcommand{\nin}{\not\in}         % no pertenece
\newcommand{\incl}{\subseteq}      % inclusion A <= B
\newcommand{\nincl}{\nsubseteq}    % no inclusion A </= B
\newcommand{\sincl}{\subset}       % inclusion estricta A < B
\newcommand{\rincl}{\supseteq}     % reverse inclusion B > A
\newcommand{\rsincl}{\supset}      % reverse strict inclusion B > A
\newcommand{\tprod}{\textstyle{\prod}}     % text prod
\newcommand{\ttprod}{\Pi}          % mini prod
\newcommand{\tcoprod}{\textstyle{\coprod}} % text coprod
\newcommand{\ttcoprod}{\amalg}    % mini coprod
\newcommand{\tsum}{\textstyle{\sum}}       % text sum
\newcommand{\sprod}{\scriptstyle{\prod}}     % script prod
\newcommand{\scoprod}{\scriptstyle{\coprod}} % script coprod
\newcommand{\ssum}{\scriptstyle{\sum}}       % script sum
\newcommand{\ssprod}{\scriptscriptstyle{\prod}}     % scriptscript prod
\newcommand{\sscoprod}{\scriptscriptstyle{\coprod}} % scriptscript coprod
\newcommand{\sssum}{\scriptscriptstyle{\sum}}       % scriptscript sum
\newcommand{\bprod}{\times}        % binary product (prod)
\newcommand{\hprod}{\otimes}       % heterogeneous product
\newcommand{\bcoprod}{\amalg}      % binary coproduct (coprod)
\newcommand{\hcoprod}{\oplus}      % heterogeneous product
\newcommand{\inter}{\bigcap}       % family intersection
\newcommand{\binter}{\cap}         % binary intersection
\newcommand{\union}{\bigcup}       % family union
\newcommand{\bunion}{\cup}         % binary union
\newcommand{\iso}{\cong}           % isomorfo
\DeclareMathOperator{\card}{card}  % cardinal card(A)
\DeclareMathOperator{\fin}{f}      % Sub_{f}(X)
\DeclareMathOperator{\lfin}{lf}    % Sub_{lf}(X)
\newcommand{\function}[4]{
	\begin{array}{@{\:}c@{\:}c@{\:}l}
		#1 &\mor& #2 \\
		#3 &\longmapsto& #4
\end{array} }
\newcommand{\nfunction}[4]
{\left\{
	\function{#1}{#2}{#3}{#4}
	\right. }

\newcommand{\smod}[1]{\mathbin{\,(\op{m\acute{o}d.}{#1})}} % m\'{o}dulo
\DeclareMathOperator{\Sat}{Sat}    % Saturacion de un cjto por una equiv
\DeclareMathOperator{\unif}{unf}   % uniforme
\DeclareMathOperator{\ClSp}{ClSp}  % Functor ClSp
\DeclareMathOperator{\Ops}{Ops}    % Sistemas de abiertos
\DeclareMathOperator{\TOps}{TOps}  % Sistemas de abiertos topologicos
\DeclareMathOperator{\Cls}{Cls}    % Sistemas de cerrados (clausura)
\DeclareMathOperator{\TCls}{TCls}  % Sistemas de cerrados (clausura) topologicos

\DeclareMathOperator{\Inop}{Inop}  % Operadores interior
\DeclareMathOperator{\TInop}{TInop}% Operadores interior topologicos
\DeclareMathOperator{\Clop}{Clop}  % Operadores clausura
\DeclareMathOperator{\TClop}{TClop}% Operadores clausura topologicos

\DeclareMathOperator{\Inr}{Inr}    % Relaciones interior
\DeclareMathOperator{\TInr}{TInr}  % Relaciones interior topologicas
\DeclareMathOperator{\Clr}{Clr}    % Relaciones clausura
\DeclareMathOperator{\TClr}{TClr}  % Relaciones clausura topologicas

\DeclareMathOperator{\Opn}{Opn}    % Sistemas de entornos con abiertos
\DeclareMathOperator{\TOpn}{TOpn}  % Sistemas de entornos con abiertos topologicos
\DeclareMathOperator{\Cln}{Cln}    % Sistemas de entornos con cerrados
\DeclareMathOperator{\TCln}{TCln}  % Sistemas de entornos con cerrados topologicos

\DeclareMathOperator{\Clsg}{Clsg}  % Sistema clausura generado
\DeclareMathOperator{\Clopg}{Clopg}% Operador clausura generado

\DeclareMathOperator{\ACls}{ACls}  % Sistemas clausura algebraicos
\DeclareMathOperator{\AClop}{AClop}% Operadores clausura algebraico
\DeclareMathOperator{\UCls}{UCls}  % Sistemas clausura uniformes
\DeclareMathOperator{\UClop}{UClop}% Operadores clausura uniformes
\DeclareMathOperator{\UACls}{UACls}% Sistemas clausura algebraicos uniformes
\DeclareMathOperator{\UAClop}{UAClop}% Operadores clausura algebraico uniformes
\DeclareMathOperator{\Fix}{Fix}    % Puntos fijos del operador
\DeclareMathOperator{\Cons}{Cons}  % Consistentes
\DeclareMathOperator{\LcClop}{LcClop}% Operadores logicamente compactos
\DeclareMathOperator{\Lop}{Lop}    % Levantamiento optimal
\DeclareMathOperator{\Lcop}{Lcop}  % Levantamiento cooptimal
\newcommand{\sleq}{\sqsubseteq}    % square menor o igual
\newcommand{\sles}{\sqsubset}      %
\newcommand{\sgeq}{\sqsupseteq}    % square mayor o igual
\newcommand{\sgrt}{\sqsupset}      %
\DeclareMathOperator{\Ub}{\Uparrow}  % upper bounds cotas superiores (doble flecha) se puede redefinir como Ub
\DeclareMathOperator{\Lb}{\Downarrow}% lower bounds cotas inferiores (doble flecha) se puede redefinir como Db
\DeclareMathOperator{\ua}{\uparrow}  % final segment flecha arriba
\DeclareMathOperator{\da}{\downarrow}% initial segment flecha abajo
\DeclareMathOperator{\Ucl}{Ucl}      % Upward closed.
\DeclareMathOperator{\Dcl}{Dcl}      % Downward closed.
%\renewcommand{\sup}{\bigvee\nolimits}  % Supremos
%\renewcommand{\inf}{\bigwedge\nolimits}% Infimos
\newcommand{\bsup}{\vee}           % supremos binarios
\newcommand{\binf}{\wedge}         % infimos binarios
\DeclareMathOperator{\st}{st}      % st strong sub\'{i}ndice categor\'{i}as
\DeclareMathOperator{\Flt}{Flt}    % Filtros de un ordenado
\DeclareMathOperator{\Fg}{Fg}      % Filtro generado
\DeclareMathOperator{\Fp}{Fp}      % Filtros principales
\DeclareMathOperator{\Idl}{Idl}    % Ideales de un ordenado
\DeclareMathOperator{\Ig}{Ig}      % Ideal generado
\DeclareMathOperator{\Ip}{Ip}      % Ideales principales
\DeclareMathOperator{\Spc}{Spc}    % Espectro: inf-irreducibles
\DeclareMathOperator{\Max}{Max}    % Maximales
\DeclareMathOperator{\Ind}{Ind}    % Inductivos
\DeclareMathOperator{\Ird}{Ird}    % Irreducibles
\DeclareMathOperator{\Prim}{Prim}  % Primos
\DeclareMathOperator{\Ex}{Ex}      % Operador clausura sobre el conjunto de tipos $S$ asociado a un operador uniforme
\DeclareMathOperator{\arityd}{ar}     % ar(iedad)
\DeclareMathOperator{\car}{car}    % car (coariedad)
\DeclareMathOperator{\bard}{bar}   % biariedad
\DeclareMathOperator{\FMon}{FMon}  % Functor monoide libre
\DeclareMathOperator{\ard}{ar}     % ar(iedad)
\newcommand{\fmon}[1]
{\ensuremath{#1^{\star}}}
\newcommand{\ffmon}[1]
{\ensuremath{#1^{\star\star}}}
\newcommand{\fffmon}[1]
{\ensuremath{#1^{\star\star\star}}}
\newcommand{\bb}[1]{\ensuremath
	{\lvert #1 \rvert}}
\DeclareMathOperator{\Cl}{Cl}      % Closed subset of an algebra. Cerrados de un algebra.
\DeclareMathOperator{\Sg}{Sg}      % Sg subalgebra generada
\DeclareMathOperator{\E}{E}        % Operador E generacion de subalgebras
\DeclareMathOperator{\Cgr}{Cgr}    % Congruencia
\DeclareMathOperator{\Cg}{Cg}      % Congruencia generada
\DeclareMathOperator{\C}{C}        %
\DeclareMathOperator{\Ccg}{Ccg}    % Congruencia cogenerada
\newcommand{\FIC}{\op{Cgr}^{\op{fi}}} % Congruencia totalemente invariante (fully invariant)
\newcommand{\FICg}{\op{Cg}^{\op{fi}}} % Congruencia totalemente invariante (fully invariant)
\DeclareMathOperator{\Ker}{Ker}    % Kernel: n\'{u}cleo de un aplicaci\'{o}n
\newcommand{\pr}{\mathrm{pr}}        % proyecci\'{o}n can\'{o}nica de A en A/Phi
\DeclareMathOperator{\gen}{gen}    % para operaciones generalizadas
\DeclareMathOperator{\Op}{Op}      % Op_{w}(_A_) Algebra op. w-arias
\DeclareMathOperator{\Os }{O}     % O(_A_)=(Opw(A)s| w,s in ……)
\DeclareMathOperator{\Pol}{Pol}    % Pol_{w}(_A_) Op. polinomicas  w-arias
\DeclareMathOperator{\Pd}{Pd}      % Pd Polinomio determinado
\newcommand{\vs}[1]{\mathbin{\downarrow}#1}
\DeclareMathOperator{\Alg}{Alg}    % Alg_{w}(_A_) Op. algebraicas w-arias
\DeclareMathOperator{\W}{W}        % Words Algebra de las palabras
\DeclareMathOperator{\Fr}{Fr}      % Free Algebra libre
\DeclareMathOperator{\F}{F}        % Free para teorias, monadas, .. ..
\DeclareMathOperator{\G}{G}        % Ground functor olvido
\DeclareMathOperator{\I}{I}        % Inclusion functor inclusion
\DeclareMathOperator{\J}{J}        % Inclusion functor inclusion
\DeclareMathOperator{\p}{p}        % p: |w --> A ^ Aw
\DeclareMathOperator{\bconcat}
{\curlywedge}
\newcommand{\cncat}{\curlywedge}
\newcommand{\concat}
{\ensuremath{\text
		{\Large $\curlywedge$}}}
\newcommand{\ext}[1]
{\ensuremath{#1^{\sharp}}}
\newcommand{\coext}[1]
{\ensuremath{#1^{\flat}}}
\newcommand{\dext}[1]
{\ensuremath{#1^{\varkappa}}}
\newcommand{\nat}{\natural}        % becuadro
\newcommand{\adg}[1]
{\ensuremath{#1^{\dagger}}}
\DeclareMathOperator{\sub}{sub}    % subterminos de un t\'{e}rmino.
\DeclareMathOperator{\var}{var}    % variables de un termino. var(P)
\DeclareMathOperator{\Ter}{Ter}    % Terminos: Fr E (X)s X en Set^S, s en S
\DeclareMathOperator{\T}{T}        % Fr E (w)s w en S* s, en S
\DeclareMathOperator{\Eq}{Eq}      % Eq(_) Ecuaciones  sobre una signatura
% \models |=
\DeclareMathOperator{\Qeq}{Qeq}    % Qeq(_) Quasiecuaciones  sobre una signatura
\DeclareMathOperator{\finS}{-\mathrm{f}}
\DeclareMathOperator{\B}{B}        % Base de Dedekind-Peano
\DeclareMathOperator{\ev}{ev}      % evaluaci\'{o}n ev: A^B x B -> A
\DeclareMathOperator{\Var}{Var}    % Variedad Var(\sg{\Sigma}) conjunto de las variedades
\DeclareMathOperator{\Th}{Th}      % Teor\'{i}a generada por una clase de algebras
\DeclareMathOperator{\Mod}{Mod}    % Clase de algebras determinada por ecuaciones
\DeclareMathOperator{\Cn}{Cn}      % Operador clausura asociado a Th o Mod
\DeclareMathOperator{\Ec}{Ec}      % Operador clausura asociado a Mod o Th
\newcommand{\Hall}{\ensuremath{\mathrm{H}}} % ecuaciones Hall
\newcommand{\Ben}{\ensuremath{\mathrm{B}}}  % ecuaciones Lawvere
\DeclareMathOperator{\DP}{DP}      % condiciones Dedekind-Peano
\DeclareMathOperator{\Mon}{Mon}    % Especificaci\'{o}n de Monoide.
\DeclareMathOperator{\der}{der}    % derivors
\DeclareMathOperator{\fuj}{fuj}    % fujiwara
\DeclareMathOperator{\Sen}{Sen}    % Sentencias
\DeclareMathOperator{\Mnd}{Mnd}    % Categoria de Monadas.
\DeclareMathOperator{\Md}{Md}      % Functor de Ad en Mnd.
\DeclareMathOperator{\cp}{cp}      % choosen proyections. Proyecciones especificadas.
\newcommand{\rotc}[2]{\rotatebox[origin=c]{#1}{\ensuremath{#2}}}
\newcommand{\sx}[2]{(#1\mid#2)}
\newcommand{\Ntr}{\ensuremath{\mathrm{N}}}
\newcommand{\cnys}{condici\'{o}n necesaria y suficiente }
\newcommand{\vacio}{\ensuremath{\varnothing}}
\newcommand{\htg}{heterog\'{e}nea }
\newcommand{\HRP}{\ensuremath
	{\mathbf{H}^{\mathrm{rp}}(\mathbb{N}^{\text{\Large{$\cdot$}}},\mathbb{N})}}
\newcommand{\hrp}{\ensuremath
	{\mathrm{H}^{\mathrm{rp}}(\mathbb{N}^{\text{\Large{$\cdot$}}},\mathbb{N})}}
\newcommand{\HPR}{\ensuremath
	{\mathbf{H}^{\mathrm{pr}}(\mathbb{N}^{\text{\Large{$\cdot$}}},\mathbb{N})}}
\newcommand{\hpr}{\ensuremath
	{\mathrm{H}^{\mathrm{pr}}(\mathbb{N}^{\text{\Large{$\cdot$}}},\mathbb{N})}}
\newcommand{\HRG}{\ensuremath
	{\mathbf{H}^{\mathrm{rg}}(\mathbb{N}^{\text{\Large{$\cdot$}}},\mathbb{N})}}
\newcommand{\hrg}{\ensuremath
	{\mathrm{H}^{\mathrm{rg}}(\mathbb{N}^{\text{\Large{$\cdot$}}},\mathbb{N})}}

\newcommand{\homo}{homomorfismo}
\newcommand{\Law}{\ensuremath{\mathsf{Law}(\underline{\Sigma},V)}}
\newcommand{\ul}{\underline}
\newcommand{\ol}{\overline}
\newcommand{\wh}{\widehat}
\newcommand{\bcal}[1]{\boldsymbol{\mathcal{#1}}}
\newcommand{\brel}{\ensuremath{\xymatrix{{}\arity@{{*}{-}{*}}[r] & {}}}}
\newcommand{\ulob}{\ensuremath{\overset{\boldsymbol{.}}{}}}
\newcommand{\dlob}{\ensuremath{\boldsymbol{.}}}
\newcommand{\nseq}[3]{\xymatrix@1@C=16pt{#1 \arity@{>}[r]_-{\scriptscriptstyle{#2}} & #3 }}

\NoCompileMatrices              % compila los diagramas \NoCompileMatrices
\def\labelstyle{\textstyle}     % tama\~{n}o de las decor(ations)
\def\twocellstyle{\textstyle}   % tama\~{n}o para twocell feature
\xymatrixrowsep = {8ex}         % o 40pt
\xymatrixcolsep = {10ex}        % o 50pt
\newcommand{\Bigmatrix}{% Grande @C=12.0ex@R=10.0ex
	\xymatrixcolsep={12ex}
	\xymatrixrowsep={10ex}}
\newcommand{\Normalmatrix}{% Normal @C=10ex@R=8ex
	\xymatrixcolsep={10ex}
	\xymatrixrowsep={8ex}}
\newcommand{\Smallmatrix}{% Peque\~{n}o @C=8ex@R=6.6ex
	\xymatrixcolsep={8ex}
	\xymatrixrowsep={6.6ex}}
\newcommand{\Smmallmatrix}{% Peque\~{n}ito @C=7ex@R=5.6ex
	\xymatrixcolsep={7ex}
	\xymatrixrowsep={5.6ex}}
\newcommand{\Smmmallmatrix}{% Peque\~{n}isimo @C=6ex@R=4.6ex
	\xymatrixcolsep={6ex}
	\xymatrixrowsep={4.6ex}}
\newcommand{\Smmmmallmatrix}{% Minusculo @C=5ex@R=3.6ex
	\xymatrixcolsep={5ex}
	\xymatrixrowsep={3.6ex}}

\newdir{<= }{:a(+25)\dir{-}*:a(-25)\dir{-}*!/:a(85)+1.8pt/:a(-25)\dir{-}}
\newdir{ = }{*@{=}}
\newdir{+>}{@{|}*@{-}!/-6.5pt/@{>}*!/-13pt/{}}
\newdir{ +}{{}*!/-5pt/@{+}}
\newdir{ >}{{}*!/-10pt/@{>}}
\newdir{.-}{{}*:a(-90)h!/2.1pt/{.}}
\newdir{.->}{:a(-90)h!/2.5pt/{.}*!/4pt/\dir{-}*!/0pt/\dir{-}*%
	!/-2pt/\dir{-}*!/-8pt/\dir{>}}
\newdir{-.->}%
{:a(-90)h!/2pt/{.}*!/8pt/\dir{-}*!/4pt/\dir{-}*!/0pt/\dir{-}*%
	!/-5pt/\dir{-}*!/-11pt/\dir{>}}
\newdir{~>}{!/4.6pt/\dir{~}*!/1pt/\dir{-}*!/-5pt/\dir{>}}
\newdir{~~>}{!/7pt/\dir{~}*!/0pt/\dir{~}*!/-5pt/\dir{>}}
\newdir{=~>}{*!/1pt/@2{~}*!/6pt/{}*!/-6pt/@2{>}}
\newdir{<~~>}{!/21pt/\dir{<}*!/21pt/\dir{-}!/12pt/\dir{~}!/5pt/\dir{~}*!/1pt/\dir{-}*!/-3pt/\dir{>}}
\newdir{=>}{!/5pt/\dir{=}!/2.5pt/\dir{=}*!/-5pt/\dir2{>}*!/7pt/\dir{ }}
\newdir{==>}{!/-2pt/\dir{=}!/-6pt/\dir{=}%
	!/-10pt/\dir{=}*!/-18pt/\dir2{>}}
\newdir{<==>}{!\dir2{<}!/-2pt/\dir{=}!/-6pt/\dir{=}%
	!/-10pt/\dir{=}*!/-17pt/\dir2{>}}
\newdir{< }{{}*!/12pt/\dir{<}}
\newdir{-o}{!/1.3pt/{}*!/0pt/{}@{o}*!/-5pt/{}}

\newsavebox{\xymor}  % --->
\newsavebox{\xymon}  % +-->
\newsavebox{\xyepi}  % --+>
\newsavebox{\xytn}   % -.->
\newsavebox{\xyrel}  % ---o
\newsavebox{\xycel}  % ===>
\newsavebox{\xymdf}  % ~~~>
\newsavebox{\xyumor} % ---`
\newsavebox{\xydmor} % ---,
\newsavebox{\xyomor} % ---<
\newsavebox{\xyemor} % >--->

\newcommand{\xynode}{\makebox[0ex]{}}
\savebox{\xymor}{\ensuremath{%
		\xymatrix@1@C=19pt{\xynode \ar@{>}[r] & \xynode }}}
\savebox{\xymon}{\ensuremath{%
		\xymatrix@1@C=19pt{\xynode \ar@{{ +}{-}{>}}[r] & \xynode }}}
\savebox{\xyepi}{\ensuremath{%
		\xymatrix@1@C=19pt{\xynode \ar@{{}{-}{+>}}[r] & \xynode }}}
\savebox{\xytn}{\ensuremath{%
		\xymatrix@1@C=19pt{\xynode \ar[r]|(.44){\object@{.-}} & \xynode
}}}
\savebox{\xyrel}{\ensuremath{%
		\xymatrix@1@C=19pt{\xynode \ar@{{}{-}{-o}}[r] & \xynode }}}
\savebox{\xycel}{\ensuremath{%
		\xymatrix@1@C=19pt{\xynode \ar@{=>}[r] & \xynode }}}
\savebox{\xymdf}{\ensuremath{%
		\xymatrix@1@C=16pt{\xynode \ar@{}[r]|{\dir{~>}} & \xynode}}}
\savebox{\xyumor}{\ensuremath{%
		\xymatrix@1@C=19pt{\xynode \ar@{{}{-}^{>}}[r] & \xynode }}}
\savebox{\xydmor}{\ensuremath{%
		\xymatrix@1@C=19pt{\xynode \ar@{{}{-}_{>}}[r] & \xynode }}}
\savebox{\xyomor}{\ensuremath{%
		\xymatrix@1@C=19pt{\xynode \ar@{{}{-}^{< }}[r] & \xynode }}}
\savebox{\xyemor}{\ensuremath{%
		\xymatrix@1@C=19pt{\xynode \ar@{{ >}{-}{>}}[r] & \xynode }}}

\newcommand{\mor}{\usebox{\xymor}}    % mor  --->
\newcommand{\mon}{\usebox{\xymon}}    % mon  +-->
\newcommand{\epi}{\usebox{\xyepi}}    % epi  --+>
\newcommand{\tn}{\usebox{\xytn}}      % tn   -.->
\newcommand{\rel}{\usebox{\xyrel}}    % rel  ---o
\newcommand{\cel}{\usebox{\xycel}}    % cel  ===>
\newcommand{\mdf}{\usebox{\xymdf}}    % mdf  ~~~>
\newcommand{\umor}{\usebox{\xyumor}}  % umor ---`
\newcommand{\dmor}{\usebox{\xydmor}}  % dmor ---,
\newcommand{\omor}{\usebox{\xyomor}}  % omor ---<
\newcommand{\emor}{\usebox{\xyemor}}  % emor >--->

\newcommand{\xyn}[1]{\save="#1"\restore}
\newcommand{\xymn}[3][0ex,5ex]{%
	\save[]+<#1>*+{#2}="#3" \restore }
\newcommand{\functor}[9]{
	\xymatrix{
		#4 \save[]+<0ex,5ex>*+{#1}="1"  \restore
		\arity[d]_{#6}  \arity@{}[rd]|{\longmapsto}
		& #5 \save[]+<0ex,5ex>*+{#3}="3"  \restore
		\arity[d]^{#7}
		\\
		#8 & #9 \arity "1";"3"^-{#2} } }
\newcommand{\functornd}[9]{
	\xymatrix{
		#4 \save[]+<0ex,5ex>*+{#1}="1"  \restore
		\arity[d]_{#6}  \arity@{}[rd]|{\longmapsto}
		& #5 \save[]+<0ex,5ex>*+{#3}="3"  \restore
		\\
		#8 & #9 \arity[u]_{#7} \arity "1";"3"^-{#2} } }
\newcommand{\functordn}[9]{
	\xymatrix{
		#4 \save[]+<0ex,5ex>*+{#1}="1"  \restore
		\arity@{}[rd]|{\longmapsto}
		& #5 \save[]+<0ex,5ex>*+{#3}="3"  \restore
		\arity[d]^{#7}
		\\
		#8  \arity[u]^{#6}  & #9 \arity "1";"3"^-{#2} } }

\newcommand{\larr}{->}
\newcommand{\rarr}{->}
\newcommand{\dirarrows}[2]{\renewcommand{\larr}{#1}\renewcommand{\rarr}{#2}}
\newcommand{\xfunctor}[9]{
	\xymatrix{
		#4 \save[]+<0ex,5ex>*+{#1}="1"  \restore
		\ifthenelse{\equal{\larr}{->}}{\arity[d]_{#6}}{}
		\ifthenelse{\equal{\larr}{<-}}{\arity[d];[]^{#6}}{}
		\ifthenelse{\equal{\larr}{-<}}{\arity@{< }[d]_{#6}}{}
		\arity@{}[rd]|{\longmapsto}
		& #5 \save[]+<0ex,5ex>*+{#3}="3"  \restore
		\ifthenelse{\equal{\rarr}{->}}{\arity[d]^{#7}}{}
		\ifthenelse{\equal{\rarr}{<-}}{\arity[d];[]_{#7}}{}
		\ifthenelse{\equal{\rarr}{-<}}{\arity@{< }[d]^{#7}}{}
		\\
		#8 & #9 \arity "1";"3"^-{#2} } }

% End of macros
%\input{climent_macros}
%\input{climent_Xymacros}

\theoremstyle{plain}
\newtheorem{theorem}{Theorem}[section]
\newtheorem{proposition}[theorem]{Proposition}
\newtheorem{classification}[theorem]{Classification}
\newtheorem{corollary}[theorem]{Corollary}
\newtheorem{lemma}[theorem]{Lemma}
\newtheorem{remark}[theorem]{Remark}

\theoremstyle{definition}
\newtheorem{definition}[theorem]{Definition}
\newtheorem{example}[theorem]{Example}
\newtheorem{examples}{\bf Examples}

\newtheorem{assumption}{\bf Assumption}
\newtheorem{assumptions}{\bf Assumptions}
\newtheorem{exercise}[theorem]{Exercise}
\newtheorem*{conventions}{Conventions}
\newtheorem*{convention}{Convention}

\begin{document}

\title{Higher-order categorial algebras}

\author[Climent]{J. Climent Vidal}
\address{Universitat de Val\`{e}ncia\\
         Departament de L\`{o}gica i Filosofia de la Ci\`{e}ncia
         }
\email{Juan.B.Climent@uv.es}

\author[Cosme]{E. Cosme Ll\'{o}pez}
\address{Universitat de Val\`{e}ncia\\
         Departament de Matem\`{a}tiques\\ and 
         Nantong University\\ %南通大学
         School of Mathematics and Statistics
         }
\email{Enric.Cosme@uv.es}

\author[Ruiz]{R. Ruiz Mora}
\address{Universitat de Val\`{e}ncia\\
         Departament de Matem\`{a}tiques
         }
\email{Raul.Ruiz-Mora@uv.es}

\author[Vago]{M. Vago Marzal}

\maketitle

\section{Preliminaries}

\subsection{Partial algebras and QE-equations}

Throughout, $S$ is a set of sorts, $S^{\star}$ the set of words on $S$, and, for an $S$-sorted set $A$ and a word $\mathbf{s}=(s_{j})_{j\in\bb{\mathbf{s}}}\in S^{\star}$, $A_{\mathbf{s}}=\prod_{j\in\bb{\mathbf{s}}}A_{s_{j}}$. Sections~\ref{SCompact} and~\ref{SExtended} present higher-order categorial algebras as models of sets of $\mathrm{QE}$-equations between partial algebras. We fix here these notions, following~\cite{pb86, pb02}.

\begin{definition}\label{DPartialAlg}
Let $\Sigma$ be an $S$-sorted signature. A \emph{partial $\Sigma$-algebra} is a pair $\mathbf{A}=(A,F)$ where $A$ is an $S$-sorted set and $F$ assigns to every $\sigma\in\Sigma_{\mathbf{s},s}$ a partial mapping $F_{\sigma}\colon A_{\mathbf{s}}\dmor A_{s}$. A \emph{homomorphism} from $\mathbf{A}=(A,F)$ to $\mathbf{B}=(B,G)$ is an $S$-sorted mapping $h\colon A\mor B$ such that, for every $\sigma\in\Sigma_{\mathbf{s},s}$ and every $(a_{j})_{j\in\bb{\mathbf{s}}}\in\mathrm{Dom}(F_{\sigma})$, the family $(h_{s_{j}}(a_{j}))_{j\in\bb{\mathbf{s}}}$ is in $\mathrm{Dom}(G_{\sigma})$ and
\[
h_{s}(F_{\sigma}((a_{j})_{j\in\bb{\mathbf{s}}}))=G_{\sigma}((h_{s_{j}}(a_{j}))_{j\in\bb{\mathbf{s}}}).
\]
Homomorphisms preserve, but need not reflect, definedness.
\end{definition}

\begin{definition}\label{DQEEq}
Let $\Sigma$ be an $S$-sorted signature and $V$ an $S$-sorted set of variables, and let $\mathrm{T}_{\Sigma}(V)$ be the $S$-sorted set of terms. For a partial $\Sigma$-algebra $\mathbf{A}$ and a valuation $v\colon V\mor A$, the value $t^{\mathbf{A}}[v]$ of a term $t$ is partial: $x^{\mathbf{A}}[v]=v(x)$ for a variable $x$, and $\sigma((t_{j})_{j})^{\mathbf{A}}[v]$ is defined if and only if every $t_{j}^{\mathbf{A}}[v]$ is defined and $(t_{j}^{\mathbf{A}}[v])_{j}\in\mathrm{Dom}(F_{\sigma})$, in which case it is $F_{\sigma}((t_{j}^{\mathbf{A}}[v])_{j})$.

An \emph{existence equation} is a pair of terms $t,t'\in\mathrm{T}_{\Sigma}(V)_{s}$ of the same sort, written $t\overset{\mathrm{e}}{=}t'$; it \emph{holds} in $\mathbf{A}$ under $v$ if $t^{\mathbf{A}}[v]$ and $t'^{\mathbf{A}}[v]$ are both defined and are equal. In particular, $t\overset{\mathrm{e}}{=}t$ states that $t$ is defined. A \emph{$\mathrm{QE}$-equation} is an expression
\[
\textstyle\bigwedge_{i\in m}(t_{i}\overset{\mathrm{e}}{=}t'_{i})\;\longrightarrow\;t\overset{\mathrm{e}}{=}t',
\]
with $m\in\mathbb{N}$; it \emph{holds} in $\mathbf{A}$ if, for every valuation $v\colon V\mor A$ under which every premise $t_{i}\overset{\mathrm{e}}{=}t'_{i}$ holds, the conclusion $t\overset{\mathrm{e}}{=}t'$ holds. We denote by $\mathrm{QE}(\Sigma)_{V}$ the set of all $\mathrm{QE}$-equations with variables in $V$.

A \emph{$\mathrm{QE}$-specification} is a quadruple $(S,\Sigma,V,\mathcal{E})$ with $\mathcal{E}\subseteq\mathrm{QE}(\Sigma)_{V}$. Its models, the partial $\Sigma$-algebras in which every equation of $\mathcal{E}$ holds, form the \emph{$\mathrm{QE}$-variety} $\mathcal{V}(\mathcal{E})$, and, together with the homomorphisms of Definition~\ref{DPartialAlg}, the category $\mathsf{PAlg}(\mathcal{E})$.
\end{definition}

\section{Derivors}

In this subsection, after defining the variety of Hall algebras, we introduce the notion of derivor between many-sorted signatures. This will allow us, using the homomorphisms between Hall algebras, to obtain $\mathsf{Sig}_{\mathfrak{d}}$, the category of many-sorted signatures and derivors. Next, after defining a suitable contravariant functor from $\mathsf{Sig}_{\mathfrak{d}}$ to $\mathsf{Cat}$, the category of $\boldsymbol{\mathcal{U}}$-locally small categories and functors, we obtain, by means of the Grothendieck construction, $\mathsf{Alg}_{\mathfrak{d}}$, the category with objects the ordered pairs $((S,\Sigma),(A,F))$, with $(S,\Sigma)$ a many-sorted signature and $(A,F)$ a $\Sigma$-algebra, and morphisms from $((S,\Sigma),(A,F))$ to  $((T,\Lambda),(B,G))$ the ordered pairs $((\varphi,d),f)$, with $(\varphi,d)$ a derivor from $(S,\Sigma)$ to $(T,\Lambda)$ and $f$ a homomorphism of $\Sigma$-algebras from $(A,F)$ to $(B_{\varphi},G^{(\varphi,d)})$, a canonically derived algebra of $(B,G)$. 

Before defining the notion of Hall algebra, we recall that
\begin{enumerate}
	\item a finitary specification is an ordered triple $(S,\Sigma,\mathcal{E})$, where $S$ is a set of sorts, $\Sigma$ an $S$-sorted signature, and $\mathcal{E}\subseteq\mathrm{Eq}_{V^{S}}(\Sigma)=
	(\mathrm{T}_{\Sigma}(X)_{s}^{2})_{(X,s)\in \mathrm{Sub}_{\mathrm{f}}(V^{S})\times S}$, i.e., a set of finitary  $\Sigma$-equations, where $V^{S}$, the $S$-sorted set of the variables, is a fixed $S$-countably infinite $S$-sorted set; and that
	\item a $\Sigma$-algebra $\mathbf{A}$ is a $(S,\Sigma,\mathcal{E})$-algebra if $\mathbf{A}\models^{\Sigma} \mathcal{E}$, i.e., if, for every $(X,s)\in \mathrm{Sub}_{\mathrm{f}}(V^{S})\times S$ and every $(P,Q)\in \mathcal{E}_{X,s}$, $\mathbf{A}$ is a model of $(P,Q)$, in symbols $\mathbf{A}\models^{\Sigma}_{X,s} (P,Q)$, which in turn means that, for every $a\in \mathrm{Hom}(X,A)$, $a^{\sharp}_{s}(P) = a^{\sharp}_{s}(Q)$.
\end{enumerate}

\begin{definition}
	Let $S$ be a set of sorts and let $V^{\mathrm{H}_{S}} = (V^{\mathrm{H}_{S}}_{\mathbf{s},s})_{(\mathbf{s},s)\in S^{\star}\times S}$ be a fixed $S^{\star}\times S$-sorted set of variables with a countably infinity of variables in each component, say $V^{\mathrm{H}_{S}}_{\mathbf{s},s} = \{\,v^{\mathbf{s},s}_{n}\mid n\in \mathbb{N}\,\}$ (which is isomorphic to 
	$\{(\mathbf{s},s)\}\times\mathbb{N}$), for every $(\mathbf{s},s)\in S^{\star}\times S$.  A \emph{Hall algebra for} $S$ is a $\mathrm{H}_{S} = (S^{\star}\times S,\Sigma^{\mathrm{H}_{S}},\mathcal{E}^{\mathrm{H}_{S}})$-algebra, where $\Sigma^{\mathrm{H}_{S}}$ is the $S^{\star}\times S$-sorted signature, i.e., the $(S^{\star}\times S)^{\star}\times (S^{\star}\times S)$-sorted set, defined as follows:
	\begin{enumerate}
		\item[$\mathrm{HS}_{1}$.] For every $\mathbf{s}\in S^{\star}$ and $i\in\bb{\mathbf{s}}$,
		$$
		\pi^{\mathbf{s}}_{i}\colon\lambda\mor(\mathbf{s},s_{i}),
		$$
		where $\bb{\mathbf{s}}$ is the \emph{length} of the word $\mathbf{s}$ and $\lambda$ the
		\emph{empty word} in the underlying set of the free monoid on
		$S^{\star}\times S$.
		
		\item[$\mathrm{HS}_{2}$.] For every $\mathbf{r}$, $\mathbf{s}\in S^{\star}$ and $s\in S$,
		$$
		\xi_{\mathbf{r},\mathbf{s},s}\colon ((\mathbf{s},s),(\mathbf{r},s_{0}),\ldots,(\mathbf{r},s_{\bb{\mathbf{s}}-1}))\mor(\mathbf{r},s);
		$$
	\end{enumerate}
	while $\mathcal{E}^{\mathrm{H}_{S}}$ is the sub-$(S^{\star}\times
	S)^{\star}\times (S^{\star}\times S)$-sorted set of
	$\mathrm{Eq}(\Sigma^{\mathrm{H}_{S}})$, where
	$$
	\mathrm{Eq}\left(\Sigma^{\mathrm{H}_{S}}\right) =
	\left(\mathrm{T}_{\Sigma^{\mathrm{H}_{S}}}(\vs{\ol{\mathbf{s}}})_{(\mathbf{r},s)}^{2}\right)_{(\ol{\mathbf{s}},(\mathbf{r},s))\in
		(S^{\star}\times S)^{\star}\times (S^{\star}\times S)},
	$$
	defined as follows:
	\begin{enumerate}
		\item[$\mathrm{H}_{1}$.] \emph{Projection}.
		For every $\mathbf{r}$, $\mathbf{s}\in S^{\star}$ and $i\in\bb{\mathbf{s}}$, the equation
		$$
		\xi_{\mathbf{r},\mathbf{s},s_{i}}(\pi^{\mathbf{s}}_{i},v^{\mathbf{r},s_{0}}_{0},\ldots,
		v^{\mathbf{r},s_{\bb{\mathbf{s}}-1}}_{\bb{\mathbf{s}}-1})=v^{\mathbf{r},s_{i}}_{i}
		$$
		of type
		$(((\mathbf{r},s_{0}),\ldots,(\mathbf{r},s_{\bb{\mathbf{s}}-1})),(\mathbf{r},s_{i})).$
		\item[$\mathrm{H}_{2}$.] \emph{Identity}.
		For every $\mathbf{r}\in S^{\star}$ and $j\in \bb{\mathbf{r}}$, the equation
		$$
		\xi_{\mathbf{r},\mathbf{r},r_{j}}(v^{\mathbf{r},r_{j}}_{j},\pi^{\mathbf{r}}_{0},\ldots,\pi^{\mathbf{r}}_{\bb{\mathbf{r}}-1})=
		v^{\mathbf{r},r_{j}}_{j}
		$$
		of type
		$(((\mathbf{r},r_{j})),(\mathbf{r},r_{j})).$
		\item[$\mathrm{H}_{3}$.] \emph{Associativity}.
		For every $\mathbf{q}$, $\mathbf{r}$, $\mathbf{s}\in S^{\star}$ and $s\in S$, the
		equation
		\begin{multline*}
			\xi_{\mathbf{q},\mathbf{r},s}(
			\xi_{\mathbf{r},\mathbf{s},s}(v^{\mathbf{s},s}_{0},v^{\mathbf{r},s_{0}}_{1},\ldots,
			v^{\mathbf{r},s_{\bb{\mathbf{s}}-1}}_{\bb{\mathbf{s}}}),
			v^{\mathbf{q},r_{0}}_{\bb{\mathbf{s}}+1},
			\ldots,v^{\mathbf{q},r_{\bb{\mathbf{r}}-1}}_{\bb{\mathbf{s}}+\bb{\mathbf{r}}}) = \\
			\xi_{\mathbf{q},\mathbf{s},s}(v^{\mathbf{s},s}_{0},
			\xi_{\mathbf{q},\mathbf{r},s_{0}}(v^{\mathbf{r},s_{0}}_{1},v^{\mathbf{q},r_{0}}_{\bb{\mathbf{s}}+1},
			\ldots,v^{\mathbf{q},r_{\bb{\mathbf{r}}-1}}_{\bb{\mathbf{s}}+\bb{\mathbf{r}}}),
			\\
			\ldots, 
			\xi_{\mathbf{q},\mathbf{r},s_{\bb{\mathbf{s}}-1}}(v^{\mathbf{r},s_{\bb{\mathbf{s}}-1}}_{\bb{\mathbf{s}}},
			v^{\mathbf{q},r_{0}}_{\bb{\mathbf{s}}+1},
			\ldots,v^{\mathbf{q},r_{\bb{\mathbf{r}}-1}}_{\bb{\mathbf{s}}+\bb{\mathbf{r}}}))
		\end{multline*}
		of type
		$(((\mathbf{s},s),(\mathbf{r},s_{0}),\ldots,(\mathbf{r},s_{\bb{\mathbf{s}}-1}),
		(\mathbf{q},r_{0}),\ldots,(\mathbf{q},r_{\bb{\mathbf{r}}-1})),(\mathbf{q},s)).$
	\end{enumerate}
	
	We call the formal constants $\pi^{\mathbf{s}}_{i}$ \emph{projections}, and the formal operations $\xi_{\mathbf{r},\mathbf{s},s}$ \emph{substitution operators}. Furthermore, we will denote by $\mathsf{Alg}(\mathrm{H}_{S})$ the category of Hall algebras for $S$ and homomorphisms between Hall algebras.  Since $\mathsf{Alg}(\mathrm{H}_{S})$ is a variety, the forgetful functor $\mathrm{G}_{\mathrm{H}_{S}}$ from $\mathsf{Alg}(\mathrm{H}_{S})$ to $\mathsf{Set}^{S^{\star}\times S}$ has a left adjoint $\mathbf{T}_{\mathrm{H}_{S}}$, situation denoted by $\mathbf{T}_{\mathrm{H}_{S}}\dashv \mathrm{G}_{\mathrm{H}_{S}}$, or diagrammatically by
	$$
	\xymatrix@=5pc{ \mathsf{Alg}(\mathrm{H}_{S})
		\ar@<1.5ex>[r]^{\mathrm{G}_{\mathrm{H}_{S}}} \ar@{}[r]|{\uadj} &
		\mathsf{Set}^{S^{\star}\times S}
		\ar@<1.5ex>[l]^{\mathbf{T}_{\mathrm{H}_{S}}} }
	$$
	which assigns to an $S^{\star}\times S$-sorted set $\Sigma$ the corresponding free Hall algebra
	$\mathbf{T}_{\mathrm{H}_{S}}(\Sigma)$.
\end{definition}

\begin{remark}
	The category $\mathsf{Alg}(\mathrm{H}_{S})$, of Hall algebras for $S$, is an algebraic version of many-sorted clone theory.
\end{remark}

\begin{remark}
	From $\mathrm{H}_{3}$, for $\mathbf{s}=\lambda$, the empty word on $S$, we get the invariance of constant functions axiom in~\cite{gm85}: For every $\mathbf{q}$, $\mathbf{r}\in S^{\star}$ and $s\in S$, we have the equation
	$$
	\xi_{\mathbf{q},\mathbf{r},s}(\xi_{\mathbf{r},\lambda,s}(v^{\lambda,s}_{0}),v^{\mathbf{q},r_{0}}_{1},
	\ldots,v^{\mathbf{q},r_{\bb{\mathbf{r}}-1}}_{\bb{\mathbf{r}}})=
	\xi_{\mathbf{q},\lambda,s}(v^{\lambda,s}_{0})
	$$
	of type $(((\lambda,s),(\mathbf{q},r_{0}),\ldots,(\mathbf{q},r_{\bb{\mathbf{r}}-1})),(\mathbf{q},s))$.
\end{remark}

For every $S$-sorted set $A$, $\mathrm{O}_{\mathrm{H}_{S}}(A) = (\mathrm{Hom}(A_{\mathbf{s}},A_{s}))_{(\mathbf{s},s)\in S^{\star}\times S}$, the $S^{\star}\times S$-sorted set of operation for $A$, is naturally equipped with a structure of Hall algebra, as stated in the following proposition, if we realize the projections as the true projections and the substitution operators as the generalized composition of mappings.





\begin{proposition}
	Let $A$ be an $S$-sorted set and $\mathbf{O}_{\mathrm{H}_{S}}(A)$ the $\Sigma^{\mathrm{H}_{S}}$-algebra with underlying many-sorted set $\mathrm{O}_{\mathrm{H}_{S}}(A)$ and algebraic structure defined as follows:
	\begin{enumerate}
		\item For every $\mathbf{s}\in S^{\star}$ and $i\in\bb{\mathbf{s}}$,
		$(\pi^{\mathbf{s}}_{i})^{\mathbf{O}_{\mathrm{H}_{S}}(A)} =
		\mathrm{pr}^{A}_{\mathbf{s},i}\colon A_{\mathbf{s}}\mor A_{s_{i}}$.
		
		\item For every $\mathbf{r},\mathbf{s}\in S^{\star}$ and $s\in S$,
		$\xi_{\mathbf{r},\mathbf{s},s}^{\mathbf{O}_{\mathrm{H}_{S}}(A)}$ is defined, for every $f\in
		A_{s}^{ A_{\mathbf{s}} }$ and $g\in A^{A_{\mathbf{r}}}_{\mathbf{s}}$, as
		$\xi_{\mathbf{r},\mathbf{s},s}^{\mathbf{O}_{\mathrm{H}_{S}}(A)}(f,g_{0},\ldots,g_{\bb{\mathbf{s}}-1})
		= f\circ\langle g_{i}\rangle_{i\in\bb{\mathbf{s}} }$, where $\langle
		g_{i}\rangle_{i\in\bb{\mathbf{s}}}$ is the unique mapping from $A_{\mathbf{r}}$ to
		$A_{\mathbf{s}}$ such that, for every $i\in \bb{\mathbf{s}}$, we have that
		$$
		\mathrm{pr}^{A}_{\mathbf{s},i}\circ\langle g_{i}\rangle_{i\in\bb{\mathbf{s}}} = g_{i}.
		$$
	\end{enumerate}
	Then $\mathbf{O}_{\mathrm{H}_{S}}(A)$ is a Hall algebra, the \emph{Hall algebra for} $(S,A)$.
\end{proposition}

\begin{remark}
	The closed sets of the Hall algebra $\mathbf{O}_{\mathrm{H}_{S}}(A)$ for $(S,A)$ are precisely the clones of (many-sorted) operations on the $S$-sorted set $A$. On the other hand, every $\Sigma$-algebra $\mathbf{A}$ has associated a Hall algebra. In fact, it suffices to consider $\mathbf{O}_{\mathrm{H}_{S}}(A)$, denoted by $\mathbf{O}_{\mathrm{H}_{S}}(\mathbf{A})$. Moreover, the finitary derived operations on $\mathbf{A}$ is a subalgebra of the Hall algebra $\mathbf{O}_{\mathrm{H}_{S}}(\mathbf{A})$.
\end{remark}



For every $S$-sorted signature $\Sigma$, $\mathrm{Tm}_{\mathrm{H}_{S}}(\Sigma) = (\mathrm{T}_{\Sigma}(\vs{\mathbf{s}})_{s})_{(\mathbf{s},s)\in S^{\star}\times S}$ is also e\-quipped with a structure of Hall algebra that formalizes the concept of substitution as stated in the following proposition.

\begin{proposition}
	Let $\Sigma$ be an $S$-sorted signature and $\mathbf{Tm}_{\mathrm{H}_{S}}(\Sigma)$ the $\Sigma^{\mathrm{H}_{S}}$-algebra
	with underlying many-sorted set $\mathrm{Tm}_{\mathrm{H}_{S}}(\Sigma)$ and algebraic structure defined as follows:
	\begin{enumerate}
		\item For every $\mathbf{s}\in S^{\star}$ and $i\in\bb{\mathbf{s}}$,
		$(\pi^{\mathbf{s}}_{i})^{\mathbf{Tm}_{\mathrm{H}_{S}}(\Sigma)}$
		is the image under $\eta^{\vs{\mathbf{s}}}_{s_{i}}$ of
		the variable $v_{i}^{s_{i}}$, where $\eta^{\vs{\mathbf{s}}} =
		(\eta^{\vs{\mathbf{s}}}_{s})_{s\in S}$ is the canonical embedding
		of $\vs{\mathbf{s}}$ into $\mathrm{T}_{\Sigma}(\vs{\mathbf{s}})$. Sometimes,
		to abbreviate, we will write $\pi^{\mathbf{s}}_{i}$ instead of
		$(\pi^{\mathbf{s}}_{i})^{\mathbf{Tm}_{\mathrm{H}_{S}}(\Sigma)}$.
		
		\item For every $\mathbf{r},\mathbf{s}\in S^{\star}$ and $s\in S$,
		$\xi_{\mathbf{r},\mathbf{s},s}^{\mathbf{Tm}_{\mathrm{H}_{S}}(\Sigma)}$
		is the mapping
		$$\xi_{\mathbf{r},\mathbf{s},s}^{\mathbf{Tm}_{\mathrm{H}_{S}}(\Sigma)}\nfunction
		{\mathrm{T}_{\Sigma}(\vs{\mathbf{s}})_{s} \times
			\mathrm{T}_{\Sigma}(\vs{\mathbf{r}})_{s_{0}} \times \cdots \times
			\mathrm{T}_{\Sigma}(\vs{\mathbf{r}})_{s_{\bb{\mathbf{s}}-1}}}
		{\mathrm{T}_{\Sigma}(\vs{\mathbf{r}})_{s}}
		{(P,(Q_{i})_{i\in\bb{\mathbf{s}}})}
		{\left(\left(\!\begin{smallmatrix}v^{s_{i}}_{i}\\Q_{i}\end{smallmatrix}\!\right)_{i\in \bb{\mathbf{s}}}\right)^{\sharp}_{s}(P)}
		$$
		where, for $\left(\!\begin{smallmatrix}v^{s_{i}}_{i}\\Q_{i}\end{smallmatrix}\!\right)_{i\in \bb{\mathbf{s}}}$, the $S$-sorted mapping from $\vs{\mathbf{s}}$
		to $\mathrm{T}_{\Sigma}(\vs{\mathbf{r}})$ canonically associated to the family
		$(Q_{i})_{i\in\bb{\mathbf{s}}}$, $\left(\left(\!\begin{smallmatrix}v^{s_{i}}_{i}\\Q_{i}\end{smallmatrix}\!\right)_{i\in \bb{\mathbf{s}}}\right)^{\sharp}$, also denoted by $\mathcal{Q}^{\sharp}$, is the
		unique homomorphism from $\mathbf{T}_{\Sigma}(\vs{\mathbf{s}})$ into
		$\mathbf{T}_{\Sigma}(\vs{\mathbf{r}})$ such that
		$$
		\left(\left(\!\begin{smallmatrix}v^{s_{i}}_{i}\\Q_{i}\end{smallmatrix}\!\right)_{i\in \bb{\mathbf{s}}}\right)^{\sharp}\circ \eta^{\vs{\mathbf{s}}} = \left(\!\begin{smallmatrix}v^{s_{i}}_{i}\\Q_{i}\end{smallmatrix}\!\right)_{i\in \bb{\mathbf{s}}}.
		$$
		Sometimes, to abbreviate, we will write $\xi_{\mathbf{r},\mathbf{s},s}$ instead of
		$\xi_{\mathbf{r},\mathbf{s},s}^{\mathbf{Tm}_{\mathrm{H}_{S}}(\Sigma)}$.
	\end{enumerate}
	
	Then $\mathbf{Tm}_{\mathrm{H}_{S}}(\Sigma)$ is a Hall algebra, the \emph{Hall algebra for} $(S,\Sigma)$.
\end{proposition}

\begin{remark}
	For every $\Sigma$-algebra $\mathbf{A}$ there exists a homomorphism of Hall algebras $\mathrm{d}^{\mathbf{A}} = (\mathrm{d}^{\mathbf{A}}_{\mathbf{\mathbf{s}},s})_{(\mathbf{\mathbf{s}},s)\in S^{\star}\times S}$ from the Hall algebra $\mathbf{Tm}_{\mathrm{H}_{S}}(\Sigma)$ to the Hall algebra $\mathbf{O}_{\mathrm{H}_{S}}(\mathbf{A})$ and its image is $\mathbf{D}_{\mathrm{H}_{S}}(\mathbf{A})$, the Hall subalgebra of $\mathbf{O}_{\mathrm{H}_{S}}(\mathbf{A})$ of the finitary derived operations on $\mathbf{A}$, which, for every $(\mathbf{\mathbf{s}},s)\in S^{\star}\times S$, is such that 
	$\mathrm{D}_{\mathrm{H}_{S}}(\mathbf{A})_{\mathbf{\mathbf{s}},s} = 
	\mathrm{D}_{\mathbf{\mathbf{s}}}(\mathbf{A})_{s}$, where $\mathrm{D}_{\mathbf{\mathbf{s}}}(\mathbf{A})_{s}$ is the $s$-th coordinate of the underlying $S$-sorted of 
	$\mathbf{D}_{\mathbf{\mathbf{s}}}(\mathbf{A})$, the $\Sigma$-algebra of the $\mathbf{\mathbf{s}}$-\emph{ary} derived operations on $\mathbf{A}$, introduced in Definition~\ref{wdop}.
\end{remark}

%\begin{proposition}
%Let $\Sigma$ be an $S$-sorted signature. Then $\mathbf{Tm}_{\mathrm{H}_{S}}(\Sigma)^{\wp}$ is a Hall algebra.
%\end{proposition}

%recall that the underlying $S^{\star}\times S$-sorted of $\mathbf{Tm}_{\mathrm{H}_{S}}(\Sigma)^{\wp}$ is $(\mathrm{Sub}(\mathrm{T}_{\Sigma}(\vs{w})_{s}))_{(w,s)\in S^{\star}\times S}$

Our next goal is to prove that, for every $S^{\star}\times S$-sorted set $\Sigma$, $\mathbf{T}_{\mathrm{H}_{S}}(\Sigma)$, the free Hall algebra on $\Sigma$, is isomorphic to $\mathbf{Tm}_{\mathrm{H}_{S}}(\Sigma)$. We remark that the existence of this isomorphism is interesting because it enables us to get a more tractable description of the terms in $\mathbf{T}_{\mathrm{H}_{S}}(\Sigma)$.
%We remark that the existence of this isomorphism is interesting because it enables us, on the one hand, to get a more tractable description of the terms in $\mathbf{T}_{\mathrm{H}_{S}}(\Sigma)$, and, on the other hand, as we will show afterwards, to state, for every $\Sigma$-algebra $\mathbf{A}$, taking into account the adjunction $\mathbf{T}_{\mathrm{H}_{S}}\dashv \mathrm{G}_{\mathrm{H}_{S}}$, the existence of a homomorphism of Hall algebras
%$\mathrm{Tr}^{\mathbf{A}}$ from $\mathbf{Tm}_{\mathrm{H}_{S}}(\Sigma)$ to $\mathbf{O}_{\mathrm{H}_{S}}(\mathbf{A}) = \mathbf{O}_{\mathrm{H}_{S}}(A)$ such that $\mathrm{Th}_{\Sigma}(\mathbf{A})$, the finitary $\Sigma$-equational theory determined by $\mathbf{A}$, is
%precisely $\mathrm{Ker}(\mathrm{Tr}^{\mathbf{A}})$, the kernel of the homomorphism $\mathrm{Tr}^{\mathbf{A}}$.

To attain the goal just stated we begin by defining, for a Hall algebra $\mathbf{A}$, an $S$-sorted signature $\Sigma$, an
$S^{\star}\times S$-mapping $f\colon\Sigma\mor A$, and a word $\mathbf{r}\in S^{\star}$, the concept of derived $\Sigma$-algebra of
$\mathbf{A}$ for $(f,\mathbf{r})$, since it will be used afterwards in the proof of the isomorphism between
$\mathbf{T}_{\mathrm{H}_{S}}(\Sigma)$ and $\mathbf{Tm}_{\mathrm{H}_{S}}(\Sigma)$.



\begin{definition}
	Let $\mathbf{A}$ be a Hall algebra and $\Sigma$ an $S$-sorted signature. Then, for every $f\colon\Sigma\mor A$ and $\mathbf{r}\in
	S^{\star}$, $\mathbf{A}^{f,\mathbf{r}}$, the \emph{derived} $\Sigma$-\emph{algebra} \emph{of} $\mathbf{A}$ \emph{for} $(f,\mathbf{r})$,
	is the $\Sigma$-algebra with underlying $S$-sorted set $A^{f,\mathbf{r}} = (A_{\mathbf{r},s})_{s\in S}$ and algebraic structure $F^{f,\mathbf{r}}$, defined, for every $(\mathbf{s},s)\in S^{\star}\times S$, as
	$$
	F^{f,\mathbf{r}}_{\mathbf{s},s} \nfunction
	{\Sigma_{\mathbf{s},s}}{\mathrm{O}_{\mathbf{s}}(A^{f,\mathbf{r}})_{s}} {\sigma}{\nfunction
		{\prod_{i\in \bb{\mathbf{s}}}A_{\mathbf{u},s_{i}}}
		{A_{\mathbf{u},s}}
		{(a_{0},\ldots,a_{\bb{\mathbf{s}}-1})}
		{\xi_{\mathbf{r},\mathbf{s},s}^{\mathbf{A}}(f_{(\mathbf{s},s)}(\sigma),a_{0},\ldots,a_{\bb{\mathbf{s}}-1})}
	}
	$$
	where $\mathrm{O}_{\mathbf{s}}(A^{f,\mathbf{r}})_{s} = A_{\mathbf{r},s}^{\prod_{i\in \bb{\mathbf{s}}}A_{\mathbf{r},s_{i}}}$.
	
	Furthermore, we will denote by $p^{\mathbf{r}}$ the $S$-sorted mapping from $\vs{\mathbf{r}}$ to $A^{f,\mathbf{r}}$ defined, for every $s\in S$ and $i\in\bb{\mathbf{r}}$, as $p^{\mathbf{r}}_{s}(v^{s}_{i}) = (\pi^{\mathbf{r}}_{i})^{\mathbf{A}}$, and by $(p^{\mathbf{r}})^{\sharp}$ the unique homomorphism from
	$\mathbf{T}_{\Sigma}(\vs{\mathbf{r}})$ to $\mathbf{A}^{f,\mathbf{r}}$ such that $(p^{\mathbf{r}})^{\sharp}\circ \eta^{\vs{\mathbf{r}}} = p^{\mathbf{r}}$.
\end{definition}

\begin{remark}
	For a $\Sigma$-algebra $\mathbf{B}=(B,G)$, we have that $G\colon\Sigma\mor\mathrm{O}_{\mathrm{H}_{S}}(B)$ and
	$\mathbf{B}\iso\mathbf{O}_{\mathrm{H}_{S}}(B)^{G,\lambda}$, where $\lambda$ is the empty word on $S$.  Besides, for every $\mathbf{r}\in S^{\star}$, we have that $\mathbf{B}^{B_{\mathbf{r}}}$, the direct $B_{\mathbf{r}}$-power of $\mathbf{B}$, is isomorphic to $\mathbf{O}_{\mathrm{H}_{S}}(B)^{G,\mathbf{r}}$.
\end{remark}



\begin{lemma}\label{L:aux}
	Let $\Sigma$ be an $S$-sorted signature, $\mathbf{A}$ a Hall algebra, $f\colon\Sigma\mor A$ and $\mathbf{q}\in S^{\star}$.  Then, for
	every $(\mathbf{r},s)\in S^{\star}\times S$, $P\in \mathrm{T}_{\Sigma}(\vs{\mathbf{r}})_{s}$ and $a\in \prod_{i\in \bb{\mathbf{r}}}A_{\mathbf{q},r_{i}}$, we have that
	$$
	P^{\mathbf{A}^{f,\mathbf{q}}}(a_{0},\ldots,a_{\bb{\mathbf{r}}-1}) =
	\xi_{\mathbf{q},\mathbf{r},s}^{\mathbf{A}}((p^{\mathbf{r}})^{\sharp}_{s}(P),a_{0},\ldots,a_{\bb{\mathbf{r}}-1}).
	$$
\end{lemma}

\begin{proof}
	By algebraic induction on the complexity of $P$.  If $P$ is a variable $v_{i}^{s}$, with $i\in\bb{\mathbf{r}}$, then
	\begin{align*}
		v_{i}^{s,\mathbf{A}^{f,\mathbf{q}}}(a_{0},\ldots,a_{\bb{\mathbf{r}}-1})\ &=
		a^{\sharp}_{r_{i}}(v^{s}_{i})\\
		&=
		a_{i} \\
		&=
		\xi_{\mathbf{q},\mathbf{r},s}^{\mathbf{A}}((\pi^{\mathbf{r}}_{i})^{\mathbf{A}},a_{0},\ldots,a_{\bb{\mathbf{r}}-1} )
		\tag{by $\mathrm{H}_{1}$} \\
		&=
		\xi_{\mathbf{q},\mathbf{r},s}^{\mathbf{A}}((p^{w})^{\sharp}_{s}(v^{s}_{i}),a_{0},\ldots,a_{\bb{\mathbf{r}}-1} ).
	\end{align*}
	Let us assume that $P = \sigma(Q_{0},\ldots,Q_{\bb{x}-1})$, with $\sigma\in \Sigma_{\mathbf{s},s}$ and that, for every $j\in\bb{x}$, $Q_{j}\in \mathrm{T}_{\Sigma}(\vs{\mathbf{r}})_{s_{j}}$ fulfils the induction hypothesis. Then we have that
	\begin{flushleft}
		$(\sigma(Q_{0},\ldots,Q_{\bb{\mathbf{s}}-1}))^{\mathbf{A}^{f,\mathbf{q}}}
		(a_{0},\ldots,a_{\bb{\mathbf{r}}-1})$
		\allowdisplaybreaks
		\begin{align*}
			&=
			\sigma^{\mathbf{A}^{f,\mathbf{q}}}
			(
			Q_{0}^{\mathbf{A}^{f,\mathbf{q}}}(a_{0},\ldots,a_{\bb{\mathbf{r}}-1}),
			\ldots,
			Q_{\bb{\mathbf{s}}-1}^{\mathbf{A}^{f,\mathbf{q}}}(a_{0},\ldots,a_{\bb{\mathbf{r}}-1})
			)\\
			&=
			\xi_{\mathbf{q},\mathbf{s},s}^{\mathbf{A}}
			(
			f(\sigma),
			Q_{0}^{\mathbf{A}^{f,\mathbf{q}}}(a_{0},\ldots,a_{\bb{\mathbf{r}}-1}),
			\ldots,
			Q_{\bb{\mathbf{s}}-1}^{\mathbf{A}^{f,\mathbf{q}}}(a_{0},\ldots,a_{\bb{\mathbf{r}}-1})
			)\\
			&=
			\xi_{\mathbf{q},\mathbf{s},s}^{\mathbf{A}}
			(
			f(\sigma),
			\xi_{\mathbf{q},\mathbf{r},s_{0}}^{\mathbf{A}}
			((p^{\mathbf{r}})^{\sharp}_{s_{0}}(Q_{0}),a_{0},\ldots,a_{\bb{\mathbf{r}}-1}),
			\ldots, 
			\\&\qquad\qquad\qquad\qquad\qquad\qquad
			\xi_{\mathbf{q},\mathbf{r},s_{\bb{\mathbf{s}}-1}}^{\mathbf{A}}
			((p^{\mathbf{r}})^{\sharp}_{s_{\bb{\mathbf{s}}-1}}(Q_{\bb{\mathbf{s}}-1}),a_{0},\ldots,a_{\bb{\mathbf{r}}-1})
			) \tag{by ind.}
			\\&=
			\xi_{\mathbf{q},\mathbf{r},s}^{\mathbf{A}}
			(
			\xi_{\mathbf{r},\mathbf{s},s}^{\mathbf{A}}
			(f(\sigma),
			(p^{\mathbf{r}})^{\sharp}_{s_{0}}(Q_{0}),
			\ldots,
			(p^{\mathbf{r}})^{\sharp}_{s_{\bb{\mathbf{s}}-1}}(Q_{\bb{\mathbf{s}}-1})
			),
			a_{0},
			\ldots,
			a_{\bb{\mathbf{r}}-1}
			) \tag{by $\mathrm{H}_{3}$}   \\
			&=
			\xi_{\mathbf{q},\mathbf{r},s}^{\mathbf{A}}
			(
			\sigma^{\mathbf{A}_{\mathbf{r}}}
			(
			(p^{\mathbf{r}})^{\sharp}_{s_{0}}(Q_{0}),
			\ldots,
			(p^{\mathbf{r}})^{\sharp}_{s_{\bb{\mathbf{s}}-1}}(Q_{\bb{\mathbf{s}}-1})
			),
			a_{0},
			\ldots,
			a_{\bb{\mathbf{r}}-1}
			) \\
			&=
			\xi_{\mathbf{q},\mathbf{r},s}^{\mathbf{A}}
			(
			(p^{\mathbf{r}})^{\sharp}_{s}(\sigma,Q_{0},\ldots,Q_{\bb{\mathbf{s}}-1}),
			a_{0},
			\ldots,
			a_{\bb{\mathbf{r}}-1}
			) \\
			&=
			\xi_{\mathbf{q},\mathbf{r},s}^{\mathbf{A}}
			(
			(p^{\mathbf{r}})^{\sharp}_{s}(P),
			a_{0},
			\ldots,
			a_{\bb{\mathbf{r}}-1}
			).
		\end{align*}
	\end{flushleft}
	This completes the proof.
\end{proof}

Next we prove that, for every $S^{\star}\times S$-sorted set $\Sigma$, the Hall algebra for $(S,\Sigma)$ is isomorphic to the free Hall algebra on $\Sigma$.


\begin{proposition}\label{iso:FrH-TerH}
	Let $\Sigma$ be an $S$-sorted signature, i.e., an $S^{\star}\times S$-sorted set.  Then the Hall algebra $\mathbf{Tm}_{\mathrm{H}_{S}}(\Sigma)$ is isomorphic to $\mathbf{T}_{\mathrm{H}_{S}}(\Sigma)$, the free Hall algebra on $\Sigma$.
\end{proposition}

\begin{proof}
	It is enough to prove that $\mathbf{Tm}_{\mathrm{H}_{S}}(\Sigma)$ has the universal property of $\mathbf{T}_{\mathrm{H}_{S}}(\Sigma)$. Therefore we have to specify an $S^{\star}\times S$-sorted mapping $\eta^{\Sigma}$ from $\Sigma$ to $\mathrm{Tm}_{\mathrm{H}_{S}}(\Sigma)$ such that, for every Hall algebra $\mathbf{A}$ and $S^{\star}\times S$-sorted mapping $f$ from $\Sigma$ to $A$, there is a unique homomorphism $\widehat{f}$ from $\mathbf{Tm}_{\mathrm{H}_{S}}(\Sigma)$ to $\mathbf{A}$ such that
	$\widehat{f}\circ \eta^{\Sigma} = f$.  Let $\eta^{\Sigma}$ be the $S^{\star}\times S$-sorted mapping defined, for every
	$(\mathbf{s},s)\in S^{\star}\times S$, as%
	$$
	\eta^{\Sigma}_{\mathbf{s},s}\nfunction
	{\Sigma_{\mathbf{s},s}}
	{\mathrm{T}_{\Sigma}(\vs{\mathbf{s}})_{s}}
	{\sigma}
	{\sigma\left(v^{s_{0}}_{0},\ldots,v^{s_{\bb{\mathbf{s}}-1}}_{\bb{\mathbf{s}}-1}\right)}
	$$
	Let $\mathbf{A}$ be a Hall algebra, $f\colon\Sigma\mor A$ an $S^{\star}\times S$-sorted mapping and $\widehat{f}$ the $S^{\star}\times S$-sorted mapping from $\mathrm{Tm}_{\mathrm{H}_{S}}(\Sigma)$ to $A$ defined, for every $(\mathbf{s},s)\in S^{\star}\times S$, as $\widehat{f}_{\mathbf{s},s} = (p^{\mathbf{s}})^{\sharp}_{s}$, where, we recall, $(p^{\mathbf{s}})^{\sharp}$ is the unique homomorphism from $\mathbf{T}_{\Sigma}(\vs{\mathbf{s}})$ to $\mathbf{A}^{f,\mathbf{s}}$ such that $(p^{\mathbf{s}})^{\sharp}\circ \eta^{\vs{\mathbf{s}}} = p^{\mathbf{s}}$. Then $\widehat{f}$ is a homomorphism of Hall algebras, because, on the one hand, for $\mathbf{s}\in S^{\star}$ and $i\in\bb{\mathbf{s}}$ we have that%
	\begin{align*}
		\widehat{f}_{\mathbf{s},s_{i}}((\pi^{\mathbf{s}}_{i})^{\mathbf{Tm}_{\mathrm{H}_{S}}(\Sigma)})
		&=
		\widehat{f}_{\mathbf{s},s_{i}}(v^{s_{i}}_{i}) \\
		&=
		p^{\mathbf{s}}_{s_{i}}(v^{s_{i}}_{i}) \\
		&=
		(\pi^{\mathbf{s}}_{i})^{\mathbf{A}},
	\end{align*}
	and, on the other hand, for $P\in \mathrm{T}_{\Sigma}(\vs{\mathbf{s}})_{s}$ and
	$\mathcal{Q} = (Q_{i})_{i\in \bb{\mathbf{s}}}\in \mathrm{T}_{\Sigma}(\vs{\mathbf{r}})_{\mathbf{s}}$ we have that %
	\begin{flushleft}
		$\widehat{f}_{\mathbf{r},s}(
		\xi_{\mathbf{r},\mathbf{s},s}^{\mathbf{Tm}_{\mathrm{H}_{S}}(\Sigma)}
		(P,Q_{0},\ldots,Q_{\bb{\mathbf{s}}-1})
		)$
		\begin{align*}
			\quad    &=
			(p^{\mathbf{r}})^{\sharp}_{s}(
			\mathcal{Q}^{\sharp}_{s}(P)
			) \\
			&=
			((p^{\mathbf{r}})^{\sharp}\circ \mathcal{Q})^{\sharp}_{s} (P)
			\tag{$(p^{\mathbf{r}})^{\sharp}\circ\mathcal{Q}^{\sharp} =
				((p^{\mathbf{r}})^{\sharp}\circ \mathcal{Q})^{\sharp}$}
			\\
			&=
			P^{\mathbf{A}^{f,\mathbf{r}}}
			((p^{\mathbf{r}})^{\sharp}_{s_{0}}(Q_{0}),\ldots,
			(p^{\mathbf{r}})^{\sharp}_{s_{\bb{\mathbf{s}}-1}}(Q_{\bb{\mathbf{s}}-1})
			) \\
			&=
			\xi_{\mathbf{r},\mathbf{s},s}^{\mathbf{A}}
			((p^{\mathbf{r}})^{\sharp}_{s}(P),
			(p^{\mathbf{r}})^{\sharp}_{s_{0}}(Q_{0}),\ldots,
			(p^{\mathbf{r}})^{\sharp}_{s_{\bb{\mathbf{s}}-1}}(Q_{\bb{\mathbf{s}}-1})
			)
			\tag{by Lemma~\ref{L:aux}}
			\\
			&=
			\xi_{\mathbf{r},\mathbf{s},s}^{\mathbf{A}}
			( \widehat{f}_{\mathbf{s},s}(P),
			\widehat{f}_{\mathbf{r},s_{0}}(Q_{0}),\ldots,
			\widehat{f}_{\mathbf{r},s_{\bb{\mathbf{s}}-1}}(Q_{\bb{\mathbf{s}}-1})
			).
		\end{align*}
	\end{flushleft}
	Therefore the $S^{\star}\times S$-sorted mapping $\widehat{f}$ is a homomorphism. Furthermore, $\widehat{f}\circ \eta^{\Sigma} = f$, because, for every $\mathbf{s}\in S^{\star}$, $s\in S$, and $\sigma\in \Sigma_{\mathbf{s},s}$, we have that%
	\begin{align*}
		\widehat{f}_{\mathbf{s},s}(\eta^{\Sigma}_{\mathbf{s},s}(\sigma)) &=
		(p^{\mathbf{s}})^{\sharp}_{s}(\sigma(v^{s_{0}}_{0},\ldots,v^{s_{\bb{\mathbf{s}}-1}}_{\bb{\mathbf{s}}-1}))
		\\&=
		\sigma^{\mathbf{A}_{\mathbf{s}}}
		(p^{\mathbf{s}}_{{s}_{0}}(v^{s_{0}}_{0}),
		\ldots,
		p^{\mathbf{s}}_{{s}_{\bb{\mathbf{s}}-1}}(v^{s_{\bb{\mathbf{s}}-1}}_{\bb{\mathbf{s}}-1})
		) \\
		&=
		\xi_{\mathbf{s},\mathbf{s},s}^{\mathbf{A}}
		( f_{(\mathbf{s},s)}(\sigma),
		(\pi^{\mathbf{s}}_{0})^{\mathbf{A}},
		\ldots,
		(\pi^{\mathbf{s}}_{\bb{\mathbf{s}}-1})^{\mathbf{A}}
		) \\
		&=
		f_{\mathbf{s},s}(\sigma) \tag{by $\mathrm{H}_{2}$}.
	\end{align*}
	It is obvious that $\widehat{f}$ is the unique homomorphism such that $\widehat{f}\circ \eta^{\Sigma} = f$.  Henceforth $\mathbf{Tm}_{\mathrm{H}_{S}}(\Sigma)$ is isomorphic to $\mathbf{T}_{\mathrm{H}_{S}}(\Sigma)$.
\end{proof}


This isomorphism, together with the adjunction $\mathbf{T}_{\mathrm{H}_{S}}\dashv\mathrm{G}_{\mathrm{H}_{S}}$,  has as an immediate consequence that, for every $S$-sorted set $A$ and every $S$-sorted signature $\Sigma$, the sets $\mathrm{Hom}({\Sigma},\mathrm{G}_{\mathrm{H}_{S}}(\mathbf{O}_{\mathrm{H}_{S}}(A)))$, in 
$\mathsf{Set}^{S^{\star}\times S}$, and $\mathrm{Hom}(\mathbf{Tm}_{\mathrm{H}_{S}}(\Sigma),\mathbf{O}_{\mathrm{H}_{S}}(A))$, in $\mathsf{Alg}(\mathrm{H}_{S})$, are naturally isomorphic.

Actually, for an $S$-sorted set $A$ and an $S$-sorted signature $\Sigma$,
\begin{enumerate}
	\item the mapping that assigns to a structure of $\Sigma$-algebra $F$ on $A$ (i.e., an $S^{\star}\times S$-sorted mapping $F$ from $\Sigma$ to $\mathrm{O}_{\mathrm{H}_{S}}(A)$) the unique  homomorphism of Hall algebras $\widehat{F} = (\widehat{F}_{w,s})_{(w,s)\in S^{\star}\times S}$ from $\mathbf{Tm}_{\mathrm{H}_{S}}(\Sigma)$ to $\mathbf{O}_{\mathrm{H}_{S}}(A)$ such that $\widehat{F}\circ \eta^{\Sigma} = F$, together with
	\item the mapping that assigns to a homomorphism $h$ from $\mathbf{Tm}_{\mathrm{H}_{S}}(\Sigma)$ to $\mathbf{O}_{\mathrm{H}_{S}}(A)$, essentially, the algebraic structure $\mathrm{G}_{\mathrm{H}_{S}}(h)\circ \eta^{\Sigma}$ on $A$, where $\eta^{\Sigma}$ is the canonical embedding of $\Sigma$ into $\mathrm{T}_{\mathrm{H}_{S}}(\Sigma)$, 
\end{enumerate}
which are natural in both variables and mutually inverse bijections, provide the mentioned natural isomorphism.

\begin{remark}
	For an $S$-sorted set $A$, a structure of $\Sigma$-algebra $F$ on $A$ and a word $\mathbf{s}\in S^{\star}$, the subfamily $(\widehat{F}_{\mathbf{s},s})_{s\in S}$ of $\widehat{F} = (\widehat{F}_{\mathbf{s},s})_{(\mathbf{s},s)\in S^{\star}\times S}$,
	denoted by $\mathrm{d}^{(A,F)} = (\mathrm{d}^{(A,F)}_{\mathbf{s},s})_{s\in S}$, is the unique homomorphism from $\mathbf{T}_{\Sigma}(\vs{\mathbf{s}})$ to $(A,F)^{A_{\mathbf{s}}}$, the direct $A_{\mathbf{s}}$-power of $(A,F)$, such that $\mathrm{d}^{(A,F)}_{\mathbf{s}}\circ \eta^{\vs{\mathbf{s}}} = \mathrm{p}^{A}_{\mathbf{s}}$, where $\mathrm{p}^{A}_{\mathbf{s}}$ is the $S$-sorted mapping from $\vs{\mathbf{s}}$ to $A^{A_{\mathbf{s}}}$ defined, for every $s\in S$ and every $v^{s}_{i}\in(\vs{\mathbf{s}})_{s}$, as $\mathrm{p}^{A}_{\mathbf{s},s}(v^{s}_{i}) = \mathrm{pr}^{A}_{\mathbf{s},i}$.
\end{remark}

For a set of sorts $S$ and an $S$-sorted signature $\Sigma$, the derived operation symbols of $\Sigma$, i.e., the elements of the components $\mathrm{T}_{\Sigma}(\vs{\mathbf{s}})_{s}$ of $\mathrm{Tm}_{\mathrm{H}_{S}}(\Sigma)$ can be considered as the operation symbols of a new signature. The mappings $\varphi$ from a set of sorts $S$ to another $T$, together with the $S^{\star}\times S$-sorted mappings $d$ that assign to operation symbols of an $S$-sorted signature $\Sigma$ derived operations symbols of a $T$-sorted signature $\Lambda$, i.e., terms of the components $\mathrm{T}_{\Lambda}(\vs{\varphi^{\star}(\mathbf{s})})_{\varphi(s)}$ of $\mathrm{Tm}_{\mathrm{H}_{T}}(\Lambda)_{\varphi^{\star}\times\varphi}$, form a new class of morphisms between many-sorted signatures which we will call \emph{derivors}.

\begin{definition}\label{DDerivor}
	Let $\mathbf{\Sigma} = (S,\Sigma)$ and $\mathbf{\Lambda} = (T,\Lambda)$ be many-sorted signatures. A \emph{derivor} from $\mathbf{\Sigma}$ to $\mathbf{\Lambda}$ is an ordered pair $\mathbf{d} = (\varphi,d)$, where $\varphi\colon S\mor T$ and  $d\colon\Sigma\mor\mathrm{Tm}_{\mathrm{H}_{T}}(\Lambda)_{\varphi^{\star}\times\varphi}$. Thus, for every $(\mathbf{s},s)\in S^{\star}\times S$, %
	$$
	d_{\mathbf{s},s}\colon \Sigma_{\mathbf{s},s}\mor\mathrm{Tm}_{\mathrm{H}_{T}}(\Lambda)_{\varphi^{\star}(\mathbf{s}),\varphi(s)}=
	\mathrm{T}_{\Lambda}(\vs{\varphi^{\star}(\mathbf{s})})_{\varphi(s)}.
	$$
	Hence, for every $\sigma\in \Sigma_{\mathbf{s},s}$, $d_{\mathbf{s},s}(\sigma)\in \mathrm{T}_{\Lambda}(\vs{\varphi^{\star}(\mathbf{s})})_{\varphi(s)}$, i.e., $d_{\mathbf{s},s}(\sigma)$ is a term of sort $\varphi(s)$ with variables in $\vs{\varphi^{\star}(\mathbf{s})}$.
\end{definition}



\begin{remark}\label{Rdop}
	Let $\varphi$ be a mapping from the set of sorts $S$ to the set of sorts $T$ and let $V^{T}$ be a fixed $T$-sorted set with a countably infinity of variables in each component, say $V^{T}_{t} = \{v^{t}_{n}\mid n\in \mathbb{N}\}$ (which is isomorphic to $\{t\}\times\mathbb{N}$). Then, for every $\mathbf{s}\in S^{\star}$, we let $\vs{\varphi^{\star}(\mathbf{s})} = ((\vs{\varphi^{\star}(\mathbf{s})})_{t})_{t\in T}$ stand for the $T$-sorted set  defined, for every $t\in T$, as:
	$$
	(\vs{\varphi^{\star}(\mathbf{s})})_{t} = \{v^{t}_{0},\ldots,v^{t}_{\bb{\varphi^{\star}(\mathbf{s})}_{t}-1}\} (\cong \varphi^{\star}(\mathbf{s})^{-1}[\{t\}] = \{i\in \bb{\varphi^{\star}(\mathbf{s})}\mid \varphi(w_{i}) = t\}).
	$$
	Thus, for every $t\in T$, $(\vs{\varphi^{\star}(\mathbf{s})})_{t}$ is an initial segment of $V^{T}_{t}$ and 
	$\vs{\varphi^{\star}(\mathbf{s})}$ is a finite subset of $V^{T}$. The latter is fulfilled because $\mathrm{supp}_{T}(\vs{\varphi^{\star}(\mathbf{s})}) = \mathrm{Im}(\varphi^{\star}(\mathbf{s}))$ and $\bb{\varphi^{\star}(\mathbf{s})} = \bb{\mathbf{s}}$ is finite. Moreover, $\mathrm{supp}_{T}(\vs{\varphi^{\star}(\mathbf{s})}) =\varphi[\mathrm{supp}_{S}(\vs{\mathbf{s}})]$. Recall  that $\varphi^{\star}(\mathbf{s}) = (\varphi(s_{i}))_{i\in \bb{\mathbf{s}}}$ and that, since $\varphi^{\star}(\lambda) = \lambda$, $\vs{\varphi^{\star}(\lambda)} = \varnothing^{T} = (\varnothing)_{t\in T}$.
	
	If $\mathbf{d} = (\varphi,d)$ is a derivor from $\mathbf{\Sigma}$ to $\mathbf{\Lambda}$, $\mathbf{B}$ a 
	$\mathbf{\Lambda}$-algebra, $(\mathbf{s},s)\in S^{\star}\times S$ and $\sigma\in \Sigma_{\mathbf{s},s}$, then the derived operation (or term operation or polynomial operation) of type $(\varphi^{\star}(\mathbf{s}),\varphi(s))$ determined by the term $d_{\mathbf{s},s}(\sigma)$ in $\mathbf{B}$, denoted by $d_{\mathbf{s},s}(\sigma)^{\mathbf{B}}$ or by $d(\sigma)^{\mathbf{B}}$ for short, is a mapping from $B_{\varphi^{\star}(\mathbf{s})}$ to $B_{\varphi(s)}$. In particular, for a $T$-sorted set $Y$, we have that the derived operation (or term operation or polynomial operation) of type $(\varphi^{\star}(\mathbf{s}),\varphi(s))$ determined by the term $d_{\mathbf{s},s}(\sigma)$ in $\mathbf{T}_{\Lambda}(Y)$, denoted by $d_{\mathbf{s},s}(\sigma)^{\mathbf{T}_{\Lambda}(Y)}$ or by  $d(\sigma)^{\mathbf{T}_{\Lambda}(Y)}$ for short, is a mapping from $\mathrm{T}_{\Lambda}(Y)_{\varphi^{\star}(\mathbf{s})}$ to $\mathrm{T}_{\Lambda}(Y)_{\varphi(s)}$. Notice that, for every $s\in S$ and every $\sigma\in \Sigma_{\lambda,s}$, $d_{\lambda,s}(\sigma)\in \mathrm{T}_{\Lambda}(\varnothing^{T})_{\varphi(s)}$, i.e., $d_{\lambda,s}(\sigma)$ is a closed term of sort $\varphi(s)$.  Therefore, for every $\mathbf{\Lambda}$-algebra $\mathbf{B}$, $d_{\lambda,s}(\sigma)^{\mathbf{B}}$ will be a mapping from $B_{\lambda}$ to $B_{\varphi(s)}$, i.e., essentially, a distinguished element of $B_{\varphi(s)}$. In particular, for a $T$-sorted set $Y$, $d_{\lambda,s}(\sigma)^{\mathbf{T}_{\Lambda}(Y)}$ will be a mapping from $\mathrm{T}_{\Lambda}(Y)_{\lambda}$ to $\mathrm{T}_{\Lambda}(Y)_{\varphi(s)}$, i.e., essentially, a distinguished element of $\mathrm{T}_{\Lambda}(Y)_{\varphi(s)}$. Let us note that $d(\sigma)^{\mathbf{B}}$ is the value of $\mathrm{d}^{\mathbf{B}}_{\varphi^{\star}(\mathbf{s}),s}$ at $d_{\mathbf{s},s}(\sigma)$ and that $d(\sigma)^{\mathbf{T}_{\Lambda}(Y)}$ is the value of $\mathrm{d}^{\mathbf{T}_{\Lambda}(Y)}_{\varphi^{\star}(\mathbf{s}),s}$ at $d_{\mathbf{s},s}(\sigma)$, where 
	$\mathrm{d}^{\mathbf{B}}_{\varphi^{\star}(\mathbf{s})}$ is the unique homomorphism from $\mathbf{T}_{\Lambda}(\vs{\varphi^{\star}(\mathbf{s})})$ to $\mathbf{B}^{B_{\varphi^{\star}(\mathbf{s})}}$ for which the left diagram below commutes, and $\mathrm{d}^{\mathbf{T}_{\Lambda}(Y)}_{\varphi^{\star}(\mathbf{s})}$ the unique homomorphism from $\mathbf{T}_{\Lambda}(\vs{\varphi^{\star}(\mathbf{s})})$ to $\mathbf{T}_{\Lambda}(Y)^{\mathrm{T}_{\Lambda}(Y)_{\varphi^{\star}(\mathbf{s})}}$ for which the right diagram below commutes.
	$$
	\begin{aligned}
		\xymatrix@C=40pt@R=33pt{
			\vs{\varphi^{\star}(\mathbf{s})}
			\ar[r]^-{\eta^{\vs{\varphi^{\star}(\mathbf{s})}}}
			\ar[rd]_-{\mathrm{p}^{B}_{\varphi^{\star}(\mathbf{s})}} &
			\mathrm{T}_{\Lambda}(\vs{\varphi^{\star}(\mathbf{s})}) \ar[d]^{\mathrm{d}^{\mathbf{B}}_{\varphi^{\star}(\mathbf{s})}}
			\\
			& B^{B_{\varphi^{\star}(\mathbf{s})}}
		}
		%
	\end{aligned}
	\quad \text{} \qquad\;
	\begin{aligned}
		%
		\xymatrix@C=40pt@R=33pt{
			\vs{\varphi^{\star}(\mathbf{s})}
			\ar[r]^-{\eta^{\vs{\varphi^{\star}(\mathbf{s})}}}
			\ar[rd]_-{\mathrm{p}^{\mathrm{T}_{\Lambda}(Y)}_{\varphi^{\star}(\mathbf{s})}} &
			\mathrm{T}_{\Lambda}(\vs{\varphi^{\star}(\mathbf{s})}) \ar[d]^{\mathrm{d}^{\mathbf{T}_{\Lambda}(Y)}_{\varphi^{\star}(\mathbf{s})}}
			\\
			& \mathrm{T}_{\Lambda}(Y)^{\mathrm{T}_{\Lambda}(Y)_{\varphi^{\star}(\mathbf{s})}}
		}
		%
	\end{aligned}
	$$
	Recall that, for every $t\in T$ and every $i\in (\vs{\varphi^{\star}(\mathbf{s})})_{t}$ (thus $(\varphi^{\star}(\mathbf{s}))(i) = \varphi(s_{i}) = t$), $\mathrm{p}^{B}_{\varphi^{\star}(\mathbf{s}),t}(i)$ is $\mathrm{pr}^{B}_{\varphi^{\star}(\mathbf{s}),i}$, i.e., the $\varphi^{\star}(\mathbf{s})$-ary, $i$-th canonical projection from $B_{\varphi^{\star}(\mathbf{s})}$ to $B_{t} = B_{(\varphi^{\star}(\mathbf{s}))(i)}$.
\end{remark}

\begin{remark}
	While derivors are not the most general type of morphism that might be considered between many-sorted signatures---for instance, one could consider polyderivors, see~\cite{CVST10b}---, they are an important class of such morphisms. One reason for its relevance is its formal properties (see below), another that there are many mathematical examples of them which are of interest (see~\cite{CVST10b}).
\end{remark}

Since by, Proposition~\ref{iso:FrH-TerH}, $\mathbf{Tm}_{\mathrm{H}_{T}}(\Lambda)$ is isomorphic to 
$\mathbf{T}_{\mathrm{H}_{T}}(\Lambda)$, the derivors can be defined, alternative, but equivalently, as ordered pairs  $\mathbf{d} = (\varphi,d)$ with $d\colon\Sigma\mor \mathrm{T}_{\mathrm{H}_{T}}(\Lambda)$.

On the other hand, every mapping $\varphi\colon S\mor T$ determines a functor from the category  $\mathsf{Alg}(\mathrm{H}_{T})$ to the category $\mathsf{Alg}(\mathrm{H}_{S})$, so $\mathrm{Tm}_{\mathrm{H}_{T}}(\Lambda)_{\varphi^{\star}\times\varphi}$ is in its turn e\-quipped with a structure of Hall algebra for $S$, which will allow us, in particular, to define the composition of derivors. We next show the existence of such a functor by defining a morphism of algebraic presentations from  $(\Sigma^{\mathrm{H}_{S}},\mathcal{E}^{\mathrm{H}_{S}})$ to $(\Sigma^{\mathrm{H}_{T}},\mathcal{E}^{\mathrm{H}_{T}})$.



\begin{proposition}
	Let $\varphi\colon S\mor T$ be a mapping between sets of sorts and $h^{\varphi}$ the $S^{\star}\times S$-sorted mapping from $\Sigma^{\mathrm{H}_{S}}$ to $\Sigma^{\mathrm{H}_{T}}_{\varphi^{\star}\times \varphi}$ defined as follows: 
	\begin{enumerate}
		\item For every $\mathbf{s}\in S^{\star}$ and every $i\in\bb{\mathbf{s}}$, $h^{\varphi}(\pi^{\mathbf{s}}_{i}) = \pi^{\varphi^{\star}(\mathbf{s})}_{i}$.
		
		\item  For every $\mathbf{r},\,\mathbf{s}\in S^{\star}$ and every $s\in S$,
		$h^{\varphi}(\xi_{\mathbf{r},\mathbf{s},s}) = \xi_{\varphi^{\star}(\mathbf{r}),\varphi^{\star}(\mathbf{s}),\varphi(s)}$.
	\end{enumerate}
	Then $(\varphi^{\star}\times \varphi, h^{\varphi}) \colon(S^{\star}\times S,\Sigma^{\mathrm{H}_{S}},\mathcal{E}^{\mathrm{H}_{S}})\mor
	(T^{\star}\times T,\Sigma^{\mathrm{H}_{T}},\mathcal{E}^{\mathrm{H}_{T}})$ is a morphism of algebraic presentations.
\end{proposition}

The morphisms of algebraic presentations determine functors in the opposite direction between the associated categories of algebras. Therefore, each mapping $\varphi\colon S\mor T$ between sets of sorts, determines a functor $(\varphi^{\star}\times \varphi, h^{\varphi})^{\ast}$ from $\mathsf{Alg}(\mathrm{H}_{T})$ to $\mathsf{Alg}(\mathrm{H}_{S})$, which transforms Hall algebras for $T$ into Hall algebras for $S$. The action of the functor on the free Hall algebra on a $T$-sorted signature $\Lambda$ is a Hall algebra for $S$, whose underlying $S^{\star}\times S$-sorted set is
$\mathrm{Tm}_{\mathrm{H}_{T}}(\Lambda)_{\varphi^{\star}\times \varphi}$.

If $\mathbf{d}\colon\mathbf{\Sigma}\mor\mathbf{\Lambda}$ is a derivor, then
$d\colon \Sigma\mor \mathrm{Tm}_{\mathrm{H}_{T}}(\Lambda)_{\varphi^{\star}\times\varphi}$ determines a homomorphism of Hall algebras $d^{\sharp}\colon\mathbf{Tm}_{\mathrm{H}_{S}}(\Sigma)\mor
\mathbf{Tm}_{\mathrm{H}_{T}}(\Lambda)_{\varphi^{\star}\times\varphi}$. Thus, for every $(\mathbf{s},s)\in S^{\star}\times S$, $d^{\sharp}_{\mathbf{s},s}$ sends terms in $\mathrm{T}_{\Sigma}(\vs{\mathbf{s}})_{s}$ to terms in $\mathrm{T}_{\Lambda}(\vs{\varphi^{\star}(\mathbf{s})})_{\varphi(s)}$.

\begin{definition}\label{DDerCompId}
	Let $\mathbf{d}\colon\mathbf{\Sigma}\mor\mathbf{\Lambda}$ and $\mathbf{e}\colon\mathbf{\Lambda}\mor\mathbf{\Omega}$ be  derivors. Then $\mathbf{e}\circ \mathbf{d} = (\psi,e)\circ(\varphi,d)$, the \emph{composition} of $\mathbf{d}$ and  $\mathbf{e}$, is the derivor $(\psi\circ\varphi,e^{\sharp}_{\varphi^{\star}\times
		\varphi}\circ d)$, where $e^{\sharp}_{\varphi^{\star}\times\varphi}\circ d$ is obtained from
	$$
	\begin{aligned}
		\xymatrix@C=40pt@R=30pt{
			\Lambda
			\ar[r]^-{\eta^{\Lambda}}
			\ar[rd]_-{e} &
			\mathrm{Tm}_{\mathrm{H}_{T}}(\Lambda) \ar[d]^{e^{\sharp}}
			\\
			& \mathrm{Tm}_{\mathrm{H}_{U}}(\Omega)_{\psi^{\star}\times \psi}
		}
		%
	\end{aligned}
	\qquad \text{as} \qquad\;\;
	\begin{aligned}
		%
		\xymatrix@C=40pt@R=30pt{
			\mathrm{Tm}_{\mathrm{H}_{T}}(\Lambda)_{\varphi^{\star}\times \varphi}
			\ar[d]^{e^{\sharp}_{\varphi^{\star}\times \varphi}} &
			\Sigma \ar[l]_-{d} \\
			{\mathrm{Tm}_{\mathrm{H}_{U}}(\Omega)_{\psi^{\star}\times\psi}}_{\varphi^{\star}\times\varphi}
		}
	\end{aligned}
	$$
	being $e^{\sharp}$ the canonical extension of $e$ to the free Hall algebra on $\Lambda$.
	
	For every many-sorted signature $\mathbf{\Sigma} = (S,\Sigma)$, the \emph{identity} at $(S,\Sigma)$
	is $(\id^{S},\eta^{\Sigma})$.
\end{definition}

The just stated definition allows us to form a category of many-sorted signatures whose morphisms are the derivors.

\begin{proposition}\label{PSigdCat}
	The many-sorted signatures together with the derivors constitute a category, denoted by $\mathsf{Sig}_{\mathfrak{d}}$.
\end{proposition}



\begin{remark}
	Let $\mathsf{Sig}$ be the category whose objects are the many-sorted signatures and whose morphisms from $\mathbf{\Sigma}$ to $\mathbf{\Lambda}$ are the ordered pairs $(\varphi,d)$, where $\varphi$ is a mapping from $S$ to $T$ and $d$ an $S^{\star}\times S$-sorted mapping from $\Sigma$ to $\Lambda_{\varphi^{\star}\times \varphi}$ (thus a mapping in $\mathsf{Sig}(S)$). Note that $\mathsf{Sig}$ is the category obtained by means of the Grothendieck construction for an  appropriate contravariant functor from $\mathsf{Set}$ to $\mathsf{Cat}$.
	We next show that the category $\mathsf{Sig}_{\mathfrak{d}}$ can be obtained as the Kleisli category for a suitable monad. In fact, for every  set of sorts $S$, we have the adjoint situation $\mathbf{T}_{\mathrm{H}_{S}}\dashv\mathrm{G}_{\mathrm{H}_{S}}$
	%$$
	%\xymatrix@=5pc{ \mathsf{Alg}(\mathrm{H}_{S})
		%\ar@<1.5ex>[r]^{\mathrm{G}_{\mathrm{H}_{S}}} \ar@{}[r]|{\uadj} &
		%\mathbf{Set}^{S^{\star}\times S}
		%\ar@<1.5ex>[l]^{\mathbf{T}_{\mathrm{H}_{S}}} }
	%$$
	and, thus, a monad on $\mathsf{Sig}(S)$ which we will denote by $\mathbb{T}_{\mathrm{H}_{S}} = (\mathrm{T}_{\mathrm{H}_{S}},\eta^{\mathrm{H}_{S}},\mu^{\mathrm{H}_{S}})$. Then the ordered triple $\boldsymbol{\mathfrak{d}} = (\mathfrak{d},\eta,\mu)$ in which (1) $\mathfrak{d}$ is the endofunctor at $\mathsf{Sig}$ that sends $(S,\Sigma)$ to $(S,\mathrm{T}_{\mathrm{H}_{S}}(\Sigma))$ and a morphism $(\varphi,d)$ from $(S,\Sigma)$ to  $(T,\Lambda)$ to the morphism $(\varphi,d^{\sharp})$ from $(S,\mathrm{T}_{\mathrm{H}_{S}}(\Sigma))$ to  $(T,\mathrm{T}_{\mathrm{H}_{T}}(\Lambda))$, (2) $\eta$ the natural transformation from $\mathrm{Id}_{\mathsf{Sig}}$ to  $\mathfrak{d}$ that sends $(S,\Sigma)$ to $\eta_{(S,\Sigma)} = (\mathrm{id},\eta^{\mathrm{H}_{S}}_{\Sigma})$, and (3) $\mu$ the natural transformation from $\mathfrak{d}\circ \mathfrak{d}$ to $\mathfrak{d}$ that sends $(S,\Sigma)$ to $\mu_{(S,\Sigma)} = (\mathrm{id},\mu^{\mathrm{H}_{S}}_{\Sigma})$, is a monad on $\mathsf{Sig}$ and $\mathsf{Kl}(\boldsymbol{\mathfrak{d}})$, the Kleisli category for $\boldsymbol{\mathfrak{d}}$, is isomorphic to $\mathsf{Sig}_{\mathfrak{d}}$. This shows, in our opinion, the mathematical naturalness of the notion of derivor. Moreover, by defining a suitable notion of transformation between derivors one can equip the category $\mathsf{Sig}_{\mathfrak{d}}$ with a structure of $2$-category (this was, in fact, already done in~\cite{CVST10b} for polyderivors).
	
\end{remark}



\begin{remark}
	Since $\mathsf{Sig}$ has coproducts, $\mathsf{Kl}(\boldsymbol{\mathfrak{d}})\cong \mathsf{Sig}_{\mathfrak{d}}$ has coproducts. In fact, let $(\mathbf{\Sigma}^{i})_{i\in I}$ be an $I$-indexed family in $\mathrm{Ob}(\mathsf{Kl}(\boldsymbol{\mathfrak{d}}))$ and $(\coprod_{i\in I}\mathbf{\Sigma}^{i},(\mathrm{in}^{\mathbf{\Sigma}^{i}})_{i\in I})$ the coproduct of $(\mathbf{\Sigma}^{i})_{i\in I}$ in $\mathsf{Sig}$, where  the first coordinate of $\coprod_{i\in I}\mathbf{\Sigma}^{i}$ is $\coprod_{i\in I}S_{i}$ and the second coordinate of $\coprod_{i\in I}\mathbf{\Sigma}^{i}$, denoted by $\coprod_{i\in I}\mathrm{\Sigma}^{i}$, is defined, for every $(((s_{i_{j}},i_{j})_{j\in n}),(s,i))\in (\coprod_{i\in I}S_{i})^{\star}\times \coprod_{i\in I}S_{i}$, as follows: 
	\begin{equation*}
		\textstyle
		(\coprod_{i\in I}\mathrm{\Sigma}^{i})_{(((s_{i_{j}},i_{j})_{j\in n}),(s,i))} = 
		\begin{cases}
			\Sigma^{i}_{\mathbf{s},s}, \text{\,\,if $\exists\,\mathbf{s}\in S_{i}^{\star}$ such that $(s_{i_{j}},i_{j})_{j\in n} = \mathrm{in}_{i}^{\star}(\mathbf{s})$;}\\
			\varnothing, \text{\,\,otherwise}.
		\end{cases}
	\end{equation*}
	Let us recall that $\mathrm{in}_{i}^{\star}$ is the monoid homomorphism from the free monoid on $S_{i}$ to the free monoid on $\coprod_{i\in I}S_{i}$ that sends $\mathbf{s} = (s_{j})_{j\in n}\in S_{i}^{\star}$ to 
	$(s_{j},i)_{j\in n}\in (\coprod_{i\in I}S_{i})^{\star}$. Therefore, 
	$(s_{i_{j}},i_{j})_{j\in n} = \mathrm{in}_{i}^{\star}(\mathbf{s})$ if and only if, for every $j\in n$, $i_{j} = i$ and $s_{i_{j}} = w_{j}$.
	Then $(\boldsymbol{\mathfrak{d}}(\coprod_{i\in I}\mathbf{\Sigma}^{i}),(\eta_{\coprod_{i\in I}\mathbf{\Sigma}^{i}}\circ \mathrm{in}^{\mathbf{\Sigma}^{i}})_{i\in I})$ will be the required coproduct of $(\mathbf{\Sigma}^{i})_{i\in I}$ in $\mathsf{Kl}(\boldsymbol{\mathfrak{d}})$. This property of the category $\mathsf{Sig}_{\mathfrak{d}}$ will be essential in our work because to be able to compare Curry-Howard towers it will be necessary, as an intermediate step, to construct, by means of the formation of coproducts, from given many-sorted  signatures $(\mathbf{\Sigma}^{i})_{i\in I}$, a new many-sorted signature $\coprod_{i\in I}\mathbf{\Sigma}^{i}$ (finally, from $\coprod_{i\in I}\mathbf{\Sigma}^{i}$, we will construct the many-sorted signature $\biguplus_{i\in I}\mathbf{\Sigma}^{i}$ which will be the one of interest for us).
	
	%$\mathbf{\Sigma}$, $\mathbf{\Omega}_{\mathrm{c}}$, $\mathbf{\Omega}_{\mathrm{c},\mathcal{A}}$, $\mathbf{\Lambda}$, and $\mathbf{\Omega}_{\mathrm{c},\mathcal{B}}$ and $\mathsf{Sig}_{\mathfrak{d}}$-morphisms 
	%$\mathbf{d}$ from $\mathbf{\Sigma}$ to $\mathbf{\Lambda}$ and $\mathbf{e}$ from $\mathbf{\Omega}_{\mathrm{c},\mathcal{A}}$ to $\mathbf{\Omega}_{\mathrm{c},\mathcal{B}}$, as, e.g.,  
	%$$
	%\textstyle
	%\mathbf{d}\amalg\mathbf{id}_{\mathbf{\Omega}_{\mathrm{c}}}\colon \mathbf{\Sigma}\amalg\mathbf{\Omega}_{\mathrm{c}}\mor \mathbf{\Lambda}\amalg\mathbf{\Omega}_{\mathrm{c}} 
	%\quad\text{and}\quad
	%\mathbf{d}\amalg\mathbf{e}\colon \mathbf{\Sigma}\amalg\mathbf{\Omega}_{\mathrm{c},\mathcal{A}}\mor \mathbf{\Lambda}\amalg\mathbf{\Omega}_{\mathrm{c},\mathcal{B}}.
	%$$
\end{remark}

\section{Generalized derivors}
\textcolor{red}{En el nostre cas no volem els sorts ordenats, no treballem en fòrmules, treballem en E-equacions. Tampoc necessitem predicats. La $d_{\mathbf{w}}$ de l'apartat 4, ha de correspondre amb la $d$ de la Definició~\ref{DDerivor}}

\begin{definition}\label{DGeneralDerivor}
	\sloppy
	Let $\mathbf{\Omega} = (S,\leq,\Sigma,\Pi)$ and $\mathbf{\Omega}' = (S',\leq',\Sigma',\Pi')$ be order-sorted signatures with predicates, and let $X$ and $X'$ be the corresponding $S$-sorted and $S'$-sorted sets of variables. Recall that a mapping $\tau\colon S\mor S'$ is \emph{monotone} if, for every $s,s'\in S$, $s\leq s'$ implies $\tau(s)\leq'\tau(s')$. A \emph{general derivor} from $\mathbf{\Omega}$ to $\mathbf{\Omega}'$ (see~\cite{SR18}) consists of
	\begin{enumerate}
		\item a monotone mapping $\tau\colon S\mor S'$,
		\item a mapping $d\colon S\mor\mathrm{Form}_{\Sigma',\Pi'}$,
		\item a family of mappings $d_{\mathbf{w},s}\colon\Sigma_{\mathbf{w},s}\mor\mathrm{Form}_{\Sigma',\Pi'}$, and
		\item a family of mappings $d_{\mathbf{w}}\colon\Pi_{\mathbf{w}}\mor\mathrm{Form}_{\Sigma',\Pi'}$,
	\end{enumerate}
	such that
	\begin{enumerate}
		\item $d(s)$ is a formula $\Delta_{s}(x)$ with at most a single free variable
		$x\in X_{s}\cap X'_{\tau(s)}$;
		\item for every $f\in\Sigma_{\mathbf{w},s}$, with $\mathbf{w} = s_{1}\cdots s_{k}$,
		$d_{\mathbf{w},s}(f)$ is a formula $\Phi^{f}_{\mathbf{w},s}(x_{1},\ldots,x_{k},y)$ in
		$\mathrm{Form}_{\Sigma',\Pi'}$ with free variables $x_{i}\in X_{s_{i}}\cap X'_{\tau(s_{i})}$,
		$1\leq i\leq k$, and $y\in X_{s}\cap X'_{\tau(s)}$;
		\item for every $P\in\Pi_{\mathbf{w}}$, with $\mathbf{w} = s_{1}\cdots s_{n}$,
		$d_{\mathbf{w}}(P)$ is a formula $\Phi^{P}_{\mathbf{w}}(x_{1},\ldots,x_{n})$ in
		$\mathrm{Form}_{\Sigma',\Pi'}$ with free variables $x_{i}\in X_{s_{i}}\cap X'_{\tau(s_{i})}$,
		$1\leq i\leq n$.
	\end{enumerate}
	When no confusion arises, we often write $s'$ instead of $\tau(s)$. Moreover, an ordinary derivor can be seen as a general derivor by taking, for every $f\in\Sigma_{\mathbf{w},s}$, the equation $y = d_{\mathbf{w},s}(f)$, with $y\in X'_{s'}$ a fresh variable, instead of the term $d_{\mathbf{w},s}(f)\in\mathrm{T}_{\Sigma'}(X')_{\tau(s)}$.
\end{definition}

\section{ep-derivors}

An equationally partial signature is an ordred triple $(S,\Sigma,\mathrm{def})$ in which $(S,\Sigma)$ is a sorted signature and, for every $(\mathbf{w},s)\in S^{\star}\times S$ and every $\sigma\in \Sigma_{\mathbf{w},s}$, $\mathrm{def}(\sigma)\subseteq \coprod_{s\in S}\mathrm{T}_{\Sigma}(\downarrow\mathbf{w})^{2}_{s}$ or, what is equivalent, a subset of $(\mathrm{T}_{\Sigma}(\downarrow\mathbf{w})^{2}_{s})_{s\in S}$. We call $\mathrm{def}(\sigma)$ the set of domain conditions associated to $\sigma$.

Note that $\mathrm{Sub}(\coprod_{s\in S}\mathrm{T}_{\Sigma}(\downarrow\mathbf{w})^{2}_{s})\cong 
\mathrm{Sub}((\mathrm{T}_{\Sigma}(\downarrow\mathbf{w})^{2}_{s})_{s\in S})$.

Let $\mathbf{A}$ be a partial $\Sigma$-algebra and $\sigma\in \Sigma_{\mathbf{w},s}$. We denote by 
$A_{\mathrm{def}(\sigma)}$ the subset of $A_{\mathbf{w}}$ defined as:
$$
A_{\mathrm{def}(\sigma)} = \{\mathbf{a}\in A_{\mathbf{w}}\mid \forall\, 
(P,Q)\in \mathrm{def}(\sigma)\, (\mathbf{A}\models_{\mathbf{a}} (P,Q))\}.
$$
Moreover, we will say that $\mathbf{A}$ is an equationally partial $\Sigma$-algebra of type $(S,\Sigma,\mathrm{def})$ if, for every $(\mathbf{w},s)\in S^{\star}\times S$ and every $\sigma\in \Sigma_{\mathbf{w},s}$, $\mathrm{Dom}(F^{\mathbf{A}}_{\sigma}) = A_{\mathrm{def}(\sigma)}$.

To define the notion of morphism from an equationally partial signature $(S,\Sigma,\mathrm{def})$ to another $(S',\Sigma',\mathrm{def}')$ we extend $\mathrm{def}$ to the terms by algebraic recursion (See Reichel). We denote by the same symbol such an extension or by $\mathrm{deft}$.

\begin{definition}\label{DDerivorEP}Let $\mathbf{\Sigma} = (S,\Sigma)$ and $\mathbf{\Lambda} = (T,\Lambda)$ be many-sorted signatures. A \emph{derivor} from $\mathbf{\Sigma}$ to $\mathbf{\Lambda}$ is an ordered pair $\mathbf{d} = (\varphi,d)$, where $\varphi\colon S\mor T$ and  $d\colon\Sigma\mor\mathrm{Tm}_{\mathrm{H}_{T}}(\Lambda)_{\varphi^{\star}\times\varphi}$. Thus, for every $(\mathbf{s},s)\in S^{\star}\times S$, %
	$$
	d_{\mathbf{s},s}\colon \Sigma_{\mathbf{s},s}\mor\mathrm{Tm}_{\mathrm{H}_{T}}(\Lambda)_{\varphi^{\star}(\mathbf{s}),\varphi(s)}=
	\mathrm{T}_{\Lambda}(\vs{\varphi^{\star}(\mathbf{s})})_{\varphi(s)}.
	$$
	Hence, for every $\sigma\in \Sigma_{\mathbf{s},s}$, $d_{\mathbf{s},s}(\sigma)\in \mathrm{T}_{\Lambda}(\vs{\varphi^{\star}(\mathbf{s})})_{\varphi(s)}$, i.e., $d_{\mathbf{s},s}(\sigma)$ is a term of sort $\varphi(s)$ with variables in $\vs{\varphi^{\star}(\mathbf{s})}$.
\end{definition}

Let $(S,\Sigma,\mathrm{def})$ and $(S',\Sigma',\mathrm{def}')$ be  equationally partial signature (ep-signatures). A morphism from $(S,\Sigma,\mathrm{def})$ to $(S',\Sigma',\mathrm{def}')$ is a derivor $\mathbf{d} = (\varphi,d)$ from $(S,\Sigma)$ to $(S',\Sigma')$ such that, for every 
$(\mathbf{w},s)\in S^{\star}\times S$ and every $\sigma\in \Sigma_{\mathbf{w},s}$, 
$d_{\mathbf{s},s}^{\sharp}\times d_{\mathbf{s},s}^{\sharp}[\mathrm{def}(\sigma)]$ is included in 
$\mathrm{Cn}(\mathrm{def}(d_{\mathbf{w},s}^{\sharp}(\sigma)))$.

I am not sure that this is the correct definition. For it to be the case, the morphism between ep-signatures should induce a functor from the category of algebras for $(S',\Sigma',\mathrm{def}')$ to the category of algebras for $(S,\Sigma,\mathrm{def})$. This may be the reason why Reichel requires the semantic consequences of $d_{\mathbf{s},s}^{\sharp}\times d_{\mathbf{s},s}^{\sharp}[\mathrm{def}(\sigma)]$ and $\mathrm{def}(d_{\mathbf{w},s}^{\sharp}(\sigma))$ to coincide.

\section{generalized ep-derivors}

\section{Compact presentation of higher-order categorial algebras}\label{SCompact}

We begin by recalling the definitions of compact category, compact $2$-category, compact $n$-category, and many-sorted categorial algebra. These are the compact (single-sorted) presentations of the corresponding categorical structures.

\begin{definition}\label{DCompactCat}
A \emph{compact category} is an ordered pair $\mathsf{C} = (C,\xi)$ where $C$ is a set and $\xi = (\#,\mathrm{sc},\mathrm{tg})$ a structure of category on $C$, i.e., $\#\colon C\times C\dmor C$ is a partial operation of \emph{composition}, and $\mathrm{sc}$, $\mathrm{tg}$ are unary operations on $C$ which assign, to every morphism in $C$, its \emph{source} and its \emph{target}, respectively, subject to satisfy the following conditions:
\begin{enumerate}
\item[(C1)]\label{C1} For every $f\in C$, 
\allowdisplaybreaks
\begin{align*}
\mathrm{sc}(\mathrm{sc}(f)) &= \mathrm{sc}(f),&
 \mathrm{sc}(\mathrm{tg}(f)) &= \mathrm{tg}(f),\\
\mathrm{tg}(\mathrm{sc}(f)) &= \mathrm{sc}(f),&
 \mathrm{tg}(\mathrm{tg}(f)) &= \mathrm{tg}(f).
\end{align*}
\item[(C2)]\label{C2} For every $f,g\in C$, 
\allowdisplaybreaks
\begin{align*}
g\#f &\mbox{ is defined}&
 \mbox{if and only if}&&
  \mathrm{sc}(g)& = \mathrm{tg}(f).
\end{align*} 
\item[(C3)]\label{C3} For every $f,g\in C$, if $\mathrm{sc}(g) = \mathrm{tg}(f)$, then 
\allowdisplaybreaks
\begin{align*}
\mathrm{sc}(g\# f) &= \mathrm{sc}(f),
&
\mathrm{tg}(g\# f) &= \mathrm{tg}(g).
\end{align*}
\item[(C4)]\label{C4} For every $f\in C$, 
\allowdisplaybreaks
\begin{align*}
f\# \mathrm{sc}(f) &= f,
&
\mathrm{tg}(f)\# f &= f.
\end{align*}
\item[(C5)]\label{C5} For every $f,g,h\in C$, if $\mathrm{sc}(h)=\mathrm{tg}(g)$ and $\mathrm{sc}(g)=\mathrm{tg}(f)$, then 
\allowdisplaybreaks
\begin{align*}
h\# (g\# f) = (h\# g)\# f.
\end{align*}
\end{enumerate}

We will call the $f\in C$ such that $\mathrm{sc}(f) = f$ or, what is equivalent, $\mathrm{tg}(f) = f$,  \emph{identities} of $C$. 

Let $\mathsf{C}=(C,\xi)$ and $\mathsf{C}'=(C',\xi')$ be two compact categories. A \emph{compact functor} from $\mathsf{C}$ to $\mathsf{C}'$ is an ordered triple $(\mathsf{C},F,\mathsf{C}')$, denoted by
$F\colon \mathsf{C}\mor\mathsf{C}'$, where $F$ is a mapping from $C$ to $C'$ such that
\begin{itemize}
\item[(Ci)]\label{Ci} If $f\in C$, then
\allowdisplaybreaks
\begin{align*}
F(\mathrm{sc}(f))&=\mathrm{sc}'(F(f)),
&
F(\mathrm{tg}(f))&=\mathrm{tg}'(F(f)).
\end{align*}
\item[(Cii)]\label{Cii} If $f,g\in C$ are such that $\mathrm{sc}(g)=\mathrm{tg}(f)$, then
\allowdisplaybreaks
\begin{align*}
F(g\#f)=F(g)\#'F(f).
\end{align*}
\end{itemize}

We will say that a compact category $\mathsf{C}$ is \emph{discrete} if every $f\in C$ is an identity (i.e., if it contains no non-identity morphisms).

Compact functors can be composed. In fact, if $F$ is a compact functor from $\mathsf{C}$ to $\mathsf{D}$ and $G$ is a compact functor from $\mathsf{D}$ to $\mathsf{E}$, then $G\circ F$ is a compact functor from $\mathsf{C}$ to $\mathsf{E}$, which we call the \emph{composite} of $F$ and $G$. This composition is associative. Moreover, for every compact category $\mathsf{C}$, there is an \emph{identity} compact functor $(\mathsf{C},\mathrm{Id}_{C},\mathsf{C})$, denoted by $\mathrm{Id}_{\mathsf{C}}\colon \mathsf{C}\mor\mathsf{C}$, where $\mathrm{Id}_{C}$ is the identity mapping  at $C$, which is the neutral element for the composition. We denote by $\mathsf{Cat}^{}$ the category consisting of compact categories and compact functors.
\end{definition}

\begin{definition}\label{DCompact2Cat}
A \emph{compact $2$-category} is an ordered pair $\mathsf{C} = (C,(\xi_{k})_{k\in 2})$
where $C$ is a set and, for every $k\in 2$, $\xi_{k}= (\#^{k},\mathrm{sc}^{k},\mathrm{tg}^{k})$ a structure of compact category on $C$, such that the following conditions are satisfied:
\begin{enumerate}
\item[(2C1)]\label{2C1} For every $f\in C$,
\allowdisplaybreaks
\begin{align*}
\mathrm{sc}^{1}(\mathrm{sc}^{0}(f))&=\mathrm{sc}^{0}(f)
=\mathrm{sc}^{0}(\mathrm{sc}^{1}(f))=\mathrm{sc}^{0}(\mathrm{tg}^{1}(f)),\\
\mathrm{tg}^{1}(\mathrm{tg}^{0}(f))&=\mathrm{tg}^{0}(f)
=\mathrm{tg}^{0}(\mathrm{tg}^{1}(f))=\mathrm{tg}^{0}(\mathrm{sc}^{1}(f)).
\end{align*}
\item[(2C2)]\label{2C2} For every $f,g\in C$, if $\mathrm{sc}^{0}(g)=\mathrm{tg}^{0}(f)$, then
\allowdisplaybreaks
\begin{align*}
\mathrm{sc}^{1}(g{\#^{0}}f)&=\mathrm{sc}^{1}(g){\#^{0}}\mathrm{sc}^{1}(f),
&
\mathrm{tg}^{1}(g{\#^{0}}f)&=\mathrm{tg}^{1}(g){\#^{0}}\mathrm{tg}^{1}(f).
\end{align*}
\item[(2C3)]\label{2C3} For every $f,f',g,g'\in C$, if 
\allowdisplaybreaks
\begin{align*}
\mathrm{sc}^{1}(g')&=\mathrm{tg}^{1}(g);
&
\mathrm{sc}^{1}(f')&=\mathrm{tg}^{1}(f);\\
\mathrm{sc}^{0}(g')&=\mathrm{tg}^{0}(f');
&
\mathrm{sc}^{0}(g)&=\mathrm{tg}^{0}(f),
\end{align*}
then
\allowdisplaybreaks
\begin{align*}
(g'\#^{0}f')\#^{1}(g\#^{0} f)
=
(g'\#^{1}g)\#^{0}(f'\#^{1}f)
.
\end{align*}
\end{enumerate}

We identify the family $(\xi_{k})_{k\in 2}$ with the ordered pair $(\xi_{0},\xi_{1})$.

Let $\mathsf{C}=(C,(\xi_{k})_{k\in 2})$ and $\mathsf{C}'=(C',(\xi'_{k})_{k\in 2})$ be two compact $2$-categories. A \emph{compact $2$-functor} from $\mathsf{C}$ to $\mathsf{C}'$ is an ordered triple $(\mathsf{C},F,\mathsf{C}')$, denoted by $F\colon \mathsf{C}\mor\mathsf{C}'$, where $F$ is a mapping from $C$ to $C'$ such that, for every $k\in 2$, $F$ is a compact functor from $(C,\xi_{k})$ to $(C',\xi'_{k})$. 

Compact $2$-functors can be composed. In fact, if $F$ is a compact $2$-functor from $\mathsf{C}$ to $\mathsf{D}$ and $G$ is a compact $2$-functor from $\mathsf{D}$ to $\mathsf{E}$, then $G\circ F$ is a compact $2$-functor from $\mathsf{C}$ to $\mathsf{E}$, which we call the \emph{composite} of $F$ and $G$. This composition is associative. Moreover, for every compact $2$-category $\mathsf{C}$, there is an \emph{identity} compact $2$-functor $(\mathsf{C},\mathrm{Id}_{C},\mathsf{C})$, denoted by $\mathrm{Id}_{\mathsf{C}}\colon \mathsf{C}\mor\mathsf{C}$, where $\mathrm{Id}_{C}$ is the identity mapping  at $C$, which is the neutral element for the composition. We denote by $\mathsf{2Cat}^{}$ the category consisting of compact 
$2$-categories and compact $2$-functors.
\end{definition}

\begin{definition}\label{DCompactnCat} 
A \emph{compact $n$-category}, for $n\in\mathbb{N}$, is an ordered pair $\mathsf{C} = (C,(\xi_{k})_{k\in n})$ where $C$ is a set and, for every $k\in n$, $\xi_{k}$ a structure of  compact category on $C$ such that, for every $j,k\in n$ with $j<k$, the ordered pair $(C,(\xi_{i})_{i\in\{j,k\}})$ is a compact $2$-category. For $j<k$, we identify $(\xi_{i})_{i\in \{j,k\}}$ with $(\xi_{j},\xi_{k})$. Thus a compact $0$-category is just a set, a compact $1$-category is a compact category in the sense of Definition~\ref{DCompactCat}, and a compact $2$-category is one in the sense of Definition~\ref{DCompact2Cat}.

For \(k\in n\) we will say that \(f\in C\) is a \emph{\(k\)-cell} if \(f=\mathrm{sc}^{k}(f)\) or, equivalently, by item~(C1) in Definition~\ref{DCompactCat}, if \(f=\mathrm{tg}^{k}(f)\). Then we denote by \(C_{k}\) the set of \(k\)-cells of \(\mathsf{C}\), i.e.,
\[
C_{k}=\{f\in C\mid f=\mathrm{sc}^{k}(f)\}.
\]
Moreover, we let \(C_{n}\) stand for \(C\) and, by analogy to the above, the elements of \(C_{n}\) will be called \emph{\(n\)-cells}.

A \emph{compact $n$-functor} $F$ from $\mathsf{C}$ to $\mathsf{C}'$, denoted by
$F\colon \mathsf{C}\mor\mathsf{C}'$, is a mapping from $C$ to $C'$ such that, for every $k\in n$, $F$ is a compact functor from $(C, \xi_{k})$ to $(C', \xi'_{k})$. 

Compact $n$-functors can be composed. In fact, if $F$ is a compact $n$-functor from $\mathsf{C}$ to $\mathsf{D}$ and $G$ is a compact $n$-functor from $\mathsf{D}$ to $\mathsf{E}$, then $G\circ F$ is a compact $n$-functor from $\mathsf{C}$ to $\mathsf{E}$, which we call the \emph{composite} of $F$ and $G$. This composition is associative. Moreover, for every compact $n$-category $\mathsf{C}$, there is an \emph{identity} compact $n$-functor $(\mathsf{C},\mathrm{Id}_{C},\mathsf{C})$, denoted by $\mathrm{Id}_{\mathsf{C}}\colon \mathsf{C}\mor\mathsf{C}$, where $\mathrm{Id}_{C}$ is the identity mapping  at $C$, which is the neutral element for the composition. We denote by $\mathsf{nCat}^{}$ the category consisting of compact $n$-categories and compact $n$-functors.
\end{definition}

\begin{definition}\label{DManySortedCatAlg}
Let $S$ be a set of sorts and $n\in\mathbb{N}$.
We will call the objects of the functor category $\mathsf{nCat}^{S}$ \emph{$S$-sorted compact $n$-categories}. Thus, an $S$-sorted compact $n$-category is an $S$-indexed family $\mathsf{A}=(\mathsf{A}_{s})_{s\in S}$ such that, for every sort $s\in S$, $\mathsf{A}_{s}$ is a compact $n$-category. Given two $S$-sorted compact $n$-categories 
$\mathsf{A}$ and $\mathsf{B}$, we will call the morphisms in $\mathsf{nCat}^{S}$ from $\mathsf{A}$ to 
$\mathsf{B}$ \emph{$S$-sorted compact $n$-functors}. Thus, an $S$-sorted compact $n$-functor from $\mathsf{A}$ to $\mathsf{B}$ is an $S$-indexed family $F=(F_{s})_{s\in S}$ such that, for every sort $s\in S$, $F_{s}$ is a compact $n$-functor from $\mathsf{A}_{s}$ to $\mathsf{B}_{s}$. We will denote by $\mathrm{nFunc}(\mathsf{A},\mathsf{B})$ the set of all $S$-sorted compact $n$-functors from $\mathsf{A}$ to $\mathsf{B}$. We consider the expressions $F\in\mathrm{nFunc}(\mathsf{A},\mathsf{B})$ and $F\colon\mathsf{A}\mor\mathsf{B}$ as synonyms.

Let $\mathsf{A}$ be an $S$-sorted compact $n$-category. The $S^{\star}\times S$-sorted set of finitary operations on $\mathsf{A}$ is $(\mathrm{nFunc}(\mathsf{A}_{\mathbf{s}},\mathsf{A}_{s}))_{(\mathbf{s},s)\in S^{\star}\times S}$, where, for every $\mathbf{s}\in S^{\star}$, $\mathsf{A}_{\mathbf{s}}=\prod_{j\in\bb{\mathbf{s}}}\mathsf{A}_{s_{j}}$, with $\bb{\mathbf{s}}$ denoting the length of the word $\mathbf{s}$ (if $\mathbf{s}=\lambda$, then $\mathsf{A}_{\lambda}$ is a final compact $n$-category). 

Let $\mathbf{\Sigma}=(S,\Sigma)$ be a many-sorted signature. 
 A \emph{structure of many-sorted $n$-categorial $\Sigma$-algebra on} an $S$-sorted $n$-category $\mathsf{A}$ is a family $F = (F_{\mathbf{s},s})_{(\mathbf{s},s)\in S^{\star}\times S}$, where, for $(\mathbf{s},s)\in S^{\star}\times S$, $F_{\mathbf{s},s}$ is a mapping from $\Sigma_{\mathbf{s},s}$ to $\mathrm{nFunc}(\mathsf{A}_{\mathbf{s}},\mathsf{A}_{s})$ (if $(\mathbf{s},s) = (\lambda,s)$ and $\sigma\in \Sigma_{\lambda,s}$, then $F_{\lambda,s}(\sigma)$ picks out a cell of $\mathsf{A}_{s}$ which is a $k$-cell for every $k\in n$, since $\mathsf{A}_{\lambda}$ is a final compact $n$-category). The finite products $\mathsf{A}_{\mathbf{s}}$ are taken componentwise: sources and targets are computed componentwise, and $(\mathsf{g}_{j})_{j}\#^{k}(\mathsf{f}_{j})_{j}$ is defined if and only if every $\mathsf{g}_{j}\#^{k}\mathsf{f}_{j}$ is defined, and then it is $(\mathsf{g}_{j}\#^{k}\mathsf{f}_{j})_{j}$.  An \emph{$S$-sorted compact $n$-categorial $\Sigma$-algebra} or, simply, a \emph{many-sorted $n$-categorial $\mathbf{\Sigma}$-algebra}, is a pair $(\mathsf{A},F)$, abbreviated to $\boldsymbol{\mathsf{A}}$, where $\mathsf{A}=(\mathsf{A}_{s})_{s\in S}$ is an $S$-sorted compact $n$-category and $F$ a structure of many-sorted $n$-categorial $\Sigma$-algebra on $\mathsf{A}$.  For a pair $(\mathbf{s},s)\in S^{\star}\times S$ and a formal operation $\sigma\in \Sigma_{\mathbf{s},s}$, in order to simplify the notation, the $n$-functor $F_{\mathbf{s},s}(\sigma)$ from $\mathsf{A}_{\mathbf{s}}$ to $\mathsf{A}_{s}$ will be written as $F_{\sigma}$, or simply as $\sigma^{\boldsymbol{\mathsf{A}}}$. 

An \emph{$S$-sorted compact $n$-categorial $\Sigma$-homomorphism} or, simply, a \emph{many-sorted $n$-categorial $\mathbf{\Sigma}$-ho\-mo\-mor\-phism}, from $\boldsymbol{\mathsf{A}}$ to $\boldsymbol{\mathsf{B}}$, where $\boldsymbol{\mathsf{B}}=(\mathsf{B},G)$, is a triple $(\boldsymbol{\mathsf{A}},F,\boldsymbol{\mathsf{B}})$, abbreviated to $F\colon\boldsymbol{\mathsf{A}}\mor\boldsymbol{\mathsf{B}}$, where $F$ is an $S$-sorted compact $n$-functor from $\mathsf{A}$ to $\mathsf{B}$ such that, for every $(\mathbf{s},s)\in S^{\star}\times S$, every $\sigma\in\Sigma_{\mathbf{s},s}$ and every family of cells $(\mathsf{a}_{j})_{j\in\bb{\mathbf{s}}}\in\mathsf{A}_{\mathbf{s}}$ (i.e., of $n$-cells, since $C_{n}=C$), we have that
$$
F_{s}\left(\sigma^{\boldsymbol{\mathsf{A}}}
\left(\left(\mathsf{a}_{j}
\right)_{j\in\bb{\mathbf{s}}}\right)\right)
=
\sigma^{\boldsymbol{\mathsf{B}}}
\left(\left(
F_{s_{j}}\left(
\mathsf{a}_{j}
\right)
\right)_{j\in\bb{\mathbf{s}}}
\right).
$$
We denote by $\mathsf{nCat}\mathsf{Alg}(\mathbf{\Sigma})$ the category of $n$-categorial  $\mathbf{\Sigma}$-algebras and $n$-categorial $\mathbf{\Sigma}$-homomorphisms.
\end{definition}

In the compact presentation, an $n$-categorial $\mathbf{\Sigma}$-algebra is regarded as a partial algebra on a single $S$-sorted set of cells, for a signature that enlarges $\Sigma$ with sources, targets and compositions of every dimension below $n$.

\begin{definition}\label{DSigmaCatCompact}
Let $\Sigma$ be an $S$-sorted signature and $n\in\mathbb{N}$. The \emph{compact $n$-categorial signature determined by $\Sigma$}, denoted by $\Sigma^{(n)}$, is the $S$-sorted signature defined, for every $(\mathbf{s},s)\in S^{\star}\times S$, as follows:
$$
{\Sigma}^{(n)}_{\mathbf{s},s}=
\begin{cases}
{\Sigma}_{\mathbf{s},s},
& \text{if $\mathbf{s}\not\in \{(s),(s,s)\}$;}\\
{\Sigma}_{\mathbf{s},s}\amalg\{\mathrm{sc}^{k}_{s}, \mathrm{tg}^{k}_{s}\mid k\in n\},
& \text{if $\mathbf{s}= (s)$;}\\
{\Sigma}_{\mathbf{s},s}\amalg\{\circ^{k}_{s}\mid k\in n\},
& \text{if $\mathbf{s}= (s,s)$.}
\end{cases}
$$
That is, ${\Sigma}^{(n)}$ is the expansion of $\Sigma$ obtained by adding, for every $k\in n$ and every sort $s\in S$, two unary operation symbols $\mathrm{sc}^{k}_{s}$ and $\mathrm{tg}^{k}_{s}$, the \emph{$k$-source} and the \emph{$k$-target}, of arity $(s)$ and coarity $s$, and a binary operation symbol $\circ^{k}_{s}$, the \emph{$k$-composition}, of arity $(s,s)$ and coarity $s$.
\end{definition}

Note that $\Sigma^{(0)}=\Sigma$ and that, for every $n,m\in\mathbb{N}$ with $n\leq m$, $\Sigma^{(n)}\subseteq\Sigma^{(m)}$. Moreover, for a rewriting system $\boldsymbol{\mathcal{A}}$, the categorial signature $\Sigma^{\boldsymbol{\mathcal{A}}}$ of Chapter~2 is $\Sigma^{(1)}$ enlarged with a constant for every rewrite rule, and, similarly, $\Sigma^{\boldsymbol{\mathcal{A}}^{(2)}}$ is $\Sigma^{(2)}$ enlarged with constants for the rewrite rules of both orders.

We next present the $n$-categorial $\mathbf{\Sigma}$-algebras as a $\mathrm{QE}$-variety of partial $\Sigma^{(n)}$-algebras. The specification follows the pattern of the specification $\boldsymbol{\mathcal{E}}^{\boldsymbol{\mathcal{A}}^{(2)}}$ of Chapter~3: the axioms (A1)--(A8), about the $0$-dimensional structure, and (B1)--(B8), about the $1$-dimensional structure, are now stated uniformly for every dimension $k\in n$, and the axioms (AB1)--(AB3), which relate dimensions $0$ and $1$, for every pair of dimensions $l<k$ in $n$. Since no rewrite rules are involved, there are no counterparts of (A9) and (B9). Informally, these equations state that the operations of $\Sigma$ are always defined, that, for every sort $s\in S$, sources, targets and compositions of sort $s$ of every dimension behave as in an $n$-category, and that the operations of $\Sigma$ are compatible with them.

\begin{definition}\label{DVarCompact}
Let $\mathbf{\Sigma}=(S,\Sigma)$ be a many-sorted signature and $n\in\mathbb{N}$. We will denote by $(S,\Sigma^{(n)},V^{S},\mathcal{C}^{\mathbf{\Sigma}^{(n)}})$, written $\boldsymbol{\mathcal{C}}^{\mathbf{\Sigma}^{(n)}}$ for short, the specification in which $\Sigma^{(n)}$ is the compact $n$-categorial signature introduced in Definition~\ref{DSigmaCatCompact}, $V^{S}$ a fixed $S$-sorted set with a countable infinity of variables in each coordinate, and $\mathcal{C}^{\mathbf{\Sigma}^{(n)}}$ the subset of $\mathrm{QE}(\Sigma^{(n)})_{V^{S}}$, the set of all $\mathrm{QE}$-equations with variables in $V^{S}$, consisting of the following equations:

For every $(\mathbf{s},s)\in S^{\star}\times S$, every operation symbol $\sigma\in \Sigma_{\mathbf{s},s}$, and every family of variables $(x_{j})_{j\in \lvert\mathbf{s}\rvert}\in V^{S}_{\mathbf{s}}$, the operation $\sigma$ applied to the family $(x_{j})_{j\in \lvert\mathbf{s}\rvert}$ is always defined. Formally,
\allowdisplaybreaks
\begin{align*}
\sigma((x_{j})_{j\in \lvert\mathbf{s}\rvert})\overset{\mathrm{e}}{=}\sigma((x_{j})_{j\in \lvert\mathbf{s}\rvert}).\label{DVarCompactA0}\tag{A0}
\end{align*}

For every $k\in n$, the following equations describe the $k$-dimensional structure.

For every sort $s\in S$ and every variable $x\in V^{S}_{s}$, the $k$-source and $k$-target of $x$ are always defined. Formally,
\allowdisplaybreaks
\begin{align*}
\mathrm{sc}^{k}_{s}(x)&\overset{\mathrm{e}}{=}\mathrm{sc}^{k}_{s}(x);
&
\mathrm{tg}^{k}_{s}(x)&\overset{\mathrm{e}}{=}\mathrm{tg}^{k}_{s}(x).\label{DVarCompactAk1}\tag{A$^{k}$1}
\end{align*}

For every sort $s\in S$ and every variable $x\in V^{S}_{s}$, we have the following equations:
\allowdisplaybreaks
\begin{align*}
\mathrm{sc}^{k}_{s}(\mathrm{sc}^{k}_{s}(x))&\overset{\mathrm{e}}{=}\mathrm{sc}^{k}_{s}(x);
&
\mathrm{sc}^{k}_{s}(\mathrm{tg}^{k}_{s}(x))&\overset{\mathrm{e}}{=}\mathrm{tg}^{k}_{s}(x);\\
\mathrm{tg}^{k}_{s}(\mathrm{sc}^{k}_{s}(x))&\overset{\mathrm{e}}{=}\mathrm{sc}^{k}_{s}(x);
&
\mathrm{tg}^{k}_{s}(\mathrm{tg}^{k}_{s}(x))&\overset{\mathrm{e}}{=}\mathrm{tg}^{k}_{s}(x).\label{DVarCompactAk2}\tag{A$^{k}$2}
\end{align*}
In other words, $\mathrm{sc}^{k}_{s}$ and $\mathrm{tg}^{k}_{s}$ are right zeros. In particular, $\mathrm{sc}^{k}_{s}$ and $\mathrm{tg}^{k}_{s}$ are idempotent.

For every sort $s\in S$ and every pair of variables $x,y\in V^{S}_{s}$, $x\circ^{k}_{s} y$ is defined if and only if the $k$-target of $y$ is equal to the $k$-source of $x$. Formally,
\allowdisplaybreaks
\begin{align*}
x\circ^{k}_{s} y\overset{\mathrm{e}}{=}x\circ^{k}_{s} y \;\;
&\longrightarrow \;\;
\mathrm{sc}^{k}_{s}(x)\overset{\mathrm{e}}{=}\mathrm{tg}^{k}_{s}(y);\\
\mathrm{sc}^{k}_{s}(x)\overset{\mathrm{e}}{=}\mathrm{tg}^{k}_{s}(y) \;\;
&\longrightarrow \;\;
x\circ^{k}_{s} y\overset{\mathrm{e}}{=}x\circ^{k}_{s}y.\label{DVarCompactAk3}\tag{A$^{k}$3}
\end{align*}

For every sort $s\in S$ and every pair of variables $x,y\in V^{S}_{s}$, if $x\circ^{k}_{s} y$ is defined, then the $k$-source of $x\circ^{k}_{s}y$ is that of $y$ and the $k$-target of $x\circ^{k}_{s}y$ is that of $x$. Formally,
\allowdisplaybreaks
\begin{align*}
x\circ^{k}_{s}y\overset{\mathrm{e}}{=}x\circ^{k}_{s}y \;\;
&\longrightarrow \;\;
\mathrm{sc}^{k}_{s}(x\circ^{k}_{s}y)\overset{\mathrm{e}}{=}\mathrm{sc}^{k}_{s}(y);\\
x\circ^{k}_{s}y\overset{\mathrm{e}}{=}x\circ^{k}_{s}y \;\;
&\longrightarrow \;\;
\mathrm{tg}^{k}_{s}(x\circ^{k}_{s}y)\overset{\mathrm{e}}{=}\mathrm{tg}^{k}_{s}(x).\label{DVarCompactAk4}\tag{A$^{k}$4}
\end{align*}

For every sort $s\in S$ and every variable $x\in V^{S}_{s}$, the compositions $x\circ^{k}_{s}\mathrm{sc}^{k}_{s}(x)$ and $\mathrm{tg}^{k}_{s}(x)\circ^{k}_{s}x$ are always defined and are equal to $x$, i.e., $\mathrm{sc}^{k}_{s}(x)$ is a right unit element for the $k$-composition with $x$ and $\mathrm{tg}^{k}_{s}(x)$ is a left unit element for the $k$-composition with $x$. Formally,
\allowdisplaybreaks
\begin{align*}
x\circ^{k}_{s}\mathrm{sc}^{k}_{s}(x)&\overset{\mathrm{e}}{=} x;
&
\mathrm{tg}^{k}_{s}(x)\circ^{k}_{s}x&\overset{\mathrm{e}}{=} x.\label{DVarCompactAk5}\tag{A$^{k}$5}
\end{align*}

For every sort $s\in S$ and every triple of variables $x, y, z\in V^{S}_{s}$, if the $k$-compositions $x\circ^{k}_{s}y$ and $y\circ^{k}_{s}z$ are defined, then the $k$-compositions $x\circ^{k}_{s}(y\circ^{k}_{s}z)$ and $(x\circ^{k}_{s}y)\circ^{k}_{s}z$ are defined and they are equal, i.e., the $k$-composition, when defined, is associative. Formally,
\allowdisplaybreaks
\begin{align*}
(x\circ^{k}_{s}y\overset{\mathrm{e}}{=} x\circ^{k}_{s}y)
\wedge
(y\circ^{k}_{s}z\overset{\mathrm{e}}{=} y\circ^{k}_{s}z)\;\;
&\longrightarrow \;\;
(x\circ^{k}_{s}y)\circ^{k}_{s}z\overset{\mathrm{e}}{=} x\circ^{k}_{s}(y\circ^{k}_{s}z).\label{DVarCompactAk6}\tag{A$^{k}$6}
\end{align*}

For every $(\mathbf{s},s)\in S^{\star}\times S$, every operation symbol $\sigma\in \Sigma_{\mathbf{s},s}$, and every family of variables $(x_{j})_{j\in \lvert\mathbf{s}\rvert}\in V^{S}_{\mathbf{s}}$, the $k$-source of $\sigma((x_{j})_{j\in \lvert\mathbf{s}\rvert})$ is equal to $\sigma$ applied to the family $((\mathrm{sc}^{k}_{s_{j}}(x_{j}))_{j\in \lvert\mathbf{s}\rvert})$, and the $k$-target of $\sigma((x_{j})_{j\in \lvert\mathbf{s}\rvert})$ is equal to $\sigma$ applied to the family $((\mathrm{tg}^{k}_{s_{j}}(x_{j}))_{j\in \lvert\mathbf{s}\rvert})$. Formally,
\allowdisplaybreaks
\begin{align*}
\mathrm{sc}^{k}_{s}(\sigma((x_{j})_{j\in \lvert\mathbf{s}\rvert}))&\overset{\mathrm{e}}{=}\sigma((\mathrm{sc}^{k}_{s_{j}}(x_{j}))_{j\in \lvert\mathbf{s}\rvert});\\
\mathrm{tg}^{k}_{s}(\sigma((x_{j})_{j\in \lvert\mathbf{s}\rvert}))&\overset{\mathrm{e}}{=}\sigma((\mathrm{tg}^{k}_{s_{j}}(x_{j}))_{j\in \lvert\mathbf{s}\rvert}).\label{DVarCompactAk7}\tag{A$^{k}$7}
\end{align*}

For every $(\mathbf{s},s)\in S^{\star}\times S$, every operation symbol $\sigma\in \Sigma_{\mathbf{s},s}$, and every pair of families of variables $(x_{j})_{j\in \lvert\mathbf{s}\rvert}, (y_{j})_{j\in \lvert\mathbf{s}\rvert}\in V^{S}_{\mathbf{s}}$, if, for every $j\in \lvert\mathbf{s}\rvert$, the $k$-compositions $x_{j}\circ^{k}_{s_{j}}y_{j}$ are defined, then the $k$-composition $\sigma((x_{j})_{j\in \lvert\mathbf{s}\rvert})\circ^{k}_{s}\sigma((y_{j})_{j\in \lvert\mathbf{s}\rvert})$ is defined and it is equal to $\sigma$ applied to the family $(x_{j}\circ^{k}_{s_{j}}y_{j})_{j\in \lvert\mathbf{s}\rvert}$. Formally,
\allowdisplaybreaks
\begin{multline*}
\textstyle
\bigwedge_{j\in \lvert\mathbf{s}\rvert}
(x_{j}\circ^{k}_{s_{j}} y_{j}\overset{\mathrm{e}}{=}x_{j}\circ^{k}_{s_{j}} y_{j})
\;\;\longrightarrow\;\;\\
\sigma((x_{j}\circ^{k}_{s_{j}}y_{j})_{j\in \lvert\mathbf{s}\rvert})
\overset{\mathrm{e}}{=}
\sigma((x_{j})_{j\in \lvert\mathbf{s}\rvert})\circ^{k}_{s}\sigma((y_{j})_{j\in \lvert\mathbf{s}\rvert}).
\label{DVarCompactAk8}\tag{A$^{k}$8}
\end{multline*}

For every $l,k\in n$ with $l<k$, the following equations relate the $l$-dimensional and the $k$-dimensional structures.

For every sort $s\in S$ and every variable $x\in V^{S}_{s}$, the elements $\mathrm{sc}^{k}_{s}(\mathrm{sc}^{l}_{s}(x))$, $\mathrm{sc}^{l}_{s}(\mathrm{sc}^{k}_{s}(x))$ and $\mathrm{sc}^{l}_{s}(\mathrm{tg}^{k}_{s}(x))$ are always defined and are equal to $\mathrm{sc}^{l}_{s}(x)$. Similarly, the elements $\mathrm{tg}^{k}_{s}(\mathrm{tg}^{l}_{s}(x))$, $\mathrm{tg}^{l}_{s}(\mathrm{tg}^{k}_{s}(x))$, and $\mathrm{tg}^{l}_{s}(\mathrm{sc}^{k}_{s}(x))$ are always defined and are equal to $\mathrm{tg}^{l}_{s}(x)$. Formally,
\allowdisplaybreaks
\begin{align*}
\mathrm{sc}^{k}_{s}(\mathrm{sc}^{l}_{s}(x))&\overset{\mathrm{e}}{=}\mathrm{sc}^{l}_{s}(x);
&
\mathrm{tg}^{k}_{s}(\mathrm{tg}^{l}_{s}(x))&\overset{\mathrm{e}}{=}\mathrm{tg}^{l}_{s}(x);\\
\mathrm{sc}^{l}_{s}(\mathrm{sc}^{k}_{s}(x))&\overset{\mathrm{e}}{=}\mathrm{sc}^{l}_{s}(x);
&
\mathrm{tg}^{l}_{s}(\mathrm{tg}^{k}_{s}(x))&\overset{\mathrm{e}}{=}\mathrm{tg}^{l}_{s}(x);\\
\mathrm{sc}^{l}_{s}(\mathrm{tg}^{k}_{s}(x))&\overset{\mathrm{e}}{=}\mathrm{sc}^{l}_{s}(x);
&
\mathrm{tg}^{l}_{s}(\mathrm{sc}^{k}_{s}(x))&\overset{\mathrm{e}}{=}\mathrm{tg}^{l}_{s}(x).\label{DVarCompactABlk1}\tag{AB$^{l,k}$1}
\end{align*}

For every sort $s\in S$ and every pair of variables $x,y\in V^{S}_{s}$, if $x\circ^{l}_{s}y$ is defined, then the $k$-source of $x\circ^{l}_{s}y$ is the $l$-composition of the $k$-source of $x$ with the $k$-source of $y$ and the $k$-target of $x\circ^{l}_{s}y$ is the $l$-composition of the $k$-target of $x$ with the $k$-target of $y$. Formally,
\allowdisplaybreaks
\begin{align*}
x\circ^{l}_{s}y&\overset{\mathrm{e}}{=}x\circ^{l}_{s}y \;\;\longrightarrow\;\;
\mathrm{sc}^{k}_{s}(x\circ^{l}_{s}y)\overset{\mathrm{e}}{=}\mathrm{sc}^{k}_{s}(x)\circ^{l}_{s}\mathrm{sc}^{k}_{s}(y);\\
x\circ^{l}_{s}y&\overset{\mathrm{e}}{=}x\circ^{l}_{s}y \;\;\longrightarrow\;\;
\mathrm{tg}^{k}_{s}(x\circ^{l}_{s}y)\overset{\mathrm{e}}{=}\mathrm{tg}^{k}_{s}(x)\circ^{l}_{s}\mathrm{tg}^{k}_{s}(y).
\label{DVarCompactABlk2}\tag{AB$^{l,k}$2}
\end{align*}

For every sort $s\in S$ and every quadruple of variables $x,y,z,t\in V^{S}_{s}$, if the $l$-compositions $x\circ^{l}_{s}y$ and $z\circ^{l}_{s}t$ are defined and the $k$-compositions $x\circ^{k}_{s}z$ and $y\circ^{k}_{s}t$ are defined, the $k$-composition of $x\circ^{l}_{s}y$ with $z\circ^{l}_{s}t$ is equal to the $l$-composition of $x\circ^{k}_{s}z$ with $y\circ^{k}_{s}t$. Formally,
\allowdisplaybreaks
\begin{multline*}
(x\circ^{l}_{s}y\overset{\mathrm{e}}{=}x\circ^{l}_{s}y)
\wedge
(z\circ^{l}_{s}t\overset{\mathrm{e}}{=}z\circ^{l}_{s}t)
\wedge
(x\circ^{k}_{s}z\overset{\mathrm{e}}{=}x\circ^{k}_{s}z)
\wedge
(y\circ^{k}_{s}t\overset{\mathrm{e}}{=}y\circ^{k}_{s}t)
\\
\;\;\longrightarrow\;\;
(x\circ^{l}_{s}y)\circ^{k}_{s}(z\circ^{l}_{s}t)\overset{\mathrm{e}}{=}
(x\circ^{k}_{s}z)\circ^{l}_{s}(y\circ^{k}_{s}t).
\label{DVarCompactABlk3}\tag{AB$^{l,k}$3}
\end{multline*}

A model of the axioms (A$^{k}$1)--(A$^{k}$6), for $k\in n$, and (AB$^{l,k}$1)--(AB$^{l,k}$3), for $l<k$ in $n$, is an $S$-sorted compact $n$-category, see Definition~\ref{DManySortedCatAlg}, presented with all the cells of a given sort, whatever their dimension, in a single set. For $n=2$, the axioms (A$^{0}$$i$), (A$^{1}$$i$) and (AB$^{0,1}$$i$) are, respectively, the axioms (A$i$), (B$i$) and (AB$i$) of the specification $\boldsymbol{\mathcal{E}}^{\boldsymbol{\mathcal{A}}^{(2)}}$ of Chapter~3.

We will let $\mathsf{PAlg}(\boldsymbol{\mathcal{C}}^{\mathbf{\Sigma}^{(n)}})$ stand for the category canonically associated to the $\mathrm{QE}$-variety $\mathcal{V}(\boldsymbol{\mathcal{C}}^{\mathbf{\Sigma}^{(n)}})$ determined by the specification $\boldsymbol{\mathcal{C}}^{\mathbf{\Sigma}^{(n)}}$.
\end{definition}

Sections~\ref{STruncation} and~\ref{SLimitCompact} also use the $\omega$-dimensional counterparts of these notions.

\begin{definition}\label{DCompactOmegaCat}
A \emph{compact $\omega$-category} is an ordered pair $\mathsf{C} = (C,(\xi_{k})_{k\in\omega})$ where $C$ is a set and, for every $k\in\omega$, $\xi_{k}$ a structure of compact category on $C$ such that, for every $j,k\in\omega$ with $j<k$, the ordered pair $(C,(\xi_{j},\xi_{k}))$ is a compact $2$-category, and such that every $f\in C$ is a $k$-cell, i.e., $\mathrm{sc}^{k}(f)=f$, for some $k\in\omega$. Note that, by~(2C1), a $j$-cell is a $k$-cell for every $k\geq j$.

A \emph{compact $\omega$-functor} is a mapping $F\colon C\mor C'$ which is a compact functor from $(C,\xi_{k})$ to $(C',\xi'_{k})$ for every $k\in\omega$. We denote by $\omega\mathsf{Cat}$ the category of compact $\omega$-categories and compact $\omega$-functors.
\end{definition}

\begin{definition}\label{DOmegaCatAlg}
Let $\mathbf{\Sigma}=(S,\Sigma)$ be a many-sorted signature. An \emph{$\omega$-categorial $\mathbf{\Sigma}$-algebra} is a pair $\boldsymbol{\mathsf{A}}=(\mathsf{A},F)$ where $\mathsf{A}=(\mathsf{A}_{s})_{s\in S}$ is an $S$-sorted compact $\omega$-category and $F$ assigns to every $\sigma\in\Sigma_{\mathbf{s},s}$ a compact $\omega$-functor $\sigma^{\boldsymbol{\mathsf{A}}}\colon\mathsf{A}_{\mathbf{s}}\mor\mathsf{A}_{s}$. Here $\mathsf{A}_{\mathbf{s}}$ is a compact $\omega$-category: if each $\mathsf{a}_{j}$ is a $k_{j}$-cell, then $(\mathsf{a}_{j})_{j\in\bb{\mathbf{s}}}$ is a $k$-cell for $k=\max_{j}k_{j}$. An \emph{$\omega$-categorial $\mathbf{\Sigma}$-homomorphism} is an $S$-sorted compact $\omega$-functor commuting with the operations, as in Definition~\ref{DManySortedCatAlg}. We denote by $\omega\mathsf{CatAlg}(\mathbf{\Sigma})$ the category of $\omega$-categorial $\mathbf{\Sigma}$-algebras and $\omega$-categorial $\mathbf{\Sigma}$-homomorphisms.
\end{definition}

%The axiom (A0) states that the operations of $\Sigma$ are total, and the axioms (A$^{k}$7) and (A$^{k}$8), for $k\in n$, state that they preserve sources, targets and compositions of every dimension, i.e., that they are $n$-functors. As in Chapters~2 and~3, the morphisms of $\mathsf{PAlg}(\boldsymbol{\mathcal{C}}^{\mathbf{\Sigma}^{(n)}})$ are the homomorphisms of partial $\Sigma^{(n)}$-algebras, i.e., the $S$-sorted mappings $h$ such that, whenever an operation is defined on a family of elements, it is also defined on the image of that family under $h$, and $h$ sends the result of the former to the result of the latter. These are the appropriate morphisms here: closed homomorphisms, which also reflect definedness, would be too restrictive, since an $n$-functor may send cells that are not composable to cells that are. Consequently, we have the following result.
%
%\begin{proposition}\label{PCompactVar}
%Let $\mathbf{\Sigma}$ be a many-sorted signature and $n\in\mathbb{N}$. Consider the functor from $\mathsf{nCatAlg}(\mathbf{\Sigma})$ to $\mathsf{PAlg}(\boldsymbol{\mathcal{C}}^{\mathbf{\Sigma}^{(n)}})$ that sends an $n$-categorial $\mathbf{\Sigma}$-algebra $\boldsymbol{\mathsf{A}}$ to the partial $\Sigma^{(n)}$-algebra whose underlying set of sort $s$ is the set of $n$-cells of $\mathsf{A}_{s}$, in which every $k$-cell, for $k\in n$, is identified with its identity $n$-cell, and whose operations are the operations of $\boldsymbol{\mathsf{A}}$ and the $k$-sources, $k$-targets and $k$-compositions of the $n$-categories $\mathsf{A}_{s}$; and that sends an $n$-categorial $\mathbf{\Sigma}$-homomorphism to its action on $n$-cells. Then this functor is an equivalence of categories, and it is an isomorphism of categories if $n$-categories are taken in their single-set presentation.
%\end{proposition}
%
%The proof reduces to the classical correspondence between the globular and the single-set presentations of $n$-categories, see, e.g.,~\cite{Str87,abgmmm23}, together with two observations: under this correspondence, the $n$-functors from $\mathsf{A}_{\mathbf{s}}$ to $\mathsf{A}_{s}$ are exactly the total operations that satisfy (A$^{k}$7) and (A$^{k}$8), for $k\in n$; and the $S$-sorted $n$-functors between the underlying $S$-sorted $n$-categories are exactly the homomorphisms of partial $\Sigma^{(n)}$-algebras defined above.
%
%In particular, the $\Sigma^{(1)}$-reducts of the partial algebras in $\mathbf{PAlg}(\boldsymbol{\mathcal{E}}^{\boldsymbol{\mathcal{A}}})$ of Chapter~2 and the $\Sigma^{(2)}$-reducts of those in $\mathbf{PAlg}(\boldsymbol{\mathcal{E}}^{\boldsymbol{\mathcal{A}}^{(2)}})$ of Chapter~3 are $1$-categorial and $2$-categorial $\mathbf{\Sigma}$-algebras, respectively: the specifications of those chapters are obtained from $\boldsymbol{\mathcal{C}}^{\mathbf{\Sigma}^{(1)}}$ and $\boldsymbol{\mathcal{C}}^{\mathbf{\Sigma}^{(2)}}$ by adding the rewrite rules as constants that are always defined.

\section{Extended presentation of higher-order categorial algebras}\label{SExtended}


We begin by recalling the definition of extended $\omega$-category from~\cite{CH0Ebis}, which presents $\omega$-categories with separate sorts for each dimension. An extended $n$-category is then an extended $\omega$-category truncated at level $n$.

\begin{definition}\label{DOmegaCatExt}
An \emph{extended $\omega$-category} is an ordered pair $\mathsf{C} = ((C_{k})_{k\in\omega},\zeta)$ consisting of an $\omega$-sorted set $(C_{k})_{k\in\omega}$ and a $\coprod_{k\in \omega}k$-indexed family
\[
\zeta = (\zeta_{j,k})_{(j,k)\in \coprod_{k\in \omega}k},
\] 
where, for every $k\in \omega$ and every $j\in k$, $\zeta_{j,k} = (\circ^{(j)},\mathrm{sc}^{(j,k)},\mathrm{tg}^{(j,k)},\mathrm{i}^{(k,j)})$ with
\begin{enumerate}
\item $\mathrm{sc}^{(j,k)}$ and $\mathrm{tg}^{(j,k)}$ mappings from $C_{k}$ to $C_{j}$  
which send $f\in C_{k}$ to its \emph{$(j,k)$-source} and \emph{$(j,k)$-target}, respectively;
\item $\mathrm{i}^{(k,j)}$ a mapping from $C_{j}$ to $C_{k}$ 
which sends $f\in C_{j}$ to its \emph{$(k,j)$-identity}; and
\item $\circ^{(j)}$ a partial binary operation 
$$
\circ^{(j)}\colon C_{k}\times C_{k} \dmor C_{k},
$$ 
whose domain of definition is
$$
\mathrm{Dom}\left(\circ^{(j)}\right) = 
\left\lbrace
(f,g)\in C_{k}\times C_{k}
\,\middle|\, \mathrm{sc}^{(j,k)}(g)=\mathrm{tg}^{(j,k)}(f) \right\rbrace,
$$ 
i.e., the fibered product of $C_{k}$ and $C_{k}$ over $\mathrm{sc}^{(j,k)}$ and $\mathrm{tg}^{(j,k)}$, and
which sends a pair $(f,g)\in \mathrm{Dom}(\circ^{(j)})$ to its \emph{$j$-composition in $C_{k}$}.
\end{enumerate}
These data are required to satisfy the following conditions for all $i,j,k\in \omega$ with $i<j<k$ (the conditions (MS2), (MS6) and (MS7) only involve the dimensions $j<k$, and are required for every such pair):
\begin{itemize}
\item[(MS1)] For every $f\in C_{k}$,
\begin{align*}
\mathrm{sc}^{(i,j)}(\mathrm{sc}^{(j,k)}(f))&=\mathrm{sc}^{(i,k)}(f)=\mathrm{sc}^{(i,j)}(\mathrm{tg}^{(j,k)}(f)),\\
\mathrm{tg}^{(i,j)}(\mathrm{tg}^{(j,k)}(f))&=\mathrm{tg}^{(i,k)}(f)=\mathrm{tg}^{(i,j)}(\mathrm{sc}^{(j,k)}(f)).
\end{align*}
\item[(MS2)] If $g,f\in C_{k}$ are such that $\mathrm{sc}^{(j,k)}(g)=\mathrm{tg}^{(j,k)}(f)$, then
\begin{align*}
\mathrm{sc}^{(j,k)}(g\circ^{(j)}f)&=\mathrm{sc}^{(j,k)}(f), &
\mathrm{tg}^{(j,k)}(g\circ^{(j)}f)&=\mathrm{tg}^{(j,k)}(g).
\end{align*}
\item[(MS3)] If $g,f\in C_{k}$ are such that $\mathrm{sc}^{(i,k)}(g)=\mathrm{tg}^{(i,k)}(f)$, then
\begin{align*}
\mathrm{sc}^{(j,k)}(g\circ^{(i)}f)&=\mathrm{sc}^{(j,k)}(g)\circ^{(i)}\mathrm{sc}^{(j,k)}(f),\\
\mathrm{tg}^{(j,k)}(g\circ^{(i)}f)&=\mathrm{tg}^{(j,k)}(g)\circ^{(i)}\mathrm{tg}^{(j,k)}(f).
\end{align*}
\item[(MS4)] If $f\in C_{i}$, then
$$
\mathrm{i}^{(k,j)}(\mathrm{i}^{(j,i)}(f))=\mathrm{i}^{(k,i)}(f).
$$
\item[(MS5)] If $f\in C_{i}$, then
$$
\mathrm{sc}^{(i,j)}(\mathrm{i}^{(j,i)}(f))=f=
\mathrm{tg}^{(i,j)}(\mathrm{i}^{(j,i)}(f)).
$$
\item[(MS6)] If $f\in C_{k}$, then
$$
f\circ^{(j)}
(\mathrm{i}^{(k,j)}(\mathrm{sc}^{(j,k)}(f)))
=
f
=
(\mathrm{i}^{(k,j)}(\mathrm{tg}^{(j,k)}(f)))
\circ^{(j)}
f.
$$
\item[(MS7)] If $h,g,f\in C_{k}$ are such that 
\begin{align*}
\mathrm{sc}^{(j,k)}(h)&=\mathrm{tg}^{(j,k)}(g); &
\mathrm{sc}^{(j,k)}(g)&=\mathrm{tg}^{(j,k)}(f),
\end{align*}
then
$$
h\circ^{(j)}(g\circ^{(j)}f)=
(h\circ^{(j)}g)\circ^{(j)}f.
$$
\item[(MS8)] If $g,f\in C_{j}$ are such that $\mathrm{sc}^{(i,j)}(g)=\mathrm{tg}^{(i,j)}(f)$, then
$$
\mathrm{i}^{(k,j)}(g\circ^{(i)}f)=\mathrm{i}^{(k,j)}(g)\circ^{(i)}\mathrm{i}^{(k,j)}(f),
$$
where on the left $\circ^{(i)}$ is the $i$-composition in $C_{j}$ and on the right the $i$-composition in $C_{k}$.
\item[(MS9)] If $g',g,f',f\in C_{k}$ are such that
\begin{align*}
\mathrm{sc}^{(j,k)}(g')&=\mathrm{tg}^{(j,k)}(g); &
\mathrm{sc}^{(j,k)}(f')&=\mathrm{tg}^{(j,k)}(f);\\
\mathrm{sc}^{(i,k)}(g')&=\mathrm{tg}^{(i,k)}(f'); &
\mathrm{sc}^{(i,k)}(g)&=\mathrm{tg}^{(i,k)}(f).
\end{align*}
then,
$$
(g'\circ^{(i)}f')\circ^{(j)}(g\circ^{(i)} f)
=
(g'\circ^{(j)}g)\circ^{(i)}(f'\circ^{(j)}f)
.
$$
\end{itemize}

An \emph{extended $\omega$-functor} $F$ from $\mathsf{C}$ to $\mathsf{C}'$ is a family $F=(F_{k})_{k\in\omega}$ of mappings $F_{k}\colon C_{k}\mor C'_{k}$ such that, for all $j,k\in\omega$ with $j<k$,
\begin{align*}
F_{j}\circ \mathrm{sc}^{(j,k)} &= \mathrm{sc}'^{(j,k)}\circ F_{k}, &
F_{j}\circ \mathrm{tg}^{(j,k)} &= \mathrm{tg}'^{(j,k)}\circ F_{k},\\
F_{k}\circ \mathrm{i}^{(k,j)} &= \mathrm{i}'^{(k,j)}\circ F_{j}, &
F_{k}(g\circ^{(j)}f) &= F_{k}(g)\circ^{(j)} F_{k}(f),
\end{align*}
the last one whenever $g\circ^{(j)}f$ is defined. We denote by $\omega\mathsf{Cat}^{\mathrm{ext}}$ the category of extended $\omega$-categories and extended $\omega$-functors.
\end{definition}

For $n\in\omega$, an \emph{extended $n$-category} is an ordered pair $\mathsf{C} = ((C_{k})_{k\in n+1},\zeta)$ consisting of an $n+1$-sorted set $(C_{k})_{k\in n+1}$ and a structure of extended $n$-category 
\[
\zeta = (\zeta_{j,k})_{(j,k)\in \coprod_{k\in n+1}k}
\] 
on $(C_{k})_{k\in n+1}$, where, for every $k\in n+1$ and every $j\in k$, $\zeta_{j,k} = (\circ^{(j)},\mathrm{sc}^{(j,k)},\mathrm{tg}^{(j,k)},\mathrm{i}^{(k,j)})$ with the same data as above, required to satisfy the same conditions truncated to the $n$-th level.

The notion of extended $n$-functor between $n$-categories is entirely analogous to that of extended $\omega$-functor between $\omega$-categories, but truncated to the $n$-th level. An extended $n$-functor $F$ from $\mathsf{C}$ to $\mathsf{C}'$ is a family $F=(F_{k})_{k\in n+1}$ of mappings $F_{k}\colon C_{k}\to C'_{k}$ such that for all $j,k\in n+1$ with $j<k$, the following equations hold (the last one whenever $g\circ^{(j)}f$ is defined):
\begin{align*}
F_{j}\circ \mathrm{sc}^{(j,k)} &= \mathrm{sc}'^{(j,k)}\circ F_{k}, &
F_{j}\circ \mathrm{tg}^{(j,k)} &= \mathrm{tg}'^{(j,k)}\circ F_{k},\\
F_{k}\circ \mathrm{i}^{(k,j)} &= \mathrm{i}'^{(k,j)}\circ F_{j}, &
F_{k}(g\circ^{(j)}f) &= F_{k}(g)\circ^{(j)} F_{k}(f).
\end{align*}
We denote by $\mathsf{nCat}^{\mathrm{ext}}$ the category consisting of extended $n$-categories and extended $n$-functors.


\begin{definition}\label{DSortednCatExt}
Let $S$ be a set of sorts and $n\in\mathbb{N}$.
We will call the objects of the functor category $(\mathsf{nCat}^{\mathrm{ext}})^{S}$ \emph{$S$-sorted extended $n$-categories}. Thus, an $S$-sorted extended $n$-category is an $S$-indexed family $\mathsf{A}=(\mathsf{A}_{s})_{s\in S}$ such that, for every sort $s\in S$, $\mathsf{A}_{s}$ is an extended $n$-category. Given two $S$-sorted extended $n$-categories 
$\mathsf{A}$ and $\mathsf{B}$, we will call the morphisms in $(\mathsf{nCat}^{\mathrm{ext}})^{S}$ from $\mathsf{A}$ to 
$\mathsf{B}$ \emph{$S$-sorted extended $n$-functors}. Thus, an $S$-sorted extended $n$-functor from $\mathsf{A}$ to $\mathsf{B}$ is an $S$-indexed family $F=(F_{s})_{s\in S}$ such that, for every sort $s\in S$, $F_{s}$ is an extended $n$-functor from $\mathsf{A}_{s}$ to $\mathsf{B}_{s}$. We will denote by $\mathrm{nFunc}^{\mathrm{ext}}(\mathsf{A},\mathsf{B})$ the set of all $S$-sorted extended $n$-functors from $\mathsf{A}$ to $\mathsf{B}$. We consider the expressions $F\in\mathrm{nFunc}^{\mathrm{ext}}(\mathsf{A},\mathsf{B})$ and $F\colon\mathsf{A}\mor\mathsf{B}$ as synonyms.
We recall that for $F\colon \mathsf{A}\mor\mathsf{B}$ and $G\colon \mathsf{B}\mor\mathsf{C}$, the \emph{composition} $G\circ F$ is $(G_{s}\circ F_{s})_{s\in S}$, and that for an $S$-sorted extended $n$-category $\mathsf{A}$, the \emph{identity} 
$S$-sorted extended $n$-functor at $\mathsf{A}$, that is, $\mathrm{Id}_{\mathsf{A}}$, is $(\mathrm{Id}_{\mathsf{A}_{s}})_{s\in S}$. 
\end{definition}

Next we define the corresponding notion of extended $n$-categorial algebra, following the style of the compact $n$-categorial algebra (Definition~\ref{DManySortedCatAlg}).

\begin{definition}\label{DExtCatAlg}
Let $\mathsf{A}$ be an $S$-sorted extended $n$-category. The $S^{\star}\times S$-sorted set of finitary operations on $\mathsf{A}$ is $(\mathrm{nFunc}^{\mathrm{ext}}(\mathsf{A}_{\mathbf{s}},\mathsf{A}_{s}))_{(\mathbf{s},s)\in S^{\star}\times S}$, where, for every $\mathbf{s}\in S^{\star}$, $\mathsf{A}_{\mathbf{s}}=\prod_{j\in\bb{\mathbf{s}}}\mathsf{A}_{s_{j}}$, with $\bb{\mathbf{s}}$ denoting the length of the word $\mathbf{s}$ (if $\mathbf{s}=\lambda$, then $\mathsf{A}_{\lambda}$ is a final extended $n$-category). 

Let $\mathbf{\Sigma}=(S,\Sigma)$ be a many-sorted signature. 
A \emph{structure of extended $n$-categorial $\Sigma$-algebra on} an $S$-sorted extended $n$-category $\mathsf{A}$ is a family $F = (F_{\mathbf{s},s})_{(\mathbf{s},s)\in S^{\star}\times S}$, where, for $(\mathbf{s},s)\in S^{\star}\times S$, $F_{\mathbf{s},s}$ is a mapping from $\Sigma_{\mathbf{s},s}$ to $\mathrm{nFunc}^{\mathrm{ext}}(\mathsf{A}_{\mathbf{s}},\mathsf{A}_{s})$ (if $(\mathbf{s},s) = (\lambda,s)$ and $\sigma\in \Sigma_{\lambda,s}$, then $F_{\lambda,s}(\sigma)$ picks out a $0$-cell $u$ of $\mathsf{A}_{s}$ together with its identities $\mathrm{i}^{(k,0)}(u)$, $k\in n+1$, since $\mathsf{A}_{\lambda}$ is a final extended $n$-category, with a single cell in each dimension). The finite products $\mathsf{A}_{\mathbf{s}}$ are taken dimensionwise: the $k$-cells of $\mathsf{A}_{\mathbf{s}}$ are the families of $k$-cells, and all the operations are computed componentwise. An \emph{$S$-sorted extended $n$-categorial $\Sigma$-algebra} or, simply, an \emph{extended $n$-categorial $\mathbf{\Sigma}$-algebra}, is a pair $(\mathsf{A},F)$, abbreviated to $\boldsymbol{\mathsf{A}}$, where $\mathsf{A}=(\mathsf{A}_{s})_{s\in S}$ is an $S$-sorted extended $n$-category and $F$ a structure of extended $n$-categorial $\Sigma$-algebra on $\mathsf{A}$. For a pair $(\mathbf{s},s)\in S^{\star}\times S$ and a formal operation $\sigma\in \Sigma_{\mathbf{s},s}$, in order to simplify the notation, the $n$-functor $F_{\mathbf{s},s}(\sigma)$ from $\mathsf{A}_{\mathbf{s}}$ to $\mathsf{A}_{s}$ will be written as $F_{\sigma}$, or simply as $\sigma^{\boldsymbol{\mathsf{A}}}$. 

An \emph{$S$-sorted extended $n$-categorial $\Sigma$-homomorphism} or, simply, an \emph{extended $n$-categorial $\mathbf{\Sigma}$-homomorphism}, from $\boldsymbol{\mathsf{A}}$ to $\boldsymbol{\mathsf{B}}$, where $\boldsymbol{\mathsf{B}}=(\mathsf{B},G)$, is a triple $(\boldsymbol{\mathsf{A}},F,\boldsymbol{\mathsf{B}})$, abbreviated to $F\colon\boldsymbol{\mathsf{A}}\mor\boldsymbol{\mathsf{B}}$, where $F$ is an $S$-sorted extended $n$-functor from $\mathsf{A}$ to $\mathsf{B}$ such that, for every $k\in n+1$, every $(\mathbf{s},s)\in S^{\star}\times S$, every $\sigma\in\Sigma_{\mathbf{s},s}$ and every family of $k$-cells $(\mathsf{a}_{j})_{j\in\bb{\mathbf{s}}}\in\mathsf{A}_{\mathbf{s}}$, we have that
$$
F_{s,k}\left(\sigma^{\boldsymbol{\mathsf{A}}}_{k}
\left(\left(\mathsf{a}_{j}
\right)_{j\in\bb{\mathbf{s}}}\right)\right)
=
\sigma^{\boldsymbol{\mathsf{B}}}_{k}
\left(\left(
F_{s_{j},k}\left(
\mathsf{a}_{j}
\right)
\right)_{j\in\bb{\mathbf{s}}}
\right),
$$
where $F_{s,k}$ is the component at $k$-cells of the extended $n$-functor $F_{s}$ and $\sigma^{\boldsymbol{\mathsf{A}}}_{k}$ that of $\sigma^{\boldsymbol{\mathsf{A}}}$. The equation is required for the $n$-cells too, the cells of top dimension.
We denote by $\mathsf{nCat}^{\mathrm{ext}}\mathsf{Alg}(\mathbf{\Sigma})$ the category of extended $n$-categorial  $\mathbf{\Sigma}$-algebras and extended $n$-categorial $\mathbf{\Sigma}$-homomorphisms.
\end{definition}

\begin{definition}\label{DSigmaCatExtended}
Let $\Sigma$ be an $S$-sorted signature and $n\in\mathbb{N}$. For every $k\in n+1$ and every $\mathbf{s}\in S^{\star}$, we let $(k,\mathbf{s})$ stand for the word $((k,s_{j}))_{j\in\lvert\mathbf{s}\rvert}$ in $((n+1)\times S)^{\star}$; in particular, $(k,\lambda)=\lambda$. The \emph{extended $n$-categorial signature determined by $\Sigma$}, denoted by $\Sigma^{\mathrm{ext}(n)}$, is the $(n+1)\times S$-sorted signature defined, for every $(\mathbf{w},(k,s))\in((n+1)\times S)^{\star}\times((n+1)\times S)$, as follows:
$$
{\Sigma}^{\mathrm{ext}(n)}_{\mathbf{w},(k,s)}=
\begin{cases}
\{\sigma_{k}\mid\sigma\in{\Sigma}_{\mathbf{s},s}\},
& \text{if $\mathbf{w}=(k,\mathbf{s})$ with $\mathbf{s}\neq(s,s)$;}\\
\{\sigma_{k}\mid\sigma\in{\Sigma}_{(s,s),s}\}\amalg\{\circ^{(j,k)}_{s}\mid j\in k\},
& \text{if $\mathbf{w}=((k,s),(k,s))$;}\\
\{\mathrm{sc}^{(k,m)}_{s}, \mathrm{tg}^{(k,m)}_{s}\},
& \text{if $\mathbf{w}=((m,s))$ with $k<m$;}\\
\{\mathrm{i}^{(k,j)}_{s}\},
& \text{if $\mathbf{w}=((j,s))$ with $j<k$;}\\
\varnothing,
& \text{otherwise.}
\end{cases}
$$
That is, the sort $(k,s)$ of ${\Sigma}^{\mathrm{ext}(n)}$ stands for the $k$-cells of sort $s$, and ${\Sigma}^{\mathrm{ext}(n)}$ is obtained from $\Sigma$ by taking, for every $k\in n+1$ and every operation symbol $\sigma\in\Sigma_{\mathbf{s},s}$, a copy $\sigma_{k}$ of $\sigma$ of arity $(k,\mathbf{s})$ and coarity $(k,s)$, which acts on the $k$-cells, and by adding, for every $j,k\in n+1$ with $j<k$ and every sort $s\in S$, two unary operation symbols $\mathrm{sc}^{(j,k)}_{s}$ and $\mathrm{tg}^{(j,k)}_{s}$, the \emph{$(j,k)$-source} and the \emph{$(j,k)$-target}, of arity $((k,s))$ and coarity $(j,s)$, a unary operation symbol $\mathrm{i}^{(k,j)}_{s}$, the \emph{$(k,j)$-identity}, of arity $((j,s))$ and coarity $(k,s)$, and a binary operation symbol $\circ^{(j,k)}_{s}$, the \emph{$j$-composition of $k$-cells}, of arity $((k,s),(k,s))$ and coarity $(k,s)$.
\end{definition}

In contrast with the compact case, $\Sigma^{\mathrm{ext}(n)}$ is not a subset of $\Sigma^{\mathrm{ext}(m)}$ for $n\leq m$, since the sets of sorts are different; however, the inclusion of $(n+1)\times S$ into $(m+1)\times S$ induces a canonical embedding of $\Sigma^{\mathrm{ext}(n)}$ into $\Sigma^{\mathrm{ext}(m)}$. Note also that, for $n=0$, $\Sigma^{\mathrm{ext}(0)}$ is, up to the identification of $1\times S$ with $S$, the signature $\Sigma$.

The specification of the extended presentation also follows the pattern of the specification $\boldsymbol{\mathcal{E}}^{\boldsymbol{\mathcal{A}}^{(2)}}$ of Chapter~3. Now, however, the role played in the compact case by a dimension $k\in n$ is played by a pair of dimensions $j<k$ in $n+1$, since the $j$-composition acts on the $k$-cells and the $(k,j)$-identities take the place of the $k$-sources and $k$-targets as unit elements; and the role played by a pair of dimensions $l<k$ in $n$ is played by a triple of dimensions $h<j<k$ in $n+1$. Informally, the equations state that the operations of $\Sigma$ are always defined, that, for every sort $s\in S$ and every pair $j<k$, the $j$-cells and the $k$-cells of sort $s$ are the objects and the morphisms of a category, that these categories fit together as in an $n$-category, and that the operations of $\Sigma$ are compatible with them. In what follows, we abbreviate $V^{(n+1)\times S}$ to $V$.

\begin{definition}\label{DVarExtended}
Let $\mathbf{\Sigma}=(S,\Sigma)$ be a many-sorted signature and $n\in\mathbb{N}$. We will denote by $((n+1)\times S,\Sigma^{\mathrm{ext}(n)},V^{(n+1)\times S},\mathcal{C}^{\mathbf{\Sigma}^{\mathrm{ext}(n)}})$, written $\boldsymbol{\mathcal{C}}^{\mathbf{\Sigma}^{\mathrm{ext}(n)}}$ for short, the specification in which $\Sigma^{\mathrm{ext}(n)}$ is the extended $n$-categorial signature introduced in Definition~\ref{DSigmaCatExtended}, $V^{(n+1)\times S}$ a fixed $(n+1)\times S$-sorted set with a countable infinity of variables in each coordinate, and $\mathcal{C}^{\mathbf{\Sigma}^{\mathrm{ext}(n)}}$ the subset of $\mathrm{QE}(\Sigma^{\mathrm{ext}(n)})_{V^{(n+1)\times S}}$, the set of all $\mathrm{QE}$-equations with variables in $V^{(n+1)\times S}$, consisting of the following equations:

For every $k\in n+1$, every $(\mathbf{s},s)\in S^{\star}\times S$, every operation symbol $\sigma\in \Sigma_{\mathbf{s},s}$, and every family of variables $(x_{p})_{p\in \lvert\mathbf{s}\rvert}\in V_{(k,\mathbf{s})}$, the operation $\sigma_{k}$ applied to the family $(x_{p})_{p\in \lvert\mathbf{s}\rvert}$ is always defined. Formally,
\allowdisplaybreaks
\begin{align*}
\sigma_{k}((x_{p})_{p\in \lvert\mathbf{s}\rvert})\overset{\mathrm{e}}{=}\sigma_{k}((x_{p})_{p\in \lvert\mathbf{s}\rvert}).\label{DVarExtendedEA0}\tag{EA0}
\end{align*}

For every $j,k\in n+1$ with $j<k$, the following equations describe the $j$-composition of $k$-cells.

For every sort $s\in S$, every variable $x\in V_{(k,s)}$ and every variable $u\in V_{(j,s)}$, the $(j,k)$-source and $(j,k)$-target of $x$ and the $(k,j)$-identity of $u$ are always defined. Formally,
\allowdisplaybreaks
\begin{align*}
\mathrm{sc}^{(j,k)}_{s}(x)&\overset{\mathrm{e}}{=}\mathrm{sc}^{(j,k)}_{s}(x);
&
\mathrm{tg}^{(j,k)}_{s}(x)&\overset{\mathrm{e}}{=}\mathrm{tg}^{(j,k)}_{s}(x);
&
\mathrm{i}^{(k,j)}_{s}(u)&\overset{\mathrm{e}}{=}\mathrm{i}^{(k,j)}_{s}(u).\label{DVarExtendedEAjk1}\tag{EA$^{j,k}$1}
\end{align*}

For every sort $s\in S$ and every variable $u\in V_{(j,s)}$, the $(j,k)$-source and the $(j,k)$-target of the $(k,j)$-identity of $u$ are equal to $u$. Formally,
\allowdisplaybreaks
\begin{align*}
\mathrm{sc}^{(j,k)}_{s}(\mathrm{i}^{(k,j)}_{s}(u))&\overset{\mathrm{e}}{=}u;
&
\mathrm{tg}^{(j,k)}_{s}(\mathrm{i}^{(k,j)}_{s}(u))&\overset{\mathrm{e}}{=}u.\label{DVarExtendedEAjk2}\tag{EA$^{j,k}$2}
\end{align*}
In other words, $\mathrm{i}^{(k,j)}_{s}$ is a common section of $\mathrm{sc}^{(j,k)}_{s}$ and $\mathrm{tg}^{(j,k)}_{s}$. This is the extended counterpart of the fact that $\mathrm{sc}^{k}_{s}$ and $\mathrm{tg}^{k}_{s}$ are right zeros.

For every sort $s\in S$ and every pair of variables $x,y\in V_{(k,s)}$, $x\circ^{(j,k)}_{s} y$ is defined if and only if the $(j,k)$-target of $y$ is equal to the $(j,k)$-source of $x$. Formally,
\allowdisplaybreaks
\begin{align*}
x\circ^{(j,k)}_{s} y\overset{\mathrm{e}}{=}x\circ^{(j,k)}_{s} y \;\;
&\longrightarrow \;\;
\mathrm{sc}^{(j,k)}_{s}(x)\overset{\mathrm{e}}{=}\mathrm{tg}^{(j,k)}_{s}(y);\\
\mathrm{sc}^{(j,k)}_{s}(x)\overset{\mathrm{e}}{=}\mathrm{tg}^{(j,k)}_{s}(y) \;\;
&\longrightarrow \;\;
x\circ^{(j,k)}_{s} y\overset{\mathrm{e}}{=}x\circ^{(j,k)}_{s}y.\label{DVarExtendedEAjk3}\tag{EA$^{j,k}$3}
\end{align*}

For every sort $s\in S$ and every pair of variables $x,y\in V_{(k,s)}$, if $x\circ^{(j,k)}_{s} y$ is defined, then the $(j,k)$-source of $x\circ^{(j,k)}_{s}y$ is that of $y$ and the $(j,k)$-target of $x\circ^{(j,k)}_{s}y$ is that of $x$. Formally,
\allowdisplaybreaks
\begin{align*}
x\circ^{(j,k)}_{s}y\overset{\mathrm{e}}{=}x\circ^{(j,k)}_{s}y \;\;
&\longrightarrow \;\;
\mathrm{sc}^{(j,k)}_{s}(x\circ^{(j,k)}_{s}y)\overset{\mathrm{e}}{=}\mathrm{sc}^{(j,k)}_{s}(y);\\
x\circ^{(j,k)}_{s}y\overset{\mathrm{e}}{=}x\circ^{(j,k)}_{s}y \;\;
&\longrightarrow \;\;
\mathrm{tg}^{(j,k)}_{s}(x\circ^{(j,k)}_{s}y)\overset{\mathrm{e}}{=}\mathrm{tg}^{(j,k)}_{s}(x).\label{DVarExtendedEAjk4}\tag{EA$^{j,k}$4}
\end{align*}

For every sort $s\in S$ and every variable $x\in V_{(k,s)}$, the compositions $x\circ^{(j,k)}_{s}\mathrm{i}^{(k,j)}_{s}(\mathrm{sc}^{(j,k)}_{s}(x))$ and $\mathrm{i}^{(k,j)}_{s}(\mathrm{tg}^{(j,k)}_{s}(x))\circ^{(j,k)}_{s}x$ are always defined and are equal to $x$, i.e., $\mathrm{i}^{(k,j)}_{s}(\mathrm{sc}^{(j,k)}_{s}(x))$ is a right unit element for the $j$-composition with $x$ and $\mathrm{i}^{(k,j)}_{s}(\mathrm{tg}^{(j,k)}_{s}(x))$ is a left unit element for the $j$-composition with $x$. Formally,
\allowdisplaybreaks
\begin{align*}
x\circ^{(j,k)}_{s}\mathrm{i}^{(k,j)}_{s}(\mathrm{sc}^{(j,k)}_{s}(x))&\overset{\mathrm{e}}{=} x;
&
\mathrm{i}^{(k,j)}_{s}(\mathrm{tg}^{(j,k)}_{s}(x))\circ^{(j,k)}_{s}x&\overset{\mathrm{e}}{=} x.\label{DVarExtendedEAjk5}\tag{EA$^{j,k}$5}
\end{align*}

For every sort $s\in S$ and every triple of variables $x, y, z\in V_{(k,s)}$, if the compositions $x\circ^{(j,k)}_{s}y$ and $y\circ^{(j,k)}_{s}z$ are defined, then the compositions $x\circ^{(j,k)}_{s}(y\circ^{(j,k)}_{s}z)$ and $(x\circ^{(j,k)}_{s}y)\circ^{(j,k)}_{s}z$ are defined and they are equal, i.e., the $j$-composition of $k$-cells, when defined, is associative. Formally,
\allowdisplaybreaks
\begin{multline*}
(x\circ^{(j,k)}_{s}y\overset{\mathrm{e}}{=} x\circ^{(j,k)}_{s}y)
\wedge
(y\circ^{(j,k)}_{s}z\overset{\mathrm{e}}{=} y\circ^{(j,k)}_{s}z)\;\;
\longrightarrow \;\;\\
(x\circ^{(j,k)}_{s}y)\circ^{(j,k)}_{s}z\overset{\mathrm{e}}{=} x\circ^{(j,k)}_{s}(y\circ^{(j,k)}_{s}z).\label{DVarExtendedEAjk6}\tag{EA$^{j,k}$6}
\end{multline*}

For every $(\mathbf{s},s)\in S^{\star}\times S$, every operation symbol $\sigma\in \Sigma_{\mathbf{s},s}$, every family of variables $(x_{p})_{p\in \lvert\mathbf{s}\rvert}\in V_{(k,\mathbf{s})}$ and every family of variables $(u_{p})_{p\in \lvert\mathbf{s}\rvert}\in V_{(j,\mathbf{s})}$, the $(j,k)$-source of $\sigma_{k}((x_{p})_{p\in \lvert\mathbf{s}\rvert})$ is equal to $\sigma_{j}$ applied to the family $((\mathrm{sc}^{(j,k)}_{s_{p}}(x_{p}))_{p\in \lvert\mathbf{s}\rvert})$, the $(j,k)$-target of $\sigma_{k}((x_{p})_{p\in \lvert\mathbf{s}\rvert})$ is equal to $\sigma_{j}$ applied to the family $((\mathrm{tg}^{(j,k)}_{s_{p}}(x_{p}))_{p\in \lvert\mathbf{s}\rvert})$, and the $(k,j)$-identity of $\sigma_{j}((u_{p})_{p\in \lvert\mathbf{s}\rvert})$ is equal to $\sigma_{k}$ applied to the family $((\mathrm{i}^{(k,j)}_{s_{p}}(u_{p}))_{p\in \lvert\mathbf{s}\rvert})$. Formally,
\allowdisplaybreaks
\begin{align*}
\mathrm{sc}^{(j,k)}_{s}(\sigma_{k}((x_{p})_{p\in \lvert\mathbf{s}\rvert}))&\overset{\mathrm{e}}{=}\sigma_{j}((\mathrm{sc}^{(j,k)}_{s_{p}}(x_{p}))_{p\in \lvert\mathbf{s}\rvert});\\
\mathrm{tg}^{(j,k)}_{s}(\sigma_{k}((x_{p})_{p\in \lvert\mathbf{s}\rvert}))&\overset{\mathrm{e}}{=}\sigma_{j}((\mathrm{tg}^{(j,k)}_{s_{p}}(x_{p}))_{p\in \lvert\mathbf{s}\rvert});\\
\mathrm{i}^{(k,j)}_{s}(\sigma_{j}((u_{p})_{p\in \lvert\mathbf{s}\rvert}))&\overset{\mathrm{e}}{=}\sigma_{k}((\mathrm{i}^{(k,j)}_{s_{p}}(u_{p}))_{p\in \lvert\mathbf{s}\rvert}).\label{DVarExtendedEAjk7}\tag{EA$^{j,k}$7}
\end{align*}

For every $(\mathbf{s},s)\in S^{\star}\times S$, every operation symbol $\sigma\in \Sigma_{\mathbf{s},s}$, and every pair of families of variables $(x_{p})_{p\in \lvert\mathbf{s}\rvert}, (y_{p})_{p\in \lvert\mathbf{s}\rvert}\in V_{(k,\mathbf{s})}$, if, for every $p\in \lvert\mathbf{s}\rvert$, the compositions $x_{p}\circ^{(j,k)}_{s_{p}}y_{p}$ are defined, then the composition $\sigma_{k}((x_{p})_{p\in \lvert\mathbf{s}\rvert})\circ^{(j,k)}_{s}\sigma_{k}((y_{p})_{p\in \lvert\mathbf{s}\rvert})$ is defined and it is equal to $\sigma_{k}$ applied to the family $(x_{p}\circ^{(j,k)}_{s_{p}}y_{p})_{p\in \lvert\mathbf{s}\rvert}$. Formally,
\allowdisplaybreaks
\begin{multline*}
\textstyle
\bigwedge_{p\in \lvert\mathbf{s}\rvert}
(x_{p}\circ^{(j,k)}_{s_{p}} y_{p}\overset{\mathrm{e}}{=}x_{p}\circ^{(j,k)}_{s_{p}} y_{p})
\;\;\longrightarrow\;\;\\
\sigma_{k}((x_{p}\circ^{(j,k)}_{s_{p}}y_{p})_{p\in \lvert\mathbf{s}\rvert})
\overset{\mathrm{e}}{=}
\sigma_{k}((x_{p})_{p\in \lvert\mathbf{s}\rvert})\circ^{(j,k)}_{s}\sigma_{k}((y_{p})_{p\in \lvert\mathbf{s}\rvert}).
\label{DVarExtendedEAjk8}\tag{EA$^{j,k}$8}
\end{multline*}

For every $h,j,k\in n+1$ with $h<j<k$, the following equations relate the $h$-, $j$- and $k$-cells.

For every sort $s\in S$ and every variable $x\in V_{(k,s)}$, the $(h,j)$-source of both the $(j,k)$-source and the $(j,k)$-target of $x$ is equal to the $(h,k)$-source of $x$, and the $(h,j)$-target of both the $(j,k)$-source and the $(j,k)$-target of $x$ is equal to the $(h,k)$-target of $x$. Formally,
\allowdisplaybreaks
\begin{align*}
\mathrm{sc}^{(h,j)}_{s}(\mathrm{sc}^{(j,k)}_{s}(x))&\overset{\mathrm{e}}{=}\mathrm{sc}^{(h,k)}_{s}(x);
&
\mathrm{tg}^{(h,j)}_{s}(\mathrm{tg}^{(j,k)}_{s}(x))&\overset{\mathrm{e}}{=}\mathrm{tg}^{(h,k)}_{s}(x);\\
\mathrm{sc}^{(h,j)}_{s}(\mathrm{tg}^{(j,k)}_{s}(x))&\overset{\mathrm{e}}{=}\mathrm{sc}^{(h,k)}_{s}(x);
&
\mathrm{tg}^{(h,j)}_{s}(\mathrm{sc}^{(j,k)}_{s}(x))&\overset{\mathrm{e}}{=}\mathrm{tg}^{(h,k)}_{s}(x).\label{DVarExtendedEAhjk1}\tag{EA$^{h,j,k}$1}
\end{align*}

For every sort $s\in S$ and every variable $u\in V_{(h,s)}$, the $(k,j)$-identity of the $(j,h)$-identity of $u$ is the $(k,h)$-identity of $u$. Formally,
\allowdisplaybreaks
\begin{align*}
\mathrm{i}^{(k,j)}_{s}(\mathrm{i}^{(j,h)}_{s}(u))&\overset{\mathrm{e}}{=}\mathrm{i}^{(k,h)}_{s}(u).\label{DVarExtendedEAhjk2}\tag{EA$^{h,j,k}$2}
\end{align*}

For every sort $s\in S$ and every pair of variables $x,y\in V_{(k,s)}$, if $x\circ^{(h,k)}_{s}y$ is defined, then the $(j,k)$-source of $x\circ^{(h,k)}_{s}y$ is the $h$-composition of the $(j,k)$-source of $x$ with the $(j,k)$-source of $y$, and the $(j,k)$-target of $x\circ^{(h,k)}_{s}y$ is the $h$-composition of the $(j,k)$-target of $x$ with the $(j,k)$-target of $y$. Formally,
\allowdisplaybreaks
\begin{align*}
x\circ^{(h,k)}_{s}y&\overset{\mathrm{e}}{=}x\circ^{(h,k)}_{s}y \;\;\longrightarrow\;\;
\mathrm{sc}^{(j,k)}_{s}(x\circ^{(h,k)}_{s}y)\overset{\mathrm{e}}{=}\mathrm{sc}^{(j,k)}_{s}(x)\circ^{(h,j)}_{s}\mathrm{sc}^{(j,k)}_{s}(y);\\
x\circ^{(h,k)}_{s}y&\overset{\mathrm{e}}{=}x\circ^{(h,k)}_{s}y \;\;\longrightarrow\;\;
\mathrm{tg}^{(j,k)}_{s}(x\circ^{(h,k)}_{s}y)\overset{\mathrm{e}}{=}\mathrm{tg}^{(j,k)}_{s}(x)\circ^{(h,j)}_{s}\mathrm{tg}^{(j,k)}_{s}(y).
\label{DVarExtendedEAhjk3}\tag{EA$^{h,j,k}$3}
\end{align*}

For every sort $s\in S$ and every quadruple of variables $x,y,z,t\in V_{(k,s)}$, if the compositions $x\circ^{(j,k)}_{s}y$, $z\circ^{(j,k)}_{s}t$, $x\circ^{(h,k)}_{s}z$ and $y\circ^{(h,k)}_{s}t$ are defined, the $(h,k)$-composition of $x\circ^{(j,k)}_{s}y$ with $z\circ^{(j,k)}_{s}t$ is equal to the $(j,k)$-composition of $x\circ^{(h,k)}_{s}z$ with $y\circ^{(h,k)}_{s}t$. Formally,
\allowdisplaybreaks
\begin{multline*}
(x\circ^{(j,k)}_{s}y\overset{\mathrm{e}}{=}x\circ^{(j,k)}_{s}y)
\wedge
(z\circ^{(j,k)}_{s}t\overset{\mathrm{e}}{=}z\circ^{(j,k)}_{s}t)
\wedge
(x\circ^{(h,k)}_{s}z\overset{\mathrm{e}}{=}x\circ^{(h,k)}_{s}z)
\wedge
(y\circ^{(h,k)}_{s}t\overset{\mathrm{e}}{=}y\circ^{(h,k)}_{s}t)
\\
\;\;\longrightarrow\;\;
(x\circ^{(j,k)}_{s}y)\circ^{(h,k)}_{s}(z\circ^{(j,k)}_{s}t)\overset{\mathrm{e}}{=}
(x\circ^{(h,k)}_{s}z)\circ^{(j,k)}_{s}(y\circ^{(h,k)}_{s}t).
\label{DVarExtendedEAhjk4}\tag{EA$^{h,j,k}$4}
\end{multline*}

For every sort $s\in S$ and every pair of variables $u,v\in V_{(j,s)}$, if $u\circ^{(h,j)}_{s}v$ is defined, then the $(k,j)$-identity of $u\circ^{(h,j)}_{s}v$ is the $h$-composition of the $(k,j)$-identities of $u$ and $v$. Formally,
\allowdisplaybreaks
\begin{align*}
u\circ^{(h,j)}_{s}v\overset{\mathrm{e}}{=}u\circ^{(h,j)}_{s}v \;\;\longrightarrow\;\;
\mathrm{i}^{(k,j)}_{s}(u\circ^{(h,j)}_{s}v)\overset{\mathrm{e}}{=}\mathrm{i}^{(k,j)}_{s}(u)\circ^{(h,k)}_{s}\mathrm{i}^{(k,j)}_{s}(v).
\label{DVarExtendedEAhjk5}\tag{EA$^{h,j,k}$5}
\end{align*}

The axioms (EA$^{h,j,k}$1)--(EA$^{h,j,k}$5) are the counterparts of the conditions (MS1), (MS4), (MS3), (MS9) and (MS8) of Definition~\ref{DOmegaCatExt}, and (EA$^{j,k}$1)--(EA$^{j,k}$6) those of (MS2), (MS5), (MS6) and (MS7). Hence a model of the axioms (EA$^{j,k}$1)--(EA$^{j,k}$6), for $j<k$ in $n+1$, and (EA$^{h,j,k}$1)--(EA$^{h,j,k}$5), for $h<j<k$ in $n+1$, is an $S$-sorted extended $n$-category, see Definition~\ref{DSortednCatExt}, presented in its globular form, i.e., with the cells of each dimension in a separate sort.

We will let $\mathsf{PAlg}(\boldsymbol{\mathcal{C}}^{\mathbf{\Sigma}^{\mathrm{ext}(n)}})$ stand for the category canonically associated to the $\mathrm{QE}$-variety $\mathcal{V}(\boldsymbol{\mathcal{C}}^{\mathbf{\Sigma}^{\mathrm{ext}(n)}})$ determined by the specification $\boldsymbol{\mathcal{C}}^{\mathbf{\Sigma}^{\mathrm{ext}(n)}}$.
\end{definition}

\begin{definition}\label{DOmegaExtCatAlg}
Let $\mathbf{\Sigma}=(S,\Sigma)$ be a many-sorted signature. An \emph{extended $\omega$-categorial $\mathbf{\Sigma}$-algebra} is a pair $\boldsymbol{\mathsf{A}}=(\mathsf{A},F)$ where $\mathsf{A}=(\mathsf{A}_{s})_{s\in S}$ is an $S$-sorted extended $\omega$-category and $F$ assigns to every $\sigma\in\Sigma_{\mathbf{s},s}$ an extended $\omega$-functor $\sigma^{\boldsymbol{\mathsf{A}}}\colon\mathsf{A}_{\mathbf{s}}\mor\mathsf{A}_{s}$. An \emph{extended $\omega$-categorial $\mathbf{\Sigma}$-homomorphism} is an $S$-sorted extended $\omega$-functor satisfying the equation of Definition~\ref{DExtCatAlg} for every $k\in\omega$. We denote by $\omega\mathsf{Cat}^{\mathrm{ext}}\mathsf{Alg}(\mathbf{\Sigma})$ the category of extended $\omega$-categorial $\mathbf{\Sigma}$-algebras and extended $\omega$-categorial $\mathbf{\Sigma}$-homomorphisms.
\end{definition}

%The axioms (EA0), (EA$^{j,k}$7) and (EA$^{j,k}$8) state, as in the compact case, that the operations of $\Sigma$ act layer by layer and compatibly with the $n$-categorial structure, i.e., that they are $n$-functors between $n$-categories in their globular presentation. The relation between the two presentations, however, does not come from these axioms, but from the correspondence between the globular and the single-set presentations of $n$-categories. A model of $\boldsymbol{\mathcal{C}}^{\mathbf{\Sigma}^{\mathrm{ext}(n)}}$ gives rise to a model of $\boldsymbol{\mathcal{C}}^{\mathbf{\Sigma}^{(n)}}$ by keeping only its $n$-cells, on which $k$-sources and $k$-targets are obtained by composing $(k,n)$-sources and $(k,n)$-targets with $(n,k)$-identities. Conversely, a model of $\boldsymbol{\mathcal{C}}^{\mathbf{\Sigma}^{(n)}}$ gives rise to a model of $\boldsymbol{\mathcal{C}}^{\mathbf{\Sigma}^{\mathrm{ext}(n)}}$ by taking, as $k$-cells of sort $s$, for $k\in n$, the cells $x$ with $\mathrm{sc}^{k}_{s}(x)=x$. Going from the compact presentation to the extended one and back gives the original algebra, whereas going from the extended presentation to the compact one and back only gives an isomorphic algebra, in which every $k$-cell has been replaced by its identity $n$-cell. It is therefore to be expected that $\mathsf{PAlg}(\boldsymbol{\mathcal{C}}^{\mathbf{\Sigma}^{\mathrm{ext}(n)}})$ is equivalent to $\mathsf{PAlg}(\boldsymbol{\mathcal{C}}^{\mathbf{\Sigma}^{(n)}})$, and hence to $\mathsf{nCatAlg}(\mathbf{\Sigma})$, although the two canonical functors are not mutually inverse; proving this equivalence in a way that is functorial in $\mathbf{\Sigma}$ is part of Objective~II.2.

\section{Graduation and Unification of Higher-Order Categorial Algebras}

\subsection{Graduation of higher-order categorial algebras}

\subsection{Unification of higher-order categorial algebras}

\subsection{On the relationship between graduation and unification}

\section{Discretization and Truncation of Higher-Order Categorial Algebras}\label{STruncation}

\subsection{Discretization of higher-order categorial algebras}

\subsection{Truncation of higher-order categorial algebras}

Truncating a higher-order categorial algebra at level $m$ keeps its cells of dimension at most $m$, together with the structure of the dimensions below $m$. In the compact presentation, where all the cells of a sort lie in a single set, these are the cells that are $k$-cells for every forgotten dimension $k$; in the extended presentation, where the cells of each dimension form a separate sort, they are the sorts of dimension at most $m$.

\begin{definition}\label{DTmnCatAlg}
Let $\mathbf{\Sigma}=(S,\Sigma)$ be a many-sorted signature and $m,n\in\mathbb{N}$ with $m\leq n$. For a compact $n$-category $\mathsf{C}=(C,(\xi_{k})_{k\in n})$, let
\[
T^{(m,n)}(C)=\{f\in C\mid \mathrm{sc}^{k}(f)=f\text{ for every }k\in n\text{ with }k\geq m\}.
\]
For $k<m$, the $k$-source and the $k$-target of any cell are in $T^{(m,n)}(C)$, by~(2C1), and, by~(2C2), so is $g\#^{k}f$ whenever $f,g\in T^{(m,n)}(C)$ and $\mathrm{sc}^{k}(g)=\mathrm{tg}^{k}(f)$. Hence $T^{(m,n)}(\mathsf{C})=(T^{(m,n)}(C),(\xi_{k}|_{T^{(m,n)}(C)})_{k\in m})$ is a compact $m$-category, and every compact $n$-functor $F\colon\mathsf{C}\mor\mathsf{C}'$, which preserves $k$-sources, restricts to a compact $m$-functor $T^{(m,n)}(F)\colon T^{(m,n)}(\mathsf{C})\mor T^{(m,n)}(\mathsf{C}')$.

The \emph{truncation functor} $T^{(m,n)}_{\mathbf{\Sigma}}\colon\mathsf{nCatAlg}(\mathbf{\Sigma})\mor\mathsf{mCatAlg}(\mathbf{\Sigma})$ sends an $n$-categorial $\mathbf{\Sigma}$-algebra $\boldsymbol{\mathsf{A}}=(\mathsf{A},F)$ to the $m$-categorial $\mathbf{\Sigma}$-algebra with underlying $S$-sorted compact $m$-category $(T^{(m,n)}(\mathsf{A}_{s}))_{s\in S}$ and with the restrictions of the operations $\sigma^{\boldsymbol{\mathsf{A}}}$, which preserve these cells because they are compact $n$-functors; and it sends an $n$-categorial $\mathbf{\Sigma}$-homomorphism to its restriction. In the same way, for $n\in\mathbb{N}$, the cells of a compact $\omega$-category that are $k$-cells for every $k\geq n$ form a compact $n$-category, and this defines the truncation functor $T^{(n,\omega)}_{\mathbf{\Sigma}}\colon\omega\mathsf{CatAlg}(\mathbf{\Sigma})\mor\mathsf{nCatAlg}(\mathbf{\Sigma})$.
\end{definition}

\begin{definition}\label{DTmnExtCatAlg}
Let $\mathbf{\Sigma}=(S,\Sigma)$ be a many-sorted signature and $m,n\in\mathbb{N}$ with $m\leq n$. For an extended $n$-category $\mathsf{C}=((C_{k})_{k\in n+1},\zeta)$, let $T^{\mathrm{ext}(m,n)}(\mathsf{C})=((C_{k})_{k\in m+1},(\zeta_{j,k})_{(j,k)\in\coprod_{k\in m+1}k})$, the extended $m$-category formed by the cells of dimension at most $m$; for an extended $n$-functor $F$, let $T^{\mathrm{ext}(m,n)}(F)=(F_{k})_{k\in m+1}$.

The \emph{truncation functor} $T^{\mathrm{ext}(m,n)}_{\mathbf{\Sigma}}\colon\mathsf{nCat}^{\mathrm{ext}}\mathsf{Alg}(\mathbf{\Sigma})\mor\mathsf{mCat}^{\mathrm{ext}}\mathsf{Alg}(\mathbf{\Sigma})$ sends an extended $n$-categorial $\mathbf{\Sigma}$-algebra $\boldsymbol{\mathsf{A}}$ to the extended $m$-categorial $\mathbf{\Sigma}$-algebra with underlying $S$-sorted extended $m$-category $(T^{\mathrm{ext}(m,n)}(\mathsf{A}_{s}))_{s\in S}$ and operations $T^{\mathrm{ext}(m,n)}(\sigma^{\boldsymbol{\mathsf{A}}})$, and an extended $n$-categorial $\mathbf{\Sigma}$-homomorphism $F$ to $(T^{\mathrm{ext}(m,n)}(F_{s}))_{s\in S}$. In the same way, keeping the cells of dimension at most $n$ of an extended $\omega$-categorial $\mathbf{\Sigma}$-algebra defines the truncation functor $T^{\mathrm{ext}(n,\omega)}_{\mathbf{\Sigma}}\colon\omega\mathsf{Cat}^{\mathrm{ext}}\mathsf{Alg}(\mathbf{\Sigma})\mor\mathsf{nCat}^{\mathrm{ext}}\mathsf{Alg}(\mathbf{\Sigma})$.
\end{definition}

\subsection{On the relationship between discretization and truncation}



\section{$\omega\mathsf{CatAlg}(\mathbf{\Sigma})$ as a projective limit}\label{SLimitCompact}

In this section we prove that the category $\omega\mathsf{CatAlg}(\mathbf{\Sigma})$ of $S$-sorted $\omega$-categorial $\Sigma$-algebras and $\omega$-categorial $\Sigma$-homomorphisms is the projective limit of the projective system $(\mathsf{nCatAlg}(\mathbf{\Sigma}),T^{(\bigcdot)}_{\mathbf{\Sigma}})$, where the functors $T^{(m,n)}_{\mathbf{\Sigma}}$ are the truncation functors defined in Definition~\ref{DTmnCatAlg} and $T^{(n,\omega)}_{\mathbf{\Sigma}}$ are the corresponding functors from $\omega\mathsf{CatAlg}(\mathbf{\Sigma})$ to $\mathsf{nCatAlg}(\mathbf{\Sigma})$.

\begin{remark}\label{RTmnCatAlgProjSys}
We let $(\bigcdot)\mathsf{CatAlg}(\mathbf{\Sigma})$ stand for the family of categories $(\mathsf{nCatAlg}(\mathbf{\Sigma}))_{n\in\omega}$ and we let $T^{(\bigcdot)}_{\mathbf{\Sigma}}$ stand for the family of functors $(T^{(m,n)}_{\mathbf{\Sigma}})_{(m,n)\in\leq}$. Note that $T^{(n,n)}_{\mathbf{\Sigma}}=\mathrm{Id}_{\mathsf{nCatAlg}(\mathbf{\Sigma})}$ and that, for every $m\leq n\leq p$,
\[
T^{(m,n)}_{\mathbf{\Sigma}}\circ T^{(n,p)}_{\mathbf{\Sigma}}=T^{(m,p)}_{\mathbf{\Sigma}},
\]
since a cell that is a $k$-cell for every $k\geq n$ is a $k$-cell for every $k\geq m$ exactly when, moreover, it is a $k$-cell for every $k$ with $m\leq k<n$.
Therefore, we obtain the projective system $((\bigcdot)\mathsf{CatAlg}(\mathbf{\Sigma}),T^{(\bigcdot)}_{\mathbf{\Sigma}})$ in $\mathsf{Cat}$ and its projective limit $\varprojlim ((\bigcdot)\mathsf{CatAlg}(\mathbf{\Sigma}),T^{(\bigcdot)}_{\mathbf{\Sigma}})$.
\end{remark}

Our next goal is to prove that such limit is isomorphic to $(\omega\mathsf{CatAlg}(\mathbf{\Sigma}),T^{(\bigcdot,\omega)}_{\mathbf{\Sigma}})$ with $T^{(\bigcdot,\omega)}_{\mathbf{\Sigma}}$ being the family of functors $(T^{(n,\omega)}_{\mathbf{\Sigma}})_{n\in\omega}$. That the composition of functors $T^{(m,n)}_{\mathbf{\Sigma}}\circ T^{(n,\omega)}_{\mathbf{\Sigma}}$ is equal to the functor $T^{(m,\omega)}_{\mathbf{\Sigma}}$ follows from the just stated consideration. Thus, all that remains to be proven is that $(\omega\mathsf{CatAlg}(\mathbf{\Sigma}),T^{(\bigcdot,\omega)}_{\mathbf{\Sigma}})$ satisfies the universal property.

\begin{proposition}\label{PUnivPropCatAlg}
Let $\mathsf{L}$ be a category and, for every $n\in\omega$, let $F_{n}$ be a functor from $\mathsf{L}$ to $\mathsf{nCatAlg}(\mathbf{\Sigma})$. If, for every $m,n\in\omega$ with $m\leq n$, $F_{m} = T^{(m,n)}_{\mathbf{\Sigma}} \circ F_{n}$, then there exists a unique functor $\langle F_{k}\rangle_{k\in\omega}$ from $\mathsf{L}$ to the category $\omega\mathsf{CatAlg}(\mathbf{\Sigma})$ such that, for every $n\in\omega$, $F_{n} = T^{(n,\omega)}_{\mathbf{\Sigma}} \circ \langle F_{k}\rangle_{k\in\omega}$. The diagram in Figure~\ref{FUnivPropCatAlg} depicts this situation.
\begin{figure}
\[
\xymatrix{
& \mathsf{L}\ar@{-->}[d]^{\exists!\langle F_{k}\rangle_{k\in\omega}}\ar@/_30pt/[ddl]_{F_{m}}\ar@/^30pt/[ddr]^{F_{n}} &
\\
& \omega\mathsf{CatAlg}(\mathbf{\Sigma})\ar[dr]^{T^{(n,\omega)}_{\mathbf{\Sigma}}}\ar[dl]_{T^{(m,\omega)}_{\mathbf{\Sigma}}} &
\\
\mathsf{mCatAlg}(\mathbf{\Sigma}) && \mathsf{nCatAlg}(\mathbf{\Sigma})\ar[ll]^{T^{(m,n)}_{\mathbf{\Sigma}}}
}
\]
\caption{Universal Property of $\omega\mathsf{CatAlg}(\mathbf{\Sigma})$.}
\label{FUnivPropCatAlg}
\end{figure}
\end{proposition}


\section{$\omega\mathsf{Cat}^{\mathrm{ext}}\mathsf{Alg}(\mathbf{\Sigma})$ as a projective limit}

In this section we prove that the category $\omega\mathsf{Cat}^{\mathrm{ext}}\mathsf{Alg}(\mathbf{\Sigma})$ of $S$-sorted extended $\omega$-categorial $\Sigma$-algebras and extended $\omega$-categorial $\Sigma$-homomorphisms is the projective limit of the projective system $(\mathsf{nCat}^{\mathrm{ext}}\mathsf{Alg}(\mathbf{\Sigma}),T^{\mathrm{ext}(\bigcdot)}_{\mathbf{\Sigma}})$, where the functors $T^{\mathrm{ext}(m,n)}_{\mathbf{\Sigma}}$ are the truncation functors defined in Definition~\ref{DTmnExtCatAlg} and $T^{\mathrm{ext}(n,\omega)}_{\mathbf{\Sigma}}$ are the corresponding functors from $\omega\mathsf{Cat}^{\mathrm{ext}}\mathsf{Alg}(\mathbf{\Sigma})$ to $\mathsf{nCat}^{\mathrm{ext}}\mathsf{Alg}(\mathbf{\Sigma})$.

\begin{remark}\label{RTmnExtCatAlgProjSys}
	We let $(\bigcdot)\mathsf{Cat}^{\mathrm{ext}}\mathsf{Alg}(\mathbf{\Sigma})$ stand for the family of categories $(\mathsf{nCat}^{\mathrm{ext}}\mathsf{Alg}(\mathbf{\Sigma}))_{n\in\omega}$ and we let $T^{\mathrm{ext}(\bigcdot)}_{\mathbf{\Sigma}}$ stand for the family of functors $(T^{\mathrm{ext}(m,n)}_{\mathbf{\Sigma}})_{(m,n)\in\leq}$. Note that $T^{\mathrm{ext}(n,n)}_{\mathbf{\Sigma}}=\mathrm{Id}_{\mathsf{nCat}^{\mathrm{ext}}\mathsf{Alg}(\mathbf{\Sigma})}$ and that, for every $m\leq n\leq p$,
	\[
	T^{\mathrm{ext}(m,n)}_{\mathbf{\Sigma}}\circ T^{\mathrm{ext}(n,p)}_{\mathbf{\Sigma}}=T^{\mathrm{ext}(m,p)}_{\mathbf{\Sigma}}.
	\]
	Therefore, we obtain the projective system $((\bigcdot)\mathsf{Cat}^{\mathrm{ext}}\mathsf{Alg}(\mathbf{\Sigma}),T^{\mathrm{ext}(\bigcdot)}_{\mathbf{\Sigma}})$ in $\mathsf{Cat}$ and its projective limit $\varprojlim ((\bigcdot)\mathsf{Cat}^{\mathrm{ext}}\mathsf{Alg}(\mathbf{\Sigma}),T^{\mathrm{ext}(\bigcdot)}_{\mathbf{\Sigma}})$.
\end{remark}

Our next goal is to prove that such limit is isomorphic to $(\omega\mathsf{Cat}^{\mathrm{ext}}\mathsf{Alg}(\mathbf{\Sigma}),T^{\mathrm{ext}(\bigcdot,\omega)}_{\mathbf{\Sigma}})$ with $T^{\mathrm{ext}(\bigcdot,\omega)}_{\mathbf{\Sigma}}$ being the family of functors $(T^{\mathrm{ext}(n,\omega)}_{\mathbf{\Sigma}})_{n\in\omega}$. That the composition of functors $T^{\mathrm{ext}(m,n)}_{\mathbf{\Sigma}}\circ T^{\mathrm{ext}(n,\omega)}_{\mathbf{\Sigma}}$ is equal to the functor $T^{\mathrm{ext}(m,\omega)}_{\mathbf{\Sigma}}$ follows from the just stated consideration. Thus, all that remains to be proven is that $(\omega\mathsf{Cat}^{\mathrm{ext}}\mathsf{Alg}(\mathbf{\Sigma}),T^{\mathrm{ext}(\bigcdot,\omega)}_{\mathbf{\Sigma}})$ satisfies the universal property.

\begin{proposition}\label{PUnivPropExtCatAlg}
	Let $\mathsf{L}$ be a category and, for every $n\in\omega$, let $F_{n}$ be a functor from $\mathsf{L}$ to $\mathsf{nCat}^{\mathrm{ext}}\mathsf{Alg}(\mathbf{\Sigma})$. If, for every $m,n\in\omega$ with $m\leq n$, $F_{m} = T^{\mathrm{ext}(m,n)}_{\mathbf{\Sigma}} \circ F_{n}$, then there exists a unique functor $\langle F_{k}\rangle_{k\in\omega}$ from $\mathsf{L}$ to the category $\omega\mathsf{Cat}^{\mathrm{ext}}\mathsf{Alg}(\mathbf{\Sigma})$ such that, for every $n\in\omega$, $F_{n} = T^{\mathrm{ext}(n,\omega)}_{\mathbf{\Sigma}} \circ \langle F_{k}\rangle_{k\in\omega}$. The diagram in Figure~\ref{FUnivPropExtCatAlg} depicts this situation.
	\begin{figure}
		\[
		\xymatrix{
			& \mathsf{L}\ar@{-->}[d]^{\exists!\langle F_{k}\rangle_{k\in\omega}}\ar@/_30pt/[ddl]_{F_{m}}\ar@/^30pt/[ddr]^{F_{n}} &
			\\
			& \omega\mathsf{Cat}^{\mathrm{ext}}\mathsf{Alg}(\mathbf{\Sigma})\ar[dr]^{T^{\mathrm{ext}(n,\omega)}_{\mathbf{\Sigma}}}\ar[dl]_{T^{\mathrm{ext}(m,\omega)}_{\mathbf{\Sigma}}} &
			\\
			\mathsf{mCat}^{\mathrm{ext}}\mathsf{Alg}(\mathbf{\Sigma}) && \mathsf{nCat}^{\mathrm{ext}}\mathsf{Alg}(\mathbf{\Sigma})\ar[ll]^{T^{\mathrm{ext}(m,n)}_{\mathbf{\Sigma}}}
		}
		\]
		\caption{Universal Property of $\omega\mathsf{Cat}^{\mathrm{ext}}\mathsf{Alg}(\mathbf{\Sigma})$.}
		\label{FUnivPropExtCatAlg}
	\end{figure}
\end{proposition}

\begin{proof}
Call a family $(\boldsymbol{\mathsf{A}}_{n})_{n\in\omega}$, with $\boldsymbol{\mathsf{A}}_{n}$ an extended $n$-categorial $\mathbf{\Sigma}$-algebra, \emph{compatible} if $T^{\mathrm{ext}(m,n)}_{\mathbf{\Sigma}}(\boldsymbol{\mathsf{A}}_{n})=\boldsymbol{\mathsf{A}}_{m}$ for all $m\leq n$. Its \emph{limit} $\varprojlim_{n}\boldsymbol{\mathsf{A}}_{n}$ is the extended $\omega$-categorial $\mathbf{\Sigma}$-algebra whose $k$-cells of sort $s$ are the $k$-cells of sort $s$ of $\boldsymbol{\mathsf{A}}_{k}$, its cells of top dimension. For $j<k$, the $(j,k)$-source of a $k$-cell $x$ is its $(j,k)$-source in $\boldsymbol{\mathsf{A}}_{k}$, which is a $j$-cell of $\boldsymbol{\mathsf{A}}_{k}$ and hence, by compatibility, a $j$-cell of $\boldsymbol{\mathsf{A}}_{j}$; the targets, identities and compositions are defined in the same way, and every $\sigma\in\Sigma_{\mathbf{s},s}$ acts on $k$-cells as $\sigma^{\boldsymbol{\mathsf{A}}_{k}}$ does. Each of the conditions of Definition~\ref{DOmegaCatExt} involves finitely many dimensions $i<j<k$, and holds because it holds in $\boldsymbol{\mathsf{A}}_{k}$. Likewise, a family of homomorphisms $(H_{n}\colon\boldsymbol{\mathsf{A}}_{n}\mor\boldsymbol{\mathsf{A}}'_{n})_{n\in\omega}$ between compatible families, with $T^{\mathrm{ext}(m,n)}_{\mathbf{\Sigma}}(H_{n})=H_{m}$ for $m\leq n$, has a limit $\varprojlim_{n}H_{n}$, which acts on $k$-cells as $H_{k}$. By construction, $T^{\mathrm{ext}(n,\omega)}_{\mathbf{\Sigma}}(\varprojlim_{m}\boldsymbol{\mathsf{A}}_{m})=\boldsymbol{\mathsf{A}}_{n}$ and $T^{\mathrm{ext}(n,\omega)}_{\mathbf{\Sigma}}(\varprojlim_{m}H_{m})=H_{n}$; and every extended $\omega$-categorial $\mathbf{\Sigma}$-algebra $\boldsymbol{\mathsf{B}}$, and every homomorphism, is the limit of its truncations: $\boldsymbol{\mathsf{B}}=\varprojlim_{n}T^{\mathrm{ext}(n,\omega)}_{\mathbf{\Sigma}}(\boldsymbol{\mathsf{B}})$.

For an object $X$ of $\mathsf{L}$ the family $(F_{n}(X))_{n\in\omega}$ is compatible, and for a morphism $f$ of $\mathsf{L}$ so is $(F_{n}(f))_{n\in\omega}$. Let $\langle F_{k}\rangle_{k\in\omega}(X)=\varprojlim_{n}F_{n}(X)$ and $\langle F_{k}\rangle_{k\in\omega}(f)=\varprojlim_{n}F_{n}(f)$. This is a functor, because every $F_{k}$ is, and $T^{\mathrm{ext}(n,\omega)}_{\mathbf{\Sigma}}\circ\langle F_{k}\rangle_{k\in\omega}=F_{n}$. If $G\colon\mathsf{L}\mor\omega\mathsf{Cat}^{\mathrm{ext}}\mathsf{Alg}(\mathbf{\Sigma})$ also satisfies $T^{\mathrm{ext}(n,\omega)}_{\mathbf{\Sigma}}\circ G=F_{n}$ for every $n$, then $G(X)=\varprojlim_{n}T^{\mathrm{ext}(n,\omega)}_{\mathbf{\Sigma}}(G(X))=\varprojlim_{n}F_{n}(X)$, and in the same way $G(f)=\varprojlim_{n}F_{n}(f)$, so that $G=\langle F_{k}\rangle_{k\in\omega}$.
\end{proof}


\section{An algebraic square}

We say that four adjunctions form an \emph{algebraic square} when there exist natural isomorphisms between the appropriate composites, as described in the following definition (taken from the Curry-Howard documents).

\begin{definition}\label{DefAlgSquare}
Let $\mathsf{C}$, $\mathsf{C}'$, $\mathsf{D}$ and $\mathsf{D}'$ be four categories. Moreover let
\begin{enumerate}
    \item
    $(G, F, \eta, \varepsilon)$ be an adjunction from $\mathsf{D}$ to $\mathsf{C}$.

    \item
    $(G', F', \eta', \varepsilon')$ be an adjunction from $\mathsf{D}'$ to $\mathsf{C}'$.

    \item
    $(K, J, \eta^{\mathsf{C}}, \varepsilon^{\mathsf{C}})$ be an adjunction from $\mathsf{C}'$ to $\mathsf{C}$.

    \item
    $(I, H, \eta^{\mathsf{D}}, \varepsilon^{\mathsf{D}})$ be an adjunction from $\mathsf{D}'$ to $\mathsf{D}$.
\end{enumerate}
We will say that these adjunctions form an \emph{algebraic square} if there exists a pair of natural isomorphisms $\alpha$ from $F' \circ J$ to $H \circ F$ and $\beta$ from $G \circ I$ to $K \circ G'$.

\begin{center}
\begin{tikzpicture}
    [ACliment/.style={-{To [angle'=45, length=5.75pt, width=4pt, round]}},
    RHACliment/.style={right hook-{To [angle'=45, length=5.75pt, width=4pt, round]}, font=\scriptsize},
    scale=1
    ]
    
    \node[] (C') at (0,0) {\textcolor{white}{$AAAA$}}; %Smu
    \node[] () at (0,0) {$\mathsf{C}'$};
    \node[] (D') at (3,0) {\textcolor{white}{$AAAA$}}; %Amu
    \node[] () at (3,0) {$\mathsf{D}'$};
    \node[] () at (1.5,0) {\rotatebox{90}{$\vdash$}};
    
    \draw[ACliment]  (C'.40) to node [above] {$F'$} (D'.140);
    \draw[ACliment]  (D'.220) to node [below] {$G'$} (C'.320);
    
    
    \node[] (C) at (0,3) {\textcolor{white}{$AAAA$}}; %Sz
    \node[] () at (0,3) {$\mathsf{C}$};
    \node[] (D) at (3,3) {\textcolor{white}{$AAAA$}}; %Az
    \node[] () at (3,3) {$\mathsf{D}$};
    \node[] () at (1.5,3) {\rotatebox{90}{$\vdash$}};
    \node[] () at (0,1.5) {\rotatebox{0}{$\vdash$}};
    \node[] () at (3,1.5) {\rotatebox{0}{$\vdash$}};

    \node[] () at (1.5,1.5) {$(\alpha, \beta)$};
    
    \draw[ACliment]  (C.40) to node [above] {$F$} (D.140);
    \draw[ACliment]  (D.220) to node [below] {$G$} (C.320);
    
    \draw[ACliment]  (C.310) to node [right] {$J$} (C'.50);
    \draw[ACliment]  (C'.140) to node [left] {$K$} (C.220);
    
    \draw[ACliment]  (D.310) to node [right] {$H$} (D'.50);
    \draw[ACliment]  (D'.140) to node [left] {$I$} (D.220);
\end{tikzpicture}
\end{center}
\end{definition}

%%%%%%%%%%%%%%%%%%%%%%%%%%%%%%%%%%%%%%%%%%%%%%%
\begin{thebibliography}{900}

%\bibitem%[ja90]
%{ja90} J. Almeida, \emph{On pseudovarieties, varieties of languages, filters of congruences, pseudoidentities and related topics}. Algebra Universalis, \textbf{27} (1990), pp.~333--350.

\bibitem%[AntFLT26]
{AntFLT26} Anthropic, \emph{Formalizing Fermat's Last Theorem}. Anthropic Research, 4 September 2026. \url{https://www.anthropic.com/research/formalizing-fermats-last-theorem}.

\bibitem%[AntFLTRepo]
{AntFLTRepo} Anthropic, \emph{Fermat's Last Theorem in Lean}. Code repository. \url{https://github.com/anthropics/fermats-last-theorem}.

\bibitem%[AGMAI26]
{AGMAI26} Advisory Group on Mathematics and Artificial Intelligence, \emph{Responsible release of AI-generated mathematics}. 29 September 2026. \url{https://agmai.org/general-sep29/}.

\bibitem%[ABS00]
{ABS00}
F.~A. Al-Agl, R. Brown and R. Steiner, \emph{Multiple categories: the equivalence of a globular and a cubical approach}. arXiv preprint math/0007009 (2000).

\bibitem%[Lei26]
{Lei26} J. Alper, M. Barany, A. Chavarri Villarello, S. Dahmen, W. Dean, K. Ganapathy, M. Harris, D. Holmes, M. Jamnik, S. Kelk, B. Kra, U. Martin, B. Naskr\c{e}cki, R. Ochigame, J. Portegies and J. Schmitt, \emph{Leiden Declaration on Artificial Intelligence and Mathematics}. June 2026. arXiv preprint arXiv:2609.35334 (2026). \url{https://leidendeclaration.ai}.

\bibitem%[an83]
{an83} H. Andr\'{e}ka and I. N\'{e}meti, \emph{Generalization of the concept of variety and quasivariety to partial algebras through category theory.} Dissertationes Math. (Rozprawy Mat.) \textbf{204} (1983).

\bibitem%[abgmmm23]
{abgmmm23} D. Ara, A. Burroni, Y. Guiraud, P. Malbos, F. M\'{e}tayer and S. Mimram, \emph{Polygraphs: from rewriting to higher categories.} arXiv preprint arXiv:2312.00429 (2023).

\bibitem%[Ash26]
{Ash26} J. Asher, \emph{Formalize everything, now!} Blog post, 28 September 2026. \url{https://justinasher.me/formalize-everything-now}.

\bibitem%[BW85]
{BW85} M. Barr and Ch. Wells, \emph{Toposes, triples and theories}. Springer-Verlag, New York, 1985. 

\bibitem%[BW12]
{BW12} M. Barr and Ch. Wells, Category theory for computing science, Prentice-Hall International Series in Computer Science (1995); reprinted in: Reprints in Theory and Applications of Categories 22 (2012).


\bibitem%[jb67]
{jb67} J. B\'{e}nabou, \emph{Introduction to bicategories}. In Reports of the Midwest Category Seminar, volume 47 of Lecture
Notes in Mathematics, pp.~1--77. Springer, 1967.

\bibitem%[jb68]
{jb68} J. B\'{e}nabou, \emph{Structures alg\'{e}briques dans les cat\'{e}gories}. Cahiers Topologie G\'{e}om. Diff\'{e}rentielle \textbf{10} (1968),  pp.~1--126.

\bibitem{BN98}
F. Baader and T. Nipkow, \emph{Term rewriting and all that}. Cambridge University Press, 1999.

\bibitem{Ben70}
D.~B. Benson, \emph{Syntax and semantics: A categorical view}. Information and Control,
\textbf{17}, 1970, pp. 145-160.

\bibitem{Ben74}
D.~B. Benson, \emph{Semantic preserving translations}. Mathematical Systems Theory \textbf{8} (1974),  pp.~105--126.

\bibitem%[db75]
{db75} D.~B. Benson, \emph{The basic algebraic structures in categories of derivations}. Information and Control, \textbf{28} (1975), pp.~2--29.

\bibitem%[gb15]
{gb15} G.~M. Bergman, \emph{An invitation to general algebra and universal constructions}. Second edition. Universitext. Springer, Cham, 2015.

\bibitem%{BKV03}
{BKV03} M. Bezem, J.~W. Klop and R.~C. De Vrijer (Eds), \emph{Term rewriting systems. Terese}, Cambridge University Press, 2003.    

\bibitem%[bl05]
{bl05} T.~S. Blyth, \emph{Lattices and ordered algebraic structures}. Universitext. Springer-Verlag London, Ltd., London, 2005.

\bibitem%[arXiv26]
{arXiv26} K. Boboris, \emph{Fair moderation, equitable access, and AI: arXiv's updated rate limit policy}. arXiv blog, 1 October 2026. \url{https://blog.arxiv.org/2026/10/01/updated-rate-limit-policy/}.

\bibitem%[BO93]
{BO93} R.~V. Book and F. Otto, \emph{String-rewriting systems}.
Texts and Monographs in Computer Science. Springer-Verlag, New York, 1993. 

\bibitem%[nB70]
{nB70} N. Bourbaki, \emph{Th{\'{e}}orie des ensembles}. Hermann, Paris, 1970.

\bibitem%[DB17]
{DB17} D. Bourn, \emph{From groups to categorial algebra. Introduction to protomodular and Mal'tsev categories}. Compact Textbooks in Mathematics. Birkh\"{a}user/Springer, Cham, 2017.

%\bibitem%[Bra69]
%{Bra69} W.~S. Brainerd, \emph{Tree generating regular systems}. IC \textbf{14} (1969), pp.~217--231.

\bibitem%[BH81]
{BH81} R. Brown and P.~J. Higgins, \emph{The equivalence of $\infty $-groupoids and crossed complexes}. Cahiers de topologie et g\'{e}om\'{e}trie diff\'{e}rentielle, 22(4) (1981), pp.~371--386.

\bibitem%[rb89]
{rb89} J.~R. B\"{u}chi, \emph{Finite automata, their algebras and grammars. Towards a theory of formal expressions.} Edited and with a preface by Dirk Siefkes. Springer-Verlag, New York, 1989.

\bibitem%[bs67]
{bs67} P. Burmeister and J. Schmidt, \emph{On the completion of partial algebras}. Colloq. Math. \textbf{17} (1967), pp.~235--245.

\bibitem%[pb70]
{pb70} P. Burmeister, \emph{Free partial algebras}. J. Reine Angew. Math. \textbf{241} (1970), pp.~75--86.

\bibitem%[pb82]
{pb82} P. Burmeister, \emph{Partial algebras--survey of a unifying approach towards a twovalued
model theory for partial algebras}. Algebra Universalis, \textbf{15} (1982), pp.~306--358.

\bibitem%[pb86]
{pb86} P. Burmeister, \emph{A model theoretic oriented approach to partial algebras. Introduction to theory and application of partial algebras. Part I}. Mathematical Research, 32. Akademie-Verlag, Berlin, 1986.

\bibitem%[pbbw92a]
{pbw92a} P. Burmeister and B. Wojdy{\l}o, \emph{The meaning of basic category-theoretical notions in some categories of partial algebras. I}. Demonstratio Math. \textbf{25} (1992), pp.~583--602.

\bibitem%[pbw92b]
{pbw92b} P. Burmeister and B. Wojdy{\l}o, \emph{The meaning of basic category theoretical notions in some categories of partial algebras. II}. Products and coproducts. Demonstratio Math. \textbf{25} (1992), 
pp.~973--994.

\bibitem%[pb93]
{pb93}P. Burmeister, Partial algebras--{A}n introductory survey. Algebras and orders (Montreal, PQ, 1991), pp.~1--70, NATO Adv. Sci. Inst. Ser. C: Math. Phys. Sci., 389, Kluwer Acad. Publ., Dordrecht, 1993.

\bibitem%[pbw96]
{pbw96} P. Burmeister and B. Wojdy{\l}o, \emph{The meaning of basic category theoretical notions in some categories of partial algebras. III}. Demonstratio Math. \textbf{29} (1996), pp.~825--848.

\bibitem%[pb02]
{pb02} P. Burmeister, \emph{Lecture Notes on Universal Algebra--Many-Sorted Partial Algebras}. Fragment of lectures hold in the summer terms of 1994, 1996, 1998, 2000 and 2002 at the Darmstadt University of Technology. Download via http://www.mathematik.tu-darmstadt.de/{\raise.17ex\hbox{$\scriptstyle\mathtt{\sim}$}}burmeister/

\bibitem%[bs81]
{bs81} S. Burris and H.~P. Sankappanavar, \emph{A course in universal algebra}. Springer-Verlag, New York-Berlin, 1981.

\bibitem%[Buz26]
{Buz26} K. Buzzard, \emph{To grieve, or not to grieve?} Xena (blog), 1 October 2026. \url{https://xenaproject.wordpress.com/2026/10/01/to-grieve-or-not-to-grieve/}.

\bibitem{Bur91}
A. Burroni, \emph{Higher-dimensional word problem.} Category theory and computer science (Paris, 1991), pp.~ 94--105, Lecture Notes in Comput. Sci., \textbf{530}, Springer, Berlin, 1991.

\bibitem{Bur93}
A. Burroni, \emph{Higher-dimensional word problems with applications to equational logic}. Theoretical computer science \textbf{115} (1993), pp.~43-62.

\bibitem%[CH06]
{CH06} F. Cardone and J.~R. Hindley, \emph{History of lambda-calculus and combinatory logic}. In: D.M. Gabbay and J. Woods (Eds.) Handbook of the History of Logic.  Vol.~\textbf{5}. Logic from Russell to Church (2009), pp.~723--817. Elsevier/North-Holland, Amsterdam, 2009.

\bibitem%[CM82]
{CM82} B. Chandler and W. Magnus, \emph{The History of Combinatorial Group Theory}. Sources in the History of Mathematics and Physics Sciences \textbf{9}, 1982. 

\bibitem%[Chu32]
{ChuI32} A. Church, \emph{A set of postulates for the foundation of logic}. Ann. of Math. \textbf{33}(2) (1932), pp.~346--366.

\bibitem%[Chu32]
{ChuII32} A. Church, \emph{A set of postulates for the foundation of logic}. (Second paper.)  Ann. of Math. \textbf{34} (2) (1932), pp.~839--864.

\bibitem%[Chu36]
{Chu36} A. Church, \emph{An unsolvable problem of elementary number theory}. American Journal of Mathematics, \textbf{58} (2) (1936), pp.~345--363.

\bibitem%[Chu41]
{Chu41} A. Church, \emph{The calculi of lambda-conversion}. Ann. of Math. Studies \textbf{6}, Princeton N.~J., 1941.

\bibitem%[cgw95]
{cgw95} I. Cla{\ss}en, M. Gro{\ss}e-Rhode and U. Wolter, \emph{Categorical concepts for parameterized partial specifications}. ADT-COMPASS Workshop and TAPSOFT Conference (Caldes de Malavella, 1992; Orsay, 1993). Math. Structures Comput. Sci. \textbf{5} (1995), pp.~153--188.

\bibitem{CM15}
F. Clerc  and S. Mimram, \emph{Presenting a category modulo a rewriting system}. 26th International Conference on Rewriting Techniques and Applications (RTA' 2015). Editor: Maribel Fern\'{a}ndez; pp.~89--105.  Schloss Dagstuhl-Leibniz-Zentrum f\"ur Informatik, 2015.

\bibitem%[BCM20]
{BCM20} K. Buzzard, J. Commelin and P. Massot, \emph{Formalising perfectoid spaces}. In Proceedings of the 9th ACM SIGPLAN International Conference on Certified Programs and Proofs (CPP 2020), pp.~299--312. ACM, 2020.

% [No citada des del 04/10/2026, en escurçar la secció 7.5]
%\bibitem%[CIM26]
%{CIM26} M. Codreanu, A. Imas, J. Mateos-Garcia et al., \emph{AI in science: early insights}. Google, Google DeepMind and MIT FutureTech, September 2026. \url{https://ai.google/static/documents/AI-in-Science.pdf}.

% [No citada des del 04/10/2026, en escurçar la secció 7.5]
%\bibitem%[AIM26]
%{AIM26} M. Colbrook, \emph{AIM: an invitation to explore mathematics together}. Guest post on T.~Tao's blog, on behalf of S. Brunton, M. Colbrook, M.~V. de Hoop, G. Stepaniants, A. Townsend and R. Ward, 2 October 2026. \url{https://terrytao.wordpress.com/2026/10/02/aim-an-invitation-to-explore-mathematics-together/}.

\bibitem%[CLG25]
{CLG25} E. Cosme Ll\'{o}pez and L. Gong, \emph{An introduction to Lean~4}. Zenodo, 2025. \url{https://doi.org/10.5281/zenodo.17079867}.

\bibitem%[CosRepo]
{CosRepo} E. Cosme Ll\'{o}pez, \emph{Mslang: Lean~4 formalisation of Eilenberg theorems for many-sorted formations}. Code repository. \url{https://github.com/encosllo/MSLang}.

\bibitem%[CosArch]
{CosArch} E. Cosme Ll\'{o}pez, \emph{Architecture: concurrent mathematical development and formal verification (local-first)}. Specification, revision~6. \url{https://github.com/encosllo/MSLang/blob/main/Architecture.md}.

\bibitem%[CVST04]
{CVST04} J. Climent Vidal and J. Soliveres Tur,\emph{ On many-sorted algebraic closure operators}. Math. Nachr. \textbf{266} (2004), pp.~81--84. 

\bibitem%[CVST10a]
{CVST10a} J. Climent Vidal and J. Soliveres Tur,\emph{ Birkhoff-Frink representations as functors}. Math. Nachr. \textbf{283} (2010), pp.~686--703. 

\bibitem%[CVST10b]
{CVST10b} J. Climent Vidal and J. Soliveres Tur,\emph{ A 2-categorical framework for the syntax and semantics of many-sorted equational logic.} Rep. Math. Logic  \textbf{45} (2010), pp.~37--95.

\bibitem%[CVCL16]
{CVCL16} J. Climent Vidal and E. Cosme Ll\'{o}pez, \emph{Eilenberg theorems for many-sorted formations}.  Houston Journal of Mathematics, \textbf{45} (2019), pp.~321--369.

\bibitem%[CVCL18]
{CVCL18} J. Climent Vidal and E. Cosme Ll\'{o}pez, \emph{Congruence based proofs of the recognizability theorems for free many-sorted algebras}. Journal of Logic and Computation, \textbf{30} (2020), pp.~561--633.

\bibitem%[CVCL23]
{CVCL23} J. Climent Vidal and E. Cosme Ll\'{o}pez, \emph{Functoriality of the Schmidt Construction}. Logic Journal of the IGPL, \textbf{31}(5) (2023), pp.~822--893.

\bibitem%[CVCLRM25]
{CVCLRM25} J. Climent Vidal, E. Cosme Ll\'{o}pez and R. Ruiz Mora, \emph{Riguet and generalized congruences on a category. Relationships and applications}. Preprint (2025), arXiv:2505.15767.

\bibitem%[CVCLRM26]
{CVCLRM26}
Climent Vidal, J., Cosme Llópez, E., Ruiz Mora, R.: From higher-order rewriting systems to higher-order categorial algebras and higher-order Curry-Howard isomorphisms (2026), arXiv:2402.12051.

%\bibitem%[CS12]
%{CS12}
%J. Climent Vidal and J. Soliveres Tur, \emph{On the morphisms and transformations of Tsuyoshi Fujiwara (as a concretion of a bidimensional many-sorted general algebra and its application to the equivalence between many-sorted clones and algebraic theories)}, eprint arXiv:1201.5239, 2012.

\bibitem%[Coh81]
{pC81} P. Cohn, \emph{Universal algebra}, D. Reidel Publishing Company, 1981.

%\bibitem%[Coq86]
%{Coq86} T. Coquand and G. Huet, \emph{The calculus of constructions}. Information and Control, \textbf{76} (1988), pp.~95--120.

\bibitem%[CurI30]
{CurI30} H.~B. Curry, \emph{Grundlagen der Kombinatorischen Logik. Teil I}. American Journal of Mathematics, \textbf{52} (3) (1930), pp.~509--536.

\bibitem%[CurII30]
{CurII30} H.~B. Curry, \emph{Grundlagen der Kombinatorischen Logik. Teil II}. American Journal of Mathematics, \textbf{52} (4) (1930), pp.~789--834.

\bibitem%[CF58]
{CF58} H.~B. Curry and R. Feys, \emph{Combinatory Logic, vol. $\mathrm{I}$}. Studies in Logic and the Foundations of Mathematics. North-Holland Publishing Co., Amsterdam, 1958.

%\bibitem%[Tata08]
%{Tata08} H. Comon, M. Dauchet, R. Gilleron, F. Jacquemard, D. Lugiez, S. Tison, C. L\"{o}ding, and M. Tommasi,
%\emph{Tree automata techniques and Applications}. Available on: {\tt http://tata.gforge.inria.fr}, 2008.

%\bibitem%[Cou89]
%{Cou89} B. Courcelle, \emph{On recognizable sets and tree automata}. In H. A\"{\i}t-Kaci and M. Nivat, editors. \emph{Algebraic Techniques: Resolution of Equations in Algebraic Structures}, Chapter 3, pp.~93--126, Academic press, Inc., 1989.

%\bibitem%[Cou96]
%{Cou96} B. Courcelle, \emph{Basic notions of universal algebra for language theory and graph grammars}, Theoret. Comput. Sci. \textbf{163} (1996), pp.~1--54.
%

\bibitem%[Die26]
{Die26} T.~G. Dietterich, Post on X, 13 September 2026. \url{https://x.com/tdietterich/status/2099197831580090622}.

\bibitem%[DP02]
{DP02} B.~A. Davey and H.~A. Priestley, \emph{Introduction to lattices and order}. Second edition. Cambridge University Press, New York, 2002.

\bibitem%[dMU21]
{dMU21} L. de Moura and S. Ullrich, \emph{The Lean~4 theorem prover and programming language}. In Automated Deduction -- CADE~28, Lecture Notes in Computer Science \textbf{12699}, pp.~625--635. Springer, 2021.

\bibitem%[D82]
{D82} M. Davis, \emph{Computability and unsolvability}. (Enl. version of the orig. publ. McGraw- Hill Book Company, New York 1958). Dover Publications, Inc. 
New York, 1982.

\bibitem%[D04]
{D04} M. Davis (ed.), \emph{The undecidable. Basic papers on undecidable propositions, unsolvable problems and computable functions}. Corrected reprint of the 1965 original. Dover Publications, Inc., Mineola, New York, 2004.

%\bibitem%[d63]
%{d63} R. Dedekind, \emph{Essays on the theory of numbers. I: Continuity and irrational numbers. II: The nature and meaning of numbers}. Authorized translation by Wooster Woodruff Beman Dover Publications, Inc., New York, 1963.

\bibitem%[MD11]
{MD11} M. Dehn, \emph{\"{U}ber unendliche diskontinuierliche Gruppen}, Mathematische Annalen \textbf{71} (1911), pp.~116--144.

%\bibitem%[Don65]
%{Don65} J.~E. Doner, \emph{Decidability of the weak second-order theory of two successors}. Notices Amer. Math. Soc. \textbf{12} (1965), Abstract 65T-468, p.~819.
%
%\bibitem%[Don70]
%{Don70} J.~E. Doner, \emph{Tree acceptors and some of their applications}. J. CSS \textbf{4} (1970), pp.~406--451.

\bibitem{Ehr64}
C. Ehresmann, \emph{Categories et structures: extraits}. Seminaire Ehresmann. Topologie et geometrie differentielle \textbf{6} (1964) pp.~1-31.

\bibitem%[ekkk74]
{ekkk74} H. Ehrig, K.-D. Kiermeier, H.-J. Kreowski, and W. K\"{u}hnel,\emph{ Universal theory of automata. A categorical approach.} Teubner Studienb\"{u}cher: Informatik. B.~G. Teubner, Stuttgart, 1974.

\bibitem%[EM45]
{EM45} S. Eilenberg and S. MacLane, \emph{General theory of natural equivalences}. Trans. Amer. Math. Soc. \textbf{58} (1945), pp.~231--294. 

\bibitem%[EW67]
{EW67} S. Eilenberg and J.~B. Wright, \emph{Automata in general algebras}. Information and Control \textbf{11} (1967),  pp.~452--470.

%\bibitem%[Eng75]
%{Eng75} J. Engelfriet, \emph{Bottom-up and top-down tree transformations. A comparison.} MST \textbf{9} (1957), pp.~198--231.

%\bibitem%[EngSch77]
%{EngSch77} J. Engelfriet and E.~M. Schmidt, \emph{IO and OI}. I. J. Comput. System Sci. \textbf{15}  (1977), pp.~328--353.

%\bibitem%[je15]
%{je15} J. Engelfriet, \emph{Tree automata and tree grammars}. Arxiv preprint arXiv:1510.02036v1, 2015.

%\bibitem%[Fuj59]
%{Fuj59}
%T. Fujiwara, \emph{On mappings between algebraic systems}, Osaka Math. J., \textbf{11}, pp.~153--172, 1959.

%\bibitem%[Fuj59]
%{Fuj60}
%T. Fujiwara, \emph{On mappings between algebraic systems II}, Osaka Math. J., \textbf{12}, pp.~253--268, 1960.

\bibitem%[end77]
{end77} H.~B. Enderton, \emph{Elements of set theory}. Academic Press [Harcourt Brace Jovanovich, Publishers], New York-London, 1977. 

\bibitem%[Formalpedia26]
{Formalpedia26} Prove2me, \emph{Formalpedia}. \url{https://prove2.me/formalpedia}. Accessed September 2026.

\bibitem%[Gar26]
{Gar26} D.~Z. Garza, Post on X, 10 September 2026. \url{https://x.com/dzackgarza/status/2098092343702548680}.

\bibitem%[GJ17]
{GJ17} N. Gambino and A. Joyal, \emph{On operads, bimodules and analytic functors}. Mem. Amer. Math. Soc.  \textbf{249} (2017), no. 1184.

\bibitem%[GK15]
{GK15} C. Ganesamoorthy and S.~G. Karpagavalli, \emph{A congruence relation in partially ordered sets}. Thai J. Math. \textbf{13} (2015), pp.~399-403.

\bibitem%[GRMSVCL26]
{GRMSVCL26} L. Gong, R. Ruiz Mora, N. Sanmart\'{\i}n Vich and E. Cosme Ll\'{o}pez, \emph{A Kleene theorem for free many-sorted algebras}. arXiv preprint arXiv:2606.29939 (2026).

\bibitem%[GS84]
{GS84} F. G\'{e}cseg and M. Steinby, \emph{Tree automata}. Akad\'{e}miai Kiad\'{o}, Budapest, 1984.

% [No citada des del 04/10/2026, en escurçar la secció 7.5]
%\bibitem%[Gut26]
%{Gut26} L. Guth, \emph{What math means to me}. arXiv preprint arXiv:2609.32028 (2026).

%\bibitem%[GS97]
%{GS97} F. G\'{e}cseg and M. Steinby, \emph{Tree languages}. In G. Rozenberg and A. Salomaa, editors.
%\emph{Handbook of formal languages. Vol. 3. Beyond words}, Chapter 1, pp.~1--68, Springer-Verlag, Berlin, 1997.

\bibitem%[G31]
{G31} K. G\"{o}del, \emph{\"{U}ber formal unentscheidbare S\"{a}tze der Principia Mathematica und verwandter Systeme. I.} 
Monatsh. f. Mathematik und Physik \textbf{38} (1931), pp.~173--198.

\bibitem%[GTW]
{gtw85} J. Goguen, J. Thatcher, and E. Wagner, \emph{An initial algebra approach to the specification, correctness, and implementation of abastract data types}. In R. Yeh, editor. \emph{Current Trends in Programming Methodology, IV}, Chapter 5, pp.~80--149, Prentice Hall, New Jersey, 1978.

\bibitem%[GM85]
{gm85} J. Goguen and J. Meseguer, \emph{Completeness of many-sorted equational logic}. Houston Journal of Mathematics, \textbf{11} (1985), pp.~307--334.

%\bibitem%[GB86]
%{gb86} J. Goguen and R. Burstall, \emph{A study in the foundations of programming methodology: Specifications, institutions, charters and parchments}. In D. Pitt et alii. eds. \emph{Proc. Summer Workshop on Category Theory and Computer Programming},
%Springer-Verlag (1986), pp.~313--333.

%\bibitem%[GW23]
%{GW23} U. G\"{o}rtz and T. Wedhorn, \emph{Algebraic geometry II: Cohomology of Schemes}. Springer Stud. Math--Master. Springer Spektrum, Wiesbaden, 2023.

\bibitem%[Grt08]
{gG08} G. Gr\"{a}tzer, \emph{Universal algebra. With appendices by Gr\"{a}tzer, Bjarni J\'{o}nsson, Walter Taylor, Robert W. Quackenbush, G\"{u}nter H. Wenzel, and Gr\"{a}tzer and W.~A. Lampe}. Revised reprint of the 1979 second edition. Springer, New York, 2008.

%\bibitem%[Gry66]
%{jwG66} J.~W. Gray, \emph{Fibred and cofibred categories.} In S. Eilenberg et alii. eds. \emph{Proc. Conf. Categorical Algebra}, Springer-Verlag (1966), pp.~21--83.

%\bibitem%[gree52]
%{gree52} J.~A. Green, \emph{A duality in abstract algebra}. J. London Math. Soc., \textbf{27} (1952), pp.~64--73.

\bibitem%[Gro71]
{Gro71} A. Grothendieck. \emph{Cat\'{e}gories fibr\'{e}es et
descente (Expos\'{e} VI)}.  In A. Grothendieck, editor,
Rev\^{e}tements \'{e}tales et groupe fondamental (SGA 1),
Springer-Verlag, Berlin $\cdot$ Heidelberg $\cdot$ New York, 1971, pp.~145--194.

\bibitem{Gui04}
Y. Guiraud, \emph{Pr\'esentations d'op\'erades et syst\`emes de r\'e\'ecriture}. PhD thesis, Universit\'e Montpellier 2, 2004.

\bibitem{GM09}
Y. Guiraud and P. Malbos, \emph{Higher-dimensional categories with finite derivation type}. Theory and Applications of Categories Vol. \textbf{22} (2009), pp.~420--478.

\bibitem{Gui19}
Y. Guiraud, \emph{Rewriting methods in higher algebra}, Mémoire d'Habilitation à Diriger des Recherches (2019), Univ. Paris 7.

\bibitem%[Har26]
{Har26} K. Hartnett, \emph{The proof in the code: how a truth machine is transforming math and AI}. Farrar, Straus and Giroux, New York, 2026.

\bibitem%[H05]
{H05} E. Harzheim, \emph{Ordered sets}. Advances in Mathematics (Springer), \textbf{7}. Springer, New York, 2005. 

\bibitem%[Her30]
{Her30} J. Herbrand, \emph{Recherches sur la th\'{e}orie de la d\'{e}monstration}, PhD Thesis, Universit\'{e} de Paris, 1930.

\bibitem%[hs73]
{hs73} H. Herrlich and G.~E. Strecker, \emph{Category theory: an introduction}. Allyn and Bacon Series in Advanced Mathematics. Allyn and Bacon Inc., Boston, Mass., 1973.


\bibitem%[Hig63]
{Hig63} P.~J. Higgins, \emph{Algebras with a scheme of operators}. Math. Nachr. \textbf{27} (1963), pp.~115--132. 

\bibitem%[h71]
{h71} P.~H. Higgins, \emph{Notes on categories and groupoids}. Van Nostrand Reinhold Co., London-New York-Melbourne, 1971.

\bibitem%[hs07]
{hs07} H.-J. Hoehnke and J. Schreckenberger, \emph{Partial algebras and their theories}. Shaker Verlag, Aachen, 2007. 

%\bibitem%[HoTT13]
%{HoTT13} \emph{Homotopy Type Theory: Univalent Foundations of Mathematics. The Univalent Foundations Program}, Institute for Advanced Study, 2013.

\bibitem%[Hook85]
{Hook85} J.~L. Hook, \emph{A note on interpretations of many-sorted theories}. The Journal of Symbolic Logic, Vol. \textbf{50} (1985), pp.~372--374.

\bibitem%[Hotz65]
{Hotz65} G. Hotz, \emph{Eine Algebraisierung des Syntheseproblems von Schaltkreisen. $\mathrm{I}$.} 
Elektronisch Informationsverarbeitung und Kybernetik, Vol. \textbf{1} (1965), pp.~185--205.

\bibitem%[Hotz65a]
{Hotz65a} G. Hotz, \emph{Eine Algebraisierung des Syntheseproblems von Schaltkreisen. $\mathrm{II}$.} 
Elektronisch Informationsverarbeitung und Kybernetik, Vol. \textbf{1} (1965), pp.~209--231.

\bibitem%[Hotz66]
{Hotz66} G. Hotz, \emph{Eindeutigkeit and Mehrdeutigkeit formaler Sprachen.}
Elektronisch Informationsverarbeitung und Kybernetik, Vol. \textbf{2} (1966), pp.~235--246

\bibitem%[H80]
{H80} W.~A. Howard, \emph{The formulae-as-types notion of construction.} In J.~P. Seldin and J.~R. Hindley, editors. \emph{To H.~B. Curry: Essays on Combinatory Logic, Lambda Calculus and Formalism}, pp.~479--490, Academic Press, London-New York, 1980.

\bibitem%[bj99]
{bj99} B. Jacobs, \emph{Categorical logic and type theory}. Stud. Logic Found. Math., \textbf{141}. North-Holland Publishing Co., Amsterdam, 1999.

\bibitem%[MJ88]
{MJ88} M. Jantzen,  
\emph{Confluent string rewriting}.
EATCS Monographs on Theoretical Computer Science, \textbf{14}. Springer-Verlag, Berlin, 1988.

\bibitem%[J87]
{J87} M. Johnson, \emph{Pasting diagrams in n-categories with applications to coherence theorems and categories of paths}, PhD thesis, University of Sydney, 1987.

\bibitem%[JY21]
{JY21} N. Johnson and D. Yau, \emph{2-dimensional categories}. Oxford University Press, Oxford, 2021.

%\bibitem%[jt51]
%{jt51} B. J\'{o}nsson and A. Tarski, \emph{Boolean algebras with operators, I}. Amer. J. Math., \textbf{73} (1951), pp.~891--939.
%
%\bibitem%[jt52]
%{jt52} B. J\'{o}nsson and A. Tarski, \emph{Boolean algebras with operators, II}. Amer. J. Math., \textbf{74} (1952), pp.~127--162.

\bibitem%[Jon72]
{bJ72}
B. J\'{o}nsson, \emph{Topics in universal algebra}, Springer-Verlag, 1972.

%\bibitem%[kl80]
%{kl80} M. Kline, \emph{Mathematics. The loss of certainty}. Oxford University Press, New York, 1980.

% [No citada des del 04/10/2026, en escurçar la secció 7.5]
%\bibitem%[Kal26]
%{Kal26} G. Kalai, \emph{``Be a Grothendieck!'': on AI and mathematics}. Combinatorics and more (blog), 22 September 2026. \url{https://gilkalai.wordpress.com/2026/09/22/be-a-grothendieck-on-ai-and-mathematics/}.

\bibitem%[kr72]
{kr72} H. Kaphengst and H. Reichel, \emph{Operative Theorien und Kategorien von operativen Systemen}. In  \emph{Studien zur Algebra und ihre Anwendungen (mit Anwendungen in der Mathematik, Physik und Rechentechnik)}, pp.~41--56. Schr. Zentralinst. Math. Mech. Akad. Wissensch. DDR, No. 16. Akademie Verlag, Berlin, 1972.

\bibitem%[kkm00]
{kkm00} M. Kilp, U. Knauer and A.~V. Mikhalev, \emph{Monoids, acts and categories. With applications to wreath products and graphs. A handbook for students and researchers}. Walter de Gruyter \!\!$\And$\!\! Co., Berlin, 2000.

\bibitem{Klo92}
J.~W. Klop, \emph{Term rewriting systems}. Handbook of Logic in Computer Science, Vol. \textbf{2}, Oxford University Press, 1992, pp.~1-117.

\bibitem%[agk67]
{agk67} A.~G. Kurosh, \emph{Alg\`{e}bre g\'{e}n\'{e}rale}. (French) Traduit du Russe par J.-P. Peaudecerf avec la participation de G. Lallement Collection Universitaire de Math\'{e}matiques, No. 22. Dunod, Paris, 1967.

\bibitem%[Kou26]
{Kou26} D. Koukoulopoulos, \emph{A CERN for AI-assisted science?} Guest post on T.~Tao's blog, 17 September 2026. \url{https://terrytao.wordpress.com/2026/09/17/a-cern-for-ai-assisted-science/}.

\bibitem{Laf03}
Y. Lafont, \emph{Towards an algebraic theory of Boolean circuits}. Journal of Pure and Applied Algebra \textbf{184} (2003), pp.~257-310.

\bibitem{LMW10}
Y. Lafont, F. M\'{e}tayer and K. Worytkiewicz, \emph{A folk model structure on omega-cat}. Advances in Mathematics \textbf{224}.3 (2010), pp.~1183--1231.

\bibitem%[Lk81]
{Lk81} J. Lambek, \emph{The Influence of Heraclitus on Modern Mathematics}. In: Agassi, J., Cohen, R.~S. (eds) Scientific Philosophy Today. Boston Studies in the Philosophy of Science, vol \textbf{67}. Springer, Dordrecht, 1981, pp.~111--121.

\bibitem%[LS86]
{LS86} J. Lambek and P.~J. Scott, \emph{Introduction to higher order categorical logic}. Cambridge Stud. Adv. Math., \textbf{7}. Cambridge University Press, Cambridge, 1986.

\bibitem{Lan21}
M. Land, \emph{Introduction to Infinity-categories}. Springer Nature, 2021.

%\bibitem%[l98]
%{l98} M. V. Lawson, \emph{Inverse semigroups. The theory of partial symmetries}. World Scientific Publishing Co., Inc., River Edge, NJ, 1998. 

\bibitem{Law03}
F.~W. Lawvere, \emph{Functorial semantics of algebraic theories}. Proceedings of the National Academy of Sciences \textbf{50} (1963), pp.~869-872.

%\bibitem%[L94]
%{L94} F. W. Lawvere, \emph{Cohesive toposes and Cantor's ``lauter Einsen''}. Categories in the foundations of mathematics and language. Philos. Math. (3) \textbf{2} (1994), pp.~5--15. 

\bibitem{Lei01}
T. Leinster, \emph{A survey of definitions of n-category}. arXiv preprint math/0107188 (2001).

\bibitem%[L02]
{L02} T. Leinster, \emph{A survey of definitions of $n$-category}. Theory Appl. Categ. \textbf{10} (2002), pp.~1--70.

\bibitem{Lei04}
T. Leinster, \emph{Higher operads, higher categories}. No. \textbf{298}. Cambridge University Press, 2004.

\bibitem%[LG17]
{LG17} S. Lucas and R. Guti\'{e}rrez, \emph{Automatic Synthesis of Logical Models for Order-Sorted First-Order Theories}. Journal of Automated Reasoning, \textbf{59} (2017), pp.~381--412. DOI: 10.1007/s10817-017-9419-3.

%\bibitem%[Lall79]
%{Lall79} G. Lallement, \emph{Semigroups and combinatorial applications}. Pure and Applied Mathematics. A Wiley-Interscience Publication. John Wiley $\!\!\And\!\!$ Sons, New York-Chichester-Brisbane, 1979.

%\bibitem%[sl94]
%{sl94} S. Lang, \emph{Algebraic number theory}. Second edition. Graduate Texts in Mathematics, 110. Springer-Verlag, New York, 1994.

%\bibitem%[sl96]
%{sl96} S. Lang, \emph{Topics in cohomology of groups.} Translated from the 1967 French original by the author. Chapter X based on letters written by John Tate. Springer-Verlag, Berlin, 1996.

%\bibitem%[law96]
%{law96} F. W. Lawvere,
%\emph{Adjoints in and among bicategories}. In \emph{Logic and algebra}.
%Papers from the International Conference in memory of Roberto Magari held in Pontignano, April 26-30, 1994. Edited by Aldo Ursini and Paolo Aglian\`{o}. Lecture Notes in Pure and Applied Mathematics, 180, pp.~181--189. Marcel Dekker, Inc., New York, 1996.
%
%\bibitem%[lc63]
%{lc63} L. Lesieur and R. Croisot, \emph{Alg\`{e}bre noeth\'{e}rienne non commutative}. 
%Gauthier-Villars \!$\And$\! Cie, \'{E}diteur-Imprimeur-Libraire, Paris, 1963. 

\bibitem%[Lax26]
{Lax26} \'{E}. Bonnet, J. Dreier and C. Kuske, \emph{Lax: the Lean archive}. \url{https://laxarchive.org}. Accessed September 2026.

\bibitem%[LaxAbout26]
{LaxAbout26} \'{E}. Bonnet, J. Dreier and C. Kuske, \emph{About Lax}. \url{https://laxarchive.org/about.html}. Accessed September 2026.

\bibitem%[LeanFRO26]
{LeanFRO26} Lean FRO, \emph{About the Lean FRO}. Accessed 4 October 2026. \url{https://lean-lang.org/fro/about/}.

\bibitem%[Lit26]
{Lit26} D. Litt, \emph{A beginning for mathematics}. Proofs and Prompts (blog), 14 September 2026. \url{https://proofsandprompts.com/2026/09/14/a-beginning-for-mathematics/}.

\bibitem%[KR35]
{KR35} S.~C. Kleene and J.~B. Rosser, \emph{The inconsistency of certain formal logics}. Ann. of Math. (2) \textbf{36}(3) (1935), pp.~630--636.

\bibitem%[lj97]
{lj97} E.~S. Ljapin and A.~E. Evseev,
\emph{The theory of partial algebraic operations}. Translated from the 1991 Russian original and with a preface by J. M. Cole. Revised by the authors. Mathematics and its Applications, 414. Kluwer Academic Publishers Group, Dordrecht, 1997. 

\bibitem%[mathlib20]
{mathlib20} The mathlib Community, \emph{The Lean mathematical library}. In Proceedings of the 9th ACM SIGPLAN International Conference on Certified Programs and Proofs (CPP 2020), pp.~367--381. ACM, 2020.

\bibitem%[mathlibDocs26]
{mathlibDocs26} The mathlib Community, \emph{Mathlib documentation}. \url{https://leanprover-community.github.io/mathlib4_docs/}. Accessed September 2026.

% [No citada des del 04/10/2026, en escurçar la secció 7.5]
%\bibitem%[MAIT26]
%{MAIT26} \emph{Mathematics and AI: an open letter on human--AI collaboration in mathematics}. Open letter, accessed 4 October 2026. \url{https://sites.google.com/view/math-and-ai-together/home}.

\bibitem%[Mac98]
{sM98} S. Mac Lane, \emph{Categories for the working mathematician.} 2nd ed. Springer-Verlag, New York, 1998.

%\bibitem{Mac13}
%S. Mac Lane, \emph{Categories for the working mathematician}. Vol. 5. Springer Science \& Business Media, 2013.

%\bibitem%[MM69]
%{MM69} M. Magidor and G. Moran, \emph{Finite automata over finite trees}. Technical report, \textbf{30}, Hebrew University, Jerusalem, 1969.

%\bibitem%[Maib74]
%{Maib74} T. S. E. Maibaum, \emph{A generalized approach to formal languages}. J. Comput. System Sci.  \textbf{8} (1974), pp.~409--439.
%
%\bibitem%[Maib77]
%{Maib77} T. S. E. Maibaum, \emph{Erratum}. J. Comput. System Sci.  \textbf{14} (1977), p.~369.

\bibitem%[em76]
{em76} E.~G. Manes, \emph{Algebraic theories}, Springer-Verlag, New York, 1976.

\bibitem%[Mar47]
{Mar47} A. Markov, \emph{On the impossibility of certain algorithms in the theory of associative systems}. C. R. (Doklady) Acad. Sci. URSS (N.S.) \textbf{55} (1947), pp.~583--586.

%\bibitem%[LS84]
%{LS84} P. Martin-L{\"o}f, \emph{Intuitionistic type theory}, \textbf{9} Bibliopolis, Naples, 1984.

\bibitem%[m76]
{m76} G. Matthiessen, \emph{Theorie der {H}eterogenen {A}lgebren}, Mathematik-{A}rbeitspapiere, Nr. 3, Universit\"{a}t Bremen Teil A, Mathematische Forchungspapiere, 1976.

%\bibitem%[MY60]
%{MY60} R. McNaughton and H. Yamada, \emph{Regular expressions and state graphs for automata}. IRE Transactions on Electronic Computers, \textbf{EC-9} (1960), pp.~39--47.

\bibitem{Met03}
F. M\'{e}tayer, \emph{Resolutions by polygraphs}. Theory and Applications of Categories \textbf{11} (2003), pp.~148--184.

%\bibitem%{MW67}
%{MW67} J. Mezei and J. Wright, \emph{Algebraic automata and context-free sets}. Information and Control, \textbf{11} (1967),   pp.~3--29.

% [No citada des del 04/10/2026, en escurçar la secció 7.5]
%\bibitem%[MDZ26]
%{MDZ26} J. Miao, J.~R. Davis, Y. Zhang, J.~K. Pritchard and J. Zou, \emph{Reimagining research papers as interactive and reliable AI agents}. Nature (2026). \url{https://www.nature.com/articles/s41586-026-11044-y}.

\bibitem%[Mi14]
{Mi14} C.~F. Miller\! III, \emph{Turing machines to word problems.} In R. Downey, editor. \emph{Turing's legacy: developments from Turing's ideas in logic,} pp.~329--385, Lect. Notes Log., \textbf{42} Association for Symbolic Logic, La Jolla, CA; Cambridge University Press, Cambridge, 2014.

\bibitem{Mim14}
S. Mimram, \emph{Towards 3-dimensional rewriting theory}. Logical Methods in Computer Science \textbf{10} (2014), pp.~1--47.

%\bibitem%[Myh57]
%{Myh57} J. Myhill, \emph{Finite automata and the representation of events.}  WADD TR-57-624, Wright Patterson AFB, Ohio (1957), pp.~112--137.
%

\bibitem%[Mul24]
{Mul24} S. M\"{u}ller-Stach, \emph{Max Dehn, Axel Thue, and the Undecidable} In J. Lorenat, J. McClearly, V.~R. Remmert, D.~E. Rowe \& M. Senechal (Eds) \emph{Max Dehn: Polyphonic portrait}, Chapter 8, pp.~181--196, American Mathematical Society, \textbf{46}, 2024.

%\bibitem%[Ner58]
%{Ner58} A. Nerode, \emph{Linear automata transformation.} In Proc. American Mathematical Society, \textbf{9} (1958), pp.~541--544.

\bibitem%[N55]
{N55}
P.~S. Novikov, \emph{On the algorithmic unsolvability of the word problem in group theory}. Akademiya Nauk S.S.S.R. Matematicheskii Institut Trudy, \textbf{44}, Moscow (1955), pp.~1--143.

%\bibitem%[po11]
%{po11} S.~M. Pinto and M.~T. Oliveira-Martins, \emph{Congruences and ideals on Boolean modules: a heterogeneous point of view}. Math. Log. Quart. \textbf{57} (2011), pp.~571--581.
%
%\bibitem%[po12]
%{po12} S.~M. Pinto and M.~T. Oliveira-Martins, \emph{Congruences and ideals on Peirce algebras: a heterogeneous/homogeneous point of view}. Math. Log. Quart. \textbf{58} (2012), pp.~252--262.

\bibitem%[bp72]
{bp72} B.~I. Plotkin, \emph{Groups of automorphisms of algebraic systems}. Translated from the Russian by K.~A. Hirsch. Wolters-Noordhoff Publishing, Groningen, 1972.

\bibitem%[OAINS26]
{OAINS26} OpenAI, \emph{On the Navier--Stokes Millennium Prize Problem}. OpenAI, 8 September 2026. \url{https://openai.com/index/navier-stokes-solution/}.

\bibitem%[OAINSRepo]
{OAINSRepo} OpenAI, \emph{NavierStokesAndEuler: Lean certificates accompanying Navier--Stokes and Euler results}. Code repository. \url{https://github.com/openai/NavierStokesAndEuler}.

\bibitem%[OAIAG26]
{OAIAG26} OpenAI, \emph{Advisory Group on Mathematics and Artificial Intelligence}. OpenAI, 21 September 2026. \url{https://openai.com/index/advisory-group-on-mathematics-and-ai/}.

\bibitem%[PT91]
{PT91}H.-E. Porst and W. Tholen, \emph{Concrete dualities}. Category theory at work (Bremen, 1990), pp.~111--136. Res. Exp. Math., \textbf{18} Heldermann Verlag, Berlin, 1991.


\bibitem%[Poz26]
{Poz26} B. Pozsgay, Post on X, 23 September 2026. \url{https://x.com/pozsgaybalazs/status/2102896555418177667}.

\bibitem%[Post43]
{Post43} E.~L. Post, \emph{Formal reductions of the general combinatorial decision problem}. Amer. J. Math.  \textbf{65} (1943), pp.~197--215.

\bibitem%[Post44]
{Post44} E.~L. Post. Recursively enumerable sets of positive integers and their decision problems. Bulletin of the American Mathematical Society \textbf{50} (1944), pp.~284--316.

\bibitem%[Post47]
{Post47} E.~L. Post, \emph{Recursive unsolvability of a problem of Thue}. The Journal of Symbolic Logic \textbf{12}(1) (1947), pp.~1--11.

\bibitem%[Post65]
{Post65}  E.~L. Post, \emph{Absolutely unsolvable problems and relatively undecidable
propositions---account of an Anticipation}. In \emph{The undecidable. Basic papers on undecidable propositions, unsolvable problems and computable functions}, edited by Martin Davis, Raven Press Books, Ltd., Hewlett, N.~Y. 1965, pp.~340--433.

\bibitem%[Post04]
{Post04}  E.~L. Post, \emph{Absolutely unsolvable problems and relatively undecidable propositions---account of an Anticipation}. In \emph{The undecidable. Basic papers on undecidable propositions, unsolvable problems and computable functions}, edited by Martin Davis, Dover Publications, Inc., Mineola, New York, 2004, pp.~340--406.

\bibitem{Pow91}
A.~J. Power, \emph{An $n$-categorical pasting theorem}. In \emph{Category theory}, volume 1488 of Lecture Notes in Mathematics, pp.~326-358, Springer, 1991.

\bibitem{Pow89}
A.~J. Power, \emph{An abstract formulation for rewrite systems}. In \emph{Category Theory and Computer Science}, Manchester, UK, September 5--8, 1989 Proceedings, pp.~300--312, Springer Berlin Heidelberg, 2005.

\bibitem{Pow14}
J.~F. Power, \emph{An Early Completion Algorithm: Thue's 1914 Paper on the Transformation of Symbol Sequences}. In: A. Beckmann, E. Csuhaj-Varj\'{u} and K. Meer  (Eds.) Language, Life, Limits. CiE 2014. Lecture Notes in Computer Science, 
vol.~\textbf{8493},  pp.~343--346, Springer, Cham, 2014.

%\bibitem%[Rab69]
%{Rab69} M.~O. Rabin, \emph{Decidability of second-order theories and automata on infinite trees}. Trans. Amer. Math. Soc., \textbf{141} (1969), pp.~1--35.

%\bibitem%[RS59]
%{RS59} M.~O. Rabin and D. Scott, \emph{Finite automata and their decision problems}. IBM J. Res. Develop. \textbf{3} (1959) pp.~114--125.
%
%\bibitem%[rs04]
%{rs04} M. Raussen and Chr. Skau, \emph{Interview with Michael Atiyah and Isadore Singer}. Eur. Math. Soc. Newsl. \textbf{53}  (2004), pp.~24--30.

\bibitem%[P2M26]
{P2M26} Prove2me, \emph{Prove2me: humans and AI agents proving theorems in Lean~4}. \url{https://prove2.me}. Accessed September 2026.

\bibitem%[P2MStart26]
{P2MStart26} Prove2me, \emph{Getting started}. \url{https://prove2.me/start.md}. Accessed September 2026.

\bibitem%[P2MEil26]
{P2MEil26} Prove2me, \emph{Mission: Eilenberg theorems for many-sorted formations}. Accepted 9 September 2026. \url{https://prove2.me/missions/Eilenberg_Theorems_for_Many-Sorted_Formations}.

\bibitem%[P2MKle26]
{P2MKle26} Prove2me, \emph{Mission: A Kleene theorem for free many-sorted algebras}. Accepted 9 September 2026. \url{https://prove2.me/missions/A_Kleene_Theorem_for_Free_Many-Sorted_Algebras}.

\bibitem%[QW26]
{QW26} Y. Qi and M. Weber, \emph{Choir: an open protocol for distributed multi-agent autoformalization}. arXiv preprint arXiv:2609.31903 (2026).

\bibitem%[Rei87]
{Rei87} H.~Reichel, \emph{Initial Computability, Algebraic Specifications, and Partial Algebras}. International Series of Monographs on Computer Science, Clarendon Press, Oxford, 1987.

\bibitem%[r87]
{r87} H. Reichel, \emph{Initial computability, algebraic specifications, and partial algebras}. The Clarendon Press, Oxford University Press, New York, 1987.

\bibitem{RV16}
E. Riehl and D. Verity, \emph{Infinity category theory from scratch}. arXiv preprint arXiv:1608.05314 (2016).

\bibitem%[rv22]
{rv22} E. Riehl and D. Verity, \emph{Elements of $\infty$-category theory}. Cambridge University Press, Cambridge, 2022.

\bibitem%[R08]
{R08} S. Roman, \emph{Lattices and ordered sets}. Springer, New York, 2008.

\bibitem%[pr18]
{pr18} P. Roquette, \emph{The Riemann Hypothesis in Characteristic $p$ in Historical Perspective}, Springer Nature Switzerland AG, 2018.

%\bibitem%[kr96]
%{kr96} K.~I. Rosenthal, \emph{The theory of quantaloids}. Longman, Harlow, 1996.

%\bibitem%[cr10]
%{cr10} C. Russo,\emph{ Quantale modules and their operators, with applications}. J. Logic Comput. \textbf{20} (2010), pp.~917--946. 

\bibitem%[Sah26]
{Sah26} A. Sahai, \emph{We're gonna need a lot more mathematicians}. Guest post on T.~Tao's blog, 24 September 2026. \url{https://terrytao.wordpress.com/2026/09/24/were-gonna-need-a-lot-more-mathematicians/}.

\bibitem%[SAIR26]
{SAIR26} T. Tao, \emph{SAIR's open math model initiative}. What's new (blog), 18 September 2026. \url{https://terrytao.wordpress.com/2026/09/18/sairs-open-math-model-initiative/}.

\bibitem%[sch66]
{sch66} J. Schmidt, \emph{A general existence theorem on partial algebras and its special cases}. Colloq. Math. \textbf{14} (1966) pp.~73--87.

\bibitem%[sch70]
{sch70} J. Schmidt, \emph{A homomorphism theorem for partial algebras}. Colloq. Math. \textbf{21} (1970) pp.~5--21.

\bibitem%[Sch24]
{Sch24} M. Sch\"{o}nfinkel, \emph{\"{U}ber die Bausteine der mathematischen Logik}. Math. Ann. \textbf{92}(3) (1924) pp.~305--316.

%\bibitem%[Sch55]
%{Sch55} M.~P. Sch{\"{u}}tzenberger, \emph{Une th\'{e}orie alg\'{e}brique du codage}. S\'{e}minaire Dubreil. Alg\`{e}bre et th\'{e}orie des nombres, \textbf{9} (1955--1956), Expos\'{e} No. 15, pp.~1--24.

\bibitem%[Sch22]
{Sch22} P. Scholze, \emph{Liquid Tensor Experiment}. Experimental Mathematics \textbf{31} (2022), pp.~349--354.

\bibitem%[Sel06]
{Sel06} J.~P. Seldin, \emph{The logic of Curry and Church}. In: D.~M. Gabbay and J. Woods (Eds.) Handbook of the History of Logic. Vol.~\textbf{5}. Logic from Russell to Church (2009), pp.~819--873. Elsevier/North-Holland, Amsterdam, 2009.

\bibitem%[sl68]
{sl68} J. S{\l}omi\'{n}ski, \emph{Peano-algebras and quasi-algebras}.
Dissertationes Math. (Rozprawy Mat.) \textbf{57} (1968).

%\bibitem%[sl74]
%{sl74} J. S{\l}omi\'{n}ski, \emph{On the greatest congruence relation contained in an equivalence relation and its applications to the algebraic theory of machines}. Colloq. Math.  \textbf{29} (1974), pp.~31--43.

\bibitem{SOK94}
C.~C. Squier, F. Otto and Y. Kobayashi, \emph{A finiteness condition for rewriting systems}. Theoretical Computer Science \textbf{131} (1994), pp.~271--294.

\bibitem{Str76}
R. Street, \emph{Limits indexed by category-valued 2-functors}. Journal of Pure and Applied Algebra \textbf{8}.2 (1976), pp.~149-181.

\bibitem{Str87}
R. Street, \emph{The algebra of oriented simplexes}. Journal of Pure and Applied Algebra \textbf{49}.3 (1987), pp.~283--335.

%\bibitem%[Ste81]
%{Ste81} M. Steinby, \emph{Some algebraic aspects of recognizability and rationality}. In F. G\'{e}cseg, editor, \emph{Fundamentals of Computation Theory: Proceedings of the 1981 International FCT-Conference, Szeged, Hungary, August 24--28 (1981)} pp.~360--372, Berlin-Heidelberg.

\bibitem%[ST00]
{ST00} M. Steinby and W. Thomas, \emph{Trees and term rewriting in 1910: on a paper by Axel Thue}. Bull. Eur. Assoc. Theor. Comput. Sci. EATCS \textbf{72} (2000), pp.~256--269.

\bibitem%[St72]
{St72} S. Stenlund, \emph{Combinators, $\lambda$-terms and proof theory}. D. Reidel Publishing Co., Dordrecht, 1972.

\bibitem%[Sum26]
{Sum26} M. Abouzaid, M. Brenner, I. Corwin, S. Fomin, D. Freed, J. G\'{o}mez-Serrano, L. Guth, J. Harris, M. Kisin, B. Kra, D. Litt, M. Matchett Wood, C. McMullen, A. Naor, M. Olsson, T. Perutz, E. Riehl, R. Rouquier, T. Scanlon, G. Staffilani, R. Vakil, R. Ward, L. Williams and C. Xu, \emph{Report of Summit on PhD Math Education in the Age of AI}. Summit held on 17--18 September 2026. \url{https://cmsa.fas.harvard.edu/media/2026/09/Summit-on-PhD-Math-Education-in-the-Age-of-AI.pdf}.

\bibitem%[T08]
{T08} H. Tietze, \emph{\"{U}ber die topologischen Invarianten mehrdimensionaler Mannigfaltigkeiten}, Monat. Math. Phys. \textbf{19} (1908), pp.~1--118.

\bibitem%[TaoFM26]
{TaoFM26} T. Tao et al., \emph{A severe misalignment of AI in mathematics}. Joint declaration of 25 Fields medallists, What's new (blog), 11 September 2026. \url{https://terrytao.wordpress.com/2026/09/11/a-severe-misalignment-of-ai-in-mathematics/}.

\bibitem%[TBG91]
{tbg91} A. Tarlecki, R. Burstall and J. Goguen, \emph{Some fundamental algebraic tools for the semantics of computation: Part 3. Indexed categories}, Theoretical Computer Science, \textbf{91} (1991), pp.~239--264.

%\bibitem%[Tei51]
%{Tei51} M. Teissier, \emph{Sur les \'{e}quivalences r\'{e}guli\`{e}res dans les demi-groupes}. C.R. Acad. Sci. Paris \textbf{232} (1951), pp.~1987--1989.

\bibitem%[Tha69]
{Tha69} J.~W. Thatcher, \emph{Transformations and translations from the point of view of generalized finite automata theory}. In \emph{Proceedings of the ACM Symposium on Theory Computing} (1969),
pp.~129--142.

%\bibitem%[Tha73]
%{Tha73} J.~W. Thatcher, \emph{Tree automata: an informal survey}. In A.~V. Aho, editor, \emph{Currents in the Theory of Computing}, Prentice-Hall, Englewoods Cliffs, N.~J., pp.~143--172, 1973.

%\bibitem%[TW65]
%{TW65} J.~W. Thatcher and J.~B. Wright, \emph{Generalized finite automata}. Notices Amer. Math. Soc. \textbf{12} (1965), Abstract No 65T-649, p.~820.
%
%\bibitem%[TW68]
%{TW68} J.~W. Thatcher and J.~B. Wright, \emph{Generalized finite automata theory with an application to a decision problem of second order logic}. MST \textbf{2} (1968), pp.~57--81.

\bibitem%[Th10]
{Th10} A. Thue, \emph{Die L\"{o}sung eines Spezialfalles eines generallen logischen Problems}. Videnskabs-Selskabets Skrifter Kristiania. I. Math.-Nat. Klasse 
No.~\textbf{8} (1910).
%, pp.~1--17. 

\bibitem%[Th12]
{Th12} A. Thue, \emph{\"{U}ber die gegenseitige Lage gleicher Teile gewisser Zeichenreihen}. Videnskabs-Selskabets Skrifter Kristiania. I. Math.-Nat. Klasse 
No.~\textbf{9} (1912).
%, pp.~1--67.

\bibitem%[Th14]
{Th14} A. Thue, \emph{Probleme \"{u}ber Ver\"{a}nderungen von Zeichenreihen nach gegebenen Regeln}. Videnskabs-Selskabets Skrifter Kristiania. I. Math.-Nat. Klasse No.~\textbf{10} (1914).
%, pp.~1--27.

\bibitem%[Th77]
{Th77} A. Thue, \emph{Selected Mathematical Papers}. With an introduction by Carl Ludwig Siegel and a biography by Viggo Brun. Edited by Trygve Nagell, Atle Selberg, Sigmund Selberg, and Knut Thalberg. Universitetsforlaget, Oslo-Bergen-Troms\o, 1977.

\bibitem%[Tur36]
{Tur36} A.~M. Turing, \emph{On Computable Numbers, with an Application to the Entscheidungsproblem}. Proc. London Math. Soc. \textbf{42} (1936), pp.~230--265.

\bibitem%[Vag26]
{Vag26} M. Vago Marzal, \emph{A Formalization of Higher-Order Categories in Lean~4} (Version 0.1.0). Zenodo, 2026. \url{https://doi.org/10.5281/zenodo.20631672}.

\bibitem%[VagRepo]
{VagRepo} M. Vago Marzal, \emph{Higher category theory in Lean~4}. Code repository and documentation. \url{https://mariovagomarzal.github.io/higher_category_theory/}.

%\bibitem%[Voe11]
%{Voe11} V. Voevodsky, \emph{Univalent foundations of mathematics}. In International Workshop on Logic, Language, Information, and Computation, pp.~4--4. Springer, 2011. 

%\bibitem%[w21]
%{w21} F. Wagemann, \emph{Crossed modules}. De Gruyter, Berlin, 2021.

\bibitem%[XTX26]
{XTX26} XTX Markets, \emph{2026 update on XTX Markets' AI philanthropy}. 16 September 2026. \url{https://www.xtxmarkets.com/news/2026-update-on-xtx-markets-ai-philanthropy/}.

\bibitem%[w92]
{w92} W. Wechler, \emph{Universal algebra for computer scientists}. Springer-Verlag, Berlin, 1992.

\bibitem%[Wol26]
{Wol26} S. Wolfram, \emph{What's the future for pure math research in the age of AI?} Stephen Wolfram Writings, 28 September 2026. \url{https://writings.stephenwolfram.com/2026/09/whats-the-future-for-pure-math-research-in-the-age-of-ai/}.

\bibitem%[uw90]
{uw90} U. Wolter, \emph{An algebraic approach to deduction in equational partial Horn theories}. Journal of Information Processing and Cybernetics. EIK, \textbf{27}, pp.~85--128, 1990.lag, Berlin, 1992.

\bibitem%[Yan26]
{Yan26} K. Yang, \emph{What happens when formalization becomes cheap?} Blog post, 22 September 2026. \url{https://yangky11.github.io/blog/ai-and-formalization/}.

%\bibitem%[y60]
%{y60} N. Yoneda, \emph{On Ext and exact sequences}. J. Fac. Sci. Univ. Tokyo Sect. I \textbf{8} (1960), pp.~ 507--576.


\end{thebibliography}



\end{document}
